\subsection{EE.ANDQ}\label{instruction:EEANDQ}
\setlength\tabcolsep{6pt}
\textbf{Instruction Word} \\
\begin{tabular}{|c|c|c|c|c|c|c|c|c|c|c|}
	\hline
	11 & qa[2:1] & 1101 & qa[0] & 011 & qy[2:1] & 00 & qx[2:1] & qy[0] & qx[0] & 0100 \\
	\hline
\end{tabular}

\textbf{Assembler Syntax} \\
EE.ANDQ qa, qx, qy \\
\textbf{Description} \\
This instruction performs a bitwise AND operation on registers \lstinline{qx} and \lstinline{qy} and writes the result of the logical operation to register \lstinline{qa}.\\
\textbf{Operation}
\begin{lstlisting}
  qa = qx & qy
\end{lstlisting}

\clearpage
\subsection{EE.BITREV}\label{instruction:EEBITREV}
\setlength\tabcolsep{6pt}
\textbf{Instruction Word} \\
\begin{tabular}{|c|c|c|c|c|c|c|}
	\hline
	11 & qa[2:1] & 1101 & qa[0] & 1111011 & as[3:0] & 0100 \\
	\hline
\end{tabular}

\textbf{Assembler Syntax} \\
EE.BITREV qa, as \\
\textbf{Description} \\
This instruction swaps the bit order of data of different bit widths according to the value of special register \lstinline{FFT_BIT_WIDTH}. Then, it compares the data before and after the swap, takes the larger value, pads the higher bits with 0 until it has 16 bits, and writes the result into the corresponding data segment of register \lstinline{qa}. In the following, Switchx is a function that represents the bit inversion of the x-bit. Switch5(0b10100) = 0b00101. Here 0b10100 means 5-bit binary data.\\
\textbf{Operation}
\begin{lstlisting}
  temp0[15:0] = as[15:0]
  temp1[15:0] = as[15:0] + 1
  temp2[15:0] = as[15:0] + 2
  temp3[15:0] = as[15:0] + 3
  temp4[15:0] = as[15:0] + 4
  temp5[15:0] = as[15:0] + 5
  temp6[15:0] = as[15:0] + 6
  temp7[15:0] = as[15:0] + 7
if FFT_BIT_WIDTH==3:
  Switch3(X[2:0]) = X[0:2]
  qa[ 15:  0] = {13'h0, max(tmp0[2:0], Switch3(tmp0[2:0]))}
  qa[ 31: 16] = {13'h0, max(tmp1[2:0], Switch3(tmp1[2:0]))}
  qa[ 47: 32] = {13'h0, max(tmp2[2:0], Switch3(tmp2[2:0]))}
  qa[ 63: 48] = {13'h0, max(tmp3[2:0], Switch3(tmp3[2:0]))}
  qa[ 79: 64] = {13'h0, max(tmp4[2:0], Switch3(tmp4[2:0]))}
  qa[ 95: 80] = {13'h0, max(tmp5[2:0], Switch3(tmp5[2:0]))}
  qa[127: 96] = 0
if FFT_BIT_WIDTH==4:
  Switch4(X[3:0]) = X[0:3]
  qa[ 15:  0] = {12'h0, max(tmp0[3:0], Switch4(tmp0[3:0]))}
  qa[ 31: 16] = {12'h0, max(tmp1[3:0], Switch4(tmp1[3:0]))}
  qa[ 47: 32] = {12'h0, max(tmp2[3:0], Switch4(tmp2[3:0]))}
  qa[ 63: 48] = {12'h0, max(tmp3[3:0], Switch4(tmp3[3:0]))}
  qa[ 79: 64] = {12'h0, max(tmp4[3:0], Switch4(tmp4[3:0]))}
  qa[ 95: 80] = {12'h0, max(tmp5[3:0], Switch4(tmp5[3:0]))}
  qa[111: 96] = {12'h0, max(tmp6[3:0], Switch4(tmp6[3:0]))}
  qa[127:112] = {12'h0, max(tmp7[3:0], Switch4(tmp7[3:0]))}
  ...
if FFT_BIT_WIDTH==10:
  Switch10(X[9:0]) = X[0:9]
  qa[ 15:  0] = {6'h0, max(tmp0[9:0], Switch10(tmp0[9:0]))}
  qa[ 31: 16] = {6'h0, max(tmp1[9:0], Switch10(tmp1[9:0]))}
  qa[ 47: 32] = {6'h0, max(tmp2[9:0], Switch10(tmp2[9:0]))}
  qa[ 63: 48] = {6'h0, max(tmp3[9:0], Switch10(tmp3[9:0]))}
  qa[ 79: 64] = {6'h0, max(tmp4[9:0], Switch10(tmp4[9:0]))}
  qa[ 95: 80] = {6'h0, max(tmp5[9:0], Switch10(tmp5[9:0]))}
  qa[111: 96] = {6'h0, max(tmp6[9:0], Switch10(tmp6[9:0]))}
  qa[127:112] = {6'h0, max(tmp7[9:0], Switch10(tmp7[9:0]))}

  as[31:0] = as[31:0] + 8
\end{lstlisting}

\clearpage
\subsection{EE.CLR\_BIT\_GPIO\_OUT}\label{instruction:EECLRBITGPIOOUT}
\setlength\tabcolsep{6pt}
\textbf{Instruction Word} \\
\begin{tabular}{|c|c|c|}
	\hline
	011101100100 & imm256[7:0] & 0100 \\
	\hline
\end{tabular}

\textbf{Assembler Syntax} \\
EE.CLR\_BIT\_GPIO\_OUT 0..255 \\
\textbf{Description} \\
It is a dedicated CPU GPIO instruction to clear certain GPIO\_OUT bits. The content to clear depends on the 8-bit immediate number \lstinline{imm256}.\\
\textbf{Operation}
\begin{lstlisting}
  GPIO_OUT[7:0] =  (GPIO_OUT[7:0] & ~imm256[7:0])
\end{lstlisting}

\clearpage
\subsection{EE.CMUL.S16}\label{instruction:EECMULS16}
\setlength\tabcolsep{6pt}
\textbf{Instruction Word} \\
\begin{tabular}{|c|c|c|c|c|c|c|c|c|c|c|}
	\hline
	10 & qz[2:1] & 1110 & qz[0] & qy[2] & 0 & qy[1:0] & qx[2:0] & 00 & sel4[1:0] & 0100 \\
	\hline
\end{tabular}

\textbf{Assembler Syntax} \\
EE.CMUL.S16 qz, qx, qy, 0..3 \\
\textbf{Description} \\
This instruction performs a 16-bit signed complex multiplication. The range of the immediate number \lstinline{sel4} is 0 \verb+~+ 3, which specifies the 32 bits in the two QR registers \lstinline{qx} and \lstinline{qy} for complex multiplication. The real and imaginary parts of complex numbers are stored in the upper 16 bits and lower 16 bits of the 32 bits respectively. The calculated real part and imaginary part results are stored in the corresponding 32 bits of register \lstinline{qz}.\\
\textbf{Operation}
\begin{lstlisting}
if sel4 == 0:
  qz[ 15:  0] = (qx[ 15:  0] * qy[ 15:  0] - qx[ 31: 16] * qy[ 31: 16]) >> SAR[5:0]
  qz[ 31: 16] = (qx[ 15:  0] * qy[ 31: 16] + qx[ 31: 16] * qy[ 15:  0]) >> SAR[5:0]
  qz[ 47: 32] = (qx[ 47: 32] * qy[ 47: 32] - qx[ 63: 48] * qy[ 63: 48]) >> SAR[5:0]
  qz[ 63: 48] = (qx[ 47: 32] * qy[ 63: 48] + qx[ 63: 48] * qy[ 47: 32]) >> SAR[5:0]
if sel4 == 1:
  qz[ 79: 64] = (qx[ 79: 64] * qy[ 79: 64] - qx[ 95: 80] * qy[ 95: 80]) >> SAR[5:0]
  qz[ 95: 80] = (qx[ 79: 64] * qy[ 95: 80] + qx[ 95: 80] * qy[ 79: 64]) >> SAR[5:0]
  qz[111: 96] = (qx[111: 96] * qy[111: 96] - qx[127:112] * qy[127:112]) >> SAR[5:0]
  qz[127:112] = (qx[111: 96] * qy[127:112] + qx[127:112] * qy[111: 96]) >> SAR[5:0]
if sel4 == 2:
  qz[ 15:  0] = (qx[ 15:  0] * qy[ 15:  0] + qx[ 31: 16] * qy[ 31: 16]) >> SAR[5:0]
  qz[ 31: 16] = (qx[ 15:  0] * qy[ 31: 16] - qx[ 31: 16] * qy[ 15:  0]) >> SAR[5:0]
  qz[ 47: 32] = (qx[ 47: 32] * qy[ 47: 32] + qx[ 63: 48] * qy[ 63: 48]) >> SAR[5:0]
  qz[ 63: 48] = (qx[ 47: 32] * qy[ 63: 48] - qx[ 63: 48] * qy[ 47: 32]) >> SAR[5:0]
if sel4 == 3:
  qz[ 79: 64] = (qx[ 79: 64] * qy[ 79: 64] + qx[ 95: 80] * qy[ 95: 80]) >> SAR[5:0]
  qz[ 95: 80] = (qx[ 79: 64] * qy[ 95: 80] - qx[ 95: 80] * qy[ 79: 64]) >> SAR[5:0]
  qz[111: 96] = (qx[111: 96] * qy[111: 96] + qx[127:112] * qy[127:112]) >> SAR[5:0]
  qz[127:112] = (qx[111: 96] * qy[127:112] - qx[127:112] * qy[111: 96]) >> SAR[5:0]
\end{lstlisting}

\clearpage
\subsection{EE.CMUL.S16.LD.INCP}\label{instruction:EECMULS16LDINCP}
\setlength\tabcolsep{6pt}
\textbf{Instruction Word} \\
\begin{tabular}{|c|c|c|c|c|c|c|c|c|c|c|c|c|}
	\hline
	111000 & qu[2:1] & qy[0] & 000 & qu[0] & qz[2:0] & qx[1:0] & qy[2:1] & 11 & sel4[1:0] & as[3:0] & 111 & qx[2] \\
	\hline
\end{tabular}

\textbf{Assembler Syntax} \\
EE.CMUL.S16.LD.INCP qu, as, qz, qx, qy, 0..3 \\
\textbf{Description} \\
This instruction performs a 16-bit signed complex multiplication. The range of the immediate number \lstinline{sel4} is 0 \verb+~+ 7, which specifies the 32 bits in the two QR registers \lstinline{qx} and \lstinline{qy} for complex multiplication. The real and imaginary parts of complex numbers are stored in the upper 16 bits and lower 16 bits of the 32 bits respectively. The calculated real part and imaginary part results are stored in the corresponding 32 bits of register \lstinline{qz}.\\
During the operation, the lower 4 bits of the access address in register \lstinline{as} are forced to be 0, and then the 16-byte data is loaded from the memory to register \lstinline{qu}. After the access, the value in register \lstinline{as} is incremented by 16.\\
\textbf{Operation}
\begin{lstlisting}
if sel4 == 0:
  qz[ 15:  0] = (qx[ 15:  0] * qy[ 15:  0] - qx[ 31: 16] * qy[ 31: 16]) >> SAR[5:0]
  qz[ 31: 16] = (qx[ 15:  0] * qy[ 31: 16] + qx[ 31: 16] * qy[ 15:  0]) >> SAR[5:0]
  qz[ 47: 32] = (qx[ 47: 32] * qy[ 47: 32] - qx[ 63: 48] * qy[ 63: 48]) >> SAR[5:0]
  qz[ 63: 48] = (qx[ 47: 32] * qy[ 63: 48] + qx[ 63: 48] * qy[ 47: 32]) >> SAR[5:0]
if sel4 == 1:
  qz[ 79: 64] = (qx[ 79: 64] * qy[ 79: 64] - qx[ 95: 80] * qy[ 95: 80]) >> SAR[5:0]
  qz[ 95: 80] = (qx[ 79: 64] * qy[ 95: 80] + qx[ 95: 80] * qy[ 79: 64]) >> SAR[5:0]
  qz[111: 96] = (qx[111: 96] * qy[111: 96] - qx[127:112] * qy[127:112]) >> SAR[5:0]
  qz[127:112] = (qx[111: 96] * qy[127:112] + qx[127:112] * qy[111: 96]) >> SAR[5:0]
if sel4 == 2:
  qz[ 15:  0] = (qx[ 15:  0] * qy[ 15:  0] + qx[ 31: 16] * qy[ 31: 16]) >> SAR[5:0]
  qz[ 31: 16] = (qx[ 15:  0] * qy[ 31: 16] - qx[ 31: 16] * qy[ 15:  0]) >> SAR[5:0]
  qz[ 47: 32] = (qx[ 47: 32] * qy[ 47: 32] + qx[ 63: 48] * qy[ 63: 48]) >> SAR[5:0]
  qz[ 63: 48] = (qx[ 47: 32] * qy[ 63: 48] - qx[ 63: 48] * qy[ 47: 32]) >> SAR[5:0]
if sel4 == 3:
  qz[ 79: 64] = (qx[ 79: 64] * qy[ 79: 64] + qx[ 95: 80] * qy[ 95: 80]) >> SAR[5:0]
  qz[ 95: 80] = (qx[ 79: 64] * qy[ 95: 80] - qx[ 95: 80] * qy[ 79: 64]) >> SAR[5:0]
  qz[111: 96] = (qx[111: 96] * qy[111: 96] + qx[127:112] * qy[127:112]) >> SAR[5:0]
  qz[127:112] = (qx[111: 96] * qy[127:112] - qx[127:112] * qy[111: 96]) >> SAR[5:0]

  qu[127:0] = load128({as[31:4],4{0}})
  as[31:0] = as[31:0] + 16
\end{lstlisting}

\clearpage
\subsection{EE.CMUL.S16.ST.INCP}\label{instruction:EECMULS16STINCP}
\setlength\tabcolsep{6pt}
\textbf{Instruction Word} \\
\begin{tabular}{|c|c|c|c|c|c|c|c|c|c|c|c|}
	\hline
	11100100 & qy[0] & qv[2:0] & 0 & qz[2:0] & qx[1:0] & qy[2:1] & 00 & sel4[1:0] & as[3:0] & 111 & qx[2] \\
	\hline
\end{tabular}

\textbf{Assembler Syntax} \\
EE.CMUL.S16.ST.INCP qv, as, qz, qx, qy, sel4 \\
\textbf{Description} \\
This instruction performs a 16-bit signed complex multiplication. The range of the immediate number \lstinline{sel4} is 0 \verb+~+ 7, which specifies the 32 bits in the two QR registers \lstinline{qx} and \lstinline{qy} for complex multiplication. The real and imaginary parts of complex numbers are stored in the upper 16 bits and lower 16 bits of the 32 bits respectively. The calculated real part and imaginary part results are stored in the corresponding 32 bits of register \lstinline{qz}. \\ During the operation, the lower 4 bits of the access address in register \lstinline{as} are forced to be 0, and then stores the 16-byte data of \lstinline{qv} to memory. After the access, the value in register \lstinline{as} is incremented by 16.\\
\textbf{Operation}
\begin{lstlisting}
if sel4 == 0:
  qz[ 15:  0] = (qx[ 15:  0] * qy[ 15:  0] - qx[ 31: 16] * qy[ 31: 16]) >> SAR[5:0]
  qz[ 31: 16] = (qx[ 15:  0] * qy[ 31: 16] + qx[ 31: 16] * qy[ 15:  0]) >> SAR[5:0]
  qz[ 47: 32] = (qx[ 47: 32] * qy[ 47: 32] - qx[ 63: 48] * qy[ 63: 48]) >> SAR[5:0]
  qz[ 63: 48] = (qx[ 47: 32] * qy[ 63: 48] + qx[ 63: 48] * qy[ 47: 32]) >> SAR[5:0]
if sel4 == 1:
  qz[ 79: 64] = (qx[ 79: 64] * qy[ 79: 64] - qx[ 95: 80] * qy[ 95: 80]) >> SAR[5:0]
  qz[ 95: 80] = (qx[ 79: 64] * qy[ 95: 80] + qx[ 95: 80] * qy[ 79: 64]) >> SAR[5:0]
  qz[111: 96] = (qx[111: 96] * qy[111: 96] - qx[127:112] * qy[127:112]) >> SAR[5:0]
  qz[127:112] = (qx[111: 96] * qy[127:112] + qx[127:112] * qy[111: 96]) >> SAR[5:0]
if sel4 == 2:
  qz[ 15:  0] = (qx[ 15:  0] * qy[ 15:  0] + qx[ 31: 16] * qy[ 31: 16]) >> SAR[5:0]
  qz[ 31: 16] = (qx[ 15:  0] * qy[ 31: 16] - qx[ 31: 16] * qy[ 15:  0]) >> SAR[5:0]
  qz[ 47: 32] = (qx[ 47: 32] * qy[ 47: 32] + qx[ 63: 48] * qy[ 63: 48]) >> SAR[5:0]
  qz[ 63: 48] = (qx[ 47: 32] * qy[ 63: 48] - qx[ 63: 48] * qy[ 47: 32]) >> SAR[5:0]
if sel4 == 3:
  qz[ 79: 64] = (qx[ 79: 64] * qy[ 79: 64] + qx[ 95: 80] * qy[ 95: 80]) >> SAR[5:0]
  qz[ 95: 80] = (qx[ 79: 64] * qy[ 95: 80] - qx[ 95: 80] * qy[ 79: 64]) >> SAR[5:0]
  qz[111: 96] = (qx[111: 96] * qy[111: 96] + qx[127:112] * qy[127:112]) >> SAR[5:0]
  qz[127:112] = (qx[111: 96] * qy[127:112] - qx[127:112] * qy[111: 96]) >> SAR[5:0]

  qv[127:0] => store128({as[31:4],4{0}})
  as[31:0] = as[31:0] + 16
\end{lstlisting}

\clearpage
\subsection{EE.FFT.AMS.S16.LD.INCP}\label{instruction:EEFFTAMSS16LDINCP}
\setlength\tabcolsep{5pt}
\textbf{Instruction Word} \\
\begin{tabular}{|c|c|c|c|c|c|c|c|c|c|c|c|c|c|}
	\hline
	110100 & sel2[0] & qz1[2] & qz[0] & qy[2:0] & qz1[1] & qm[2:0] & qx[1:0] & qz[2:1] & qz1[0] & qu[2:0] & as[3:0] & 111 & qx[2] \\
	\hline
\end{tabular}

\textbf{Assembler Syntax} \\
EE.FFT.AMS.S16.LD.INCP qu, as, qz, qz1, qx, qy, qm, sel2 \\
\textbf{Description} \\
It is a dedicated FFT instruction to perform addition, subtraction, multiplication, addition and subtraction, and shift operations on 16-bit data segments.\\
During the operation, the lower 4 bits of the access address in register \lstinline{as} are forced to be 0, and then the 16-byte data is loaded from the memory to register \lstinline{qu}. After the access, the value in register \lstinline{as} is incremented by 16.\\
\textbf{Operation}
\begin{lstlisting}
  temp0[15:0] = qx[ 47: 32] + qy[ 47: 32]
  temp1[15:0] = qx[ 63: 48] - qy[ 63: 48]

if sel2==0:
  temp2[15:0] = ((qx[ 47: 32] - qy[ 47: 32]) * qm[ 47: 32] - (qx[ 63: 48] + qy[ 63: 48]) * qm[ 63: 48]) >> SAR
  temp3[15:0] = ((qx[ 47: 32] - qy[ 47: 32]) * qm[ 63: 48] + (qx[ 63: 48] + qy[ 63: 48]) * qm[ 47: 32]) >> SAR
if sel2==1:
  temp2[15:0] = ((qx[ 63: 48] + qy[ 63: 48]) * qm[ 63: 48] + (qx[ 47: 32] - qy[ 47: 32]) * qm[ 47: 32]) >> SAR
  temp3[15:0] = ((qx[ 63: 48] + qy[ 63: 48]) * qm[ 47: 32] - (qx[ 47: 32] - qy[ 47: 32]) * qm[ 63: 48]) >> SAR

  qz[47: 32] = temp0[15:0] + temp2[15:0]
  qz[63: 48] = temp1[15:0] + temp3[15:0]
  qz1[47: 32] = temp0[15:0] - temp2[15:0]
  qz1[63: 48] = temp3[15:0] - temp1[15:0]
  qu = load128({as[31:4],4{0}})
  as[31:0] = as[31:0] + 16
\end{lstlisting}

\clearpage
\subsection{EE.FFT.AMS.S16.LD.INCP.UAUP}\label{instruction:EEFFTAMSS16LDINCPUAUP}
\setlength\tabcolsep{5pt}
\textbf{Instruction Word} \\
\begin{tabular}{|c|c|c|c|c|c|c|c|c|c|c|c|c|c|}
	\hline
	110101 & sel2[0] & qz1[2] & qz[0] & qy[2:0] & qz1[1] & qm[2:0] & qx[1:0] & qz[2:1] & qz1[0] & qu[2:0] & as[3:0] & 111 & qx[2] \\
	\hline
\end{tabular}

\textbf{Assembler Syntax} \\
EE.FFT.AMS.S16.LD.INCP.UAUP qu, as, qz, qz1, qx, qy, qm, sel2 \\
\textbf{Description} \\
It is a dedicated FFT instruction to perform addition, subtraction, multiplication, addition and subtraction, and shift operations on 16-bit data segments.\\
During the operation, the lower 4 bits of the access address in the register \lstinline{as} are forced to be 0, and then the 16-byte data is loaded from the memory. The instruction joins the loaded data and the data in special register \lstinline{UA_STATE} into 32-byte data, right-shifts the 32-byte data by the result of the \lstinline{SAR_BYTE} value multiplied by 8, and assigns the lower 128 bits of the shifted result to register \lstinline{qu}. Meanwhile, register \lstinline{UA_STATE} is updated with the loaded 16-byte data. After the access, the value in register \lstinline{as} is incremented by 16.\\
\textbf{Operation}
\begin{lstlisting}
  temp0[15:0] = qx[ 15:  0] + qy[ 15:  0]
  temp1[15:0] = qx[ 31: 16] - qy[ 31: 16]

if sel2==0:
  temp2[15:0] = ((qx[ 15:  0] - qy[ 15:  0]) * qm[ 15:  0] - (qx[ 31: 16] + qy[ 31: 16]) * qm[ 31: 16]) >> SAR
  temp3[15:0] = ((qx[ 15:  0] - qy[ 15:  0]) * qm[ 31: 16] + (qx[ 31: 16] + qy[ 31: 16]) * qm[ 15:  0]) >> SAR
if sel2==1:
  temp2[15:0] = ((qx[ 31: 16] + qy[ 31: 16]) * qm[ 31: 16] + (qx[ 15:  0] - qy[ 15:  0]) * qm[ 15:  0]) >> SAR
  temp3[15:0] = ((qx[ 31: 16] + qy[ 31: 16]) * qm[ 15:  0] - (qx[ 15:  0] - qy[ 15:  0]) * qm[ 31: 16]) >> SAR

  dataIn[127:0] = load128({as[31:4],4{0}})
  qz[15:  0] = temp0[15:0] + temp2[15:0]
  qz[31: 16] = temp1[15:0] + temp3[15:0]
  qz1[15: 0] = temp0[15:0] - temp2[15:0]
  qz1[31:16] = temp3[15:0] - temp1[15:0]
  qu[127: 0] = {dataIn[127:0], UA_STATE[127:0]} >> {SAR_BYTE[3:0] << 3}
  UA_STATE[127:0] = dataIn[127:0]
  as[31:0] = as[31:0] + 16
\end{lstlisting}

\clearpage
\subsection{EE.FFT.AMS.S16.LD.R32.DECP}\label{instruction:EEFFTAMSS16LDR32DECP}
\setlength\tabcolsep{5pt}
\textbf{Instruction Word} \\
\begin{tabular}{|c|c|c|c|c|c|c|c|c|c|c|c|c|c|}
	\hline
	110110 & sel2[0] & qz1[2] & qz[0] & qy[2:0] & qz1[1] & qm[2:0] & qx[1:0] & qz[2:1] & qz1[0] & qu[2:0] & as[3:0] & 111 & qx[2] \\
	\hline
\end{tabular}

\textbf{Assembler Syntax} \\
EE.FFT.AMS.S16.LD.R32.DECP qu, as, qz, qz1, qx, qy, qm, sel2 \\
\textbf{Description} \\
It is a dedicated FFT instruction to perform addition, subtraction, multiplication, addition and subtraction, and shift operations on 16-bit data segments.\\
During the operation, the instruction forces the lower 4 bits of the access address in register \lstinline{as} to 0 and loads the 16-byte data from the memory to register \lstinline{qu} in the big-endian word order, namely, loads the segment [127: 96] of the data to [31:0] of \lstinline{qu}. After the access is completed, the value in register \lstinline{as} is decreased by 16.\\
\textbf{Operation}
\begin{lstlisting}
  temp0[15:0] = qx[ 79: 64] + qy[ 79: 64]
  temp1[15:0] = qx[ 95: 80] - qy[ 95: 80]

if sel2==0:
  temp2[15:0] = ((qx[ 79: 64] - qy[ 79: 64]) * qm[ 79: 64] - (qx[ 95: 80] + qy[ 95: 80]) * qm[ 95: 80]) >> SAR
  temp3[15:0] = ((qx[ 79: 64] - qy[ 79: 64]) * qm[ 95: 80] + (qx[ 95: 80] + qy[ 95: 80]) * qm[ 79: 64]) >> SAR
if sel2==1:
  temp2[15:0] = ((qx[ 95: 80] + qy[ 95: 80]) * qm[ 95: 80] + (qx[ 79: 64] - qy[ 79: 64]) * qm[ 79: 64]) >> SAR
  temp3[15:0] = ((qx[ 95: 80] + qy[ 95: 80]) * qm[ 79: 64] - (qx[ 79: 64] - qy[ 79: 64]) * qm[ 95: 80]) >> SAR

  qz[79: 64] = temp0[15:0] + temp2[15:0]
  qz[95: 80] = temp1[15:0] + temp3[15:0]
  qz1[79:64] = temp0[15:0] - temp2[15:0]
  qz1[95:80] = temp3[15:0] - temp1[15:0]
  {qu[31: 0], qu[63: 32], qu[95: 64], qu[127: 96]} = load128({as[31:4],4{0}})
  as[31:0] = as[31:0] - 16
\end{lstlisting}

\clearpage
\subsection{EE.FFT.AMS.S16.ST.INCP}\label{instruction:EEFFTAMSS16STINCP}
\setlength\tabcolsep{6pt}
\textbf{Instruction Word} \\
\begin{tabular}{|c|c|c|c|c|c|c|c|c|c|c|c|c|}
	\hline
	10100 & sel2[0] & qz1[2:1] & qm[0] & qv[2:0] & qz1[0] & qx[2:0] & qy[1:0] & qm[2:1] & as[3:0] & at[3:0] & 111 & qy[2] \\
	\hline
\end{tabular}

\textbf{Assembler Syntax} \\
EE.FFT.AMS.S16.ST.INCP qv, qz1, at, as, qx, qy, qm, sel2 \\
\textbf{Description} \\
It is a dedicated FFT instruction to perform addition, subtraction, multiplication, addition and subtraction, and shift operations on 16-bit data segments.\\
During the operation, the instruction forces the lower 4 bits of the access address in register \lstinline{as} to 0, splices the values in registers \lstinline{qv} and \lstinline{at} into 16-byte data, and stores the spliced result to memory. After the access is completed, the value in register \lstinline{as} is incremented by 16. Besides, the operation result is updated to register \lstinline{at}.\\
\textbf{Operation}
\begin{lstlisting}
  temp0[15:0] = qx[111: 96] + qy[111: 96]
  temp1[15:0] = qx[127:112] - qy[127:112]

if sel2==0:
  temp2[15:0] = ((qx[111: 96] - qy[111: 96]) * qm[111: 96] - (qx[127:112] + qy[127:112]) * qm[127:112]) >> SAR[5:0]
  temp3[15:0] = ((qx[111: 96] - qy[111: 96]) * qm[127:112] + (qx[127:112] + qy[127:112]) * qm[111: 96]) >> SAR[5:0]
  {qv[ 95: 80] >> 1, qv[ 79: 64] >> 1, qv[ 63: 48] >> 1, qv[ 47: 32] >> 1, qv[ 31: 16] >> 1, qv[ 15:  0] >> 1, at[31: 16] >> 1, at[15:0] >> 1} => store128({as[31:4],4{0}})
if sel2==1:
  temp2[15:0] = ((qx[127:112] + qy[127:112]) * qm[127:112] + (qx[111: 96] - qy[111: 96]) * qm[111: 96]) >> SAR[5:0]
  temp3[15:0] = ((qx[127:112] + qy[127:112]) * qm[111: 96] - (qx[111: 96] - qy[111: 96]) * qm[127:112]) >> SAR[5:0]
  {qv[ 95: 64], qv[ 63: 32], qv[ 31:  0], at[31:0]} => store128({as[31:4],4{0}})

  temp4[16:0] = temp1[15:0] + temp3[15:0]
  temp5[16:0] = temp0[15:0] + temp2[15:0]
  qz1[111: 96] = temp0[15:0] - temp2[15:0]
  qz1[127:112] = temp3[15:0] - temp1[15:0]
  at = {temp4[15:0], temp5[15:0]}
  as[31:0] = as[31:0] +16
\end{lstlisting}

\clearpage
\subsection{EE.FFT.CMUL.S16.LD.XP}\label{instruction:EEFFTCMULS16LDXP}
\setlength\tabcolsep{6pt}
\textbf{Instruction Word} \\
\begin{tabular}{|c|c|c|c|c|c|c|c|c|c|c|c|}
	\hline
	110111 & sel8[2:1] & qu[0] & qy[2:0] & sel8[0] & qz[2:0] & qx[1:0] & qu[2:1] & ad[3:0] & as[3:0] & 111 & qx[2] \\
	\hline
\end{tabular}

\textbf{Assembler Syntax} \\
EE.FFT.CMUL.S16.LD.XP qu, as, ad, qz, qx, qy, sel8 \\
\textbf{Description} \\
This instruction performs a 16-bit signed complex multiplication. It is similar to \hyperref[instruction:EECMULS16]{EE.CMUL.S16} except that the order of registers \lstinline{qx} and \lstinline{qy} is reversed.\\
During the operation, the lower 4 bits of the access address in register \lstinline{as} are forced to be 0, and then the 16-byte data is loaded from the memory to register \lstinline{qu}. After the access is completed, the value in register \lstinline{as} is incremented by the value in register \lstinline{ad}.\\
\textbf{Operation}
\begin{lstlisting}
if sel8 == 0:
  qz[ 15:  0] = (qx[ 15:  0] * qy[ 15:  0] + qx[ 31: 16] * qy[ 31: 16]) >> SAR[5:0]
  qz[ 31: 16] = (qx[ 31: 16] * qy[ 15:  0] - qx[ 15:  0] * qy[ 31: 16]) >> SAR[5:0]
if sel8 == 1:
  qz[ 15:  0] = (qx[ 15:  0] * qy[ 15:  0] - qx[ 31: 16] * qy[ 31: 16]) >> SAR[5:0]
  qz[ 31: 16] = (qx[ 31: 16] * qy[ 15:  0] + qx[ 15:  0] * qy[ 31: 16]) >> SAR[5:0]
if sel8 == 2:
  qz[ 47: 32] = (qx[ 47: 32] * qy[ 47: 32] + qx[ 63: 48] * qy[ 63: 48]) >> SAR[5:0]
  qz[ 63: 48] = (qx[ 63: 48] * qy[ 47: 32] - qx[ 47: 32] * qy[ 63: 48]) >> SAR[5:0]
if sel8 == 3:
  qz[ 47: 32] = (qx[ 47: 32] * qy[ 47: 32] - qx[ 63: 48] * qy[ 63: 48]) >> SAR[5:0]
  qz[ 63: 48] = (qx[ 63: 48] * qy[ 47: 32] + qx[ 47: 32] * qy[ 63: 48]) >> SAR[5:0]
if sel8 == 4:
  qz[ 79: 64] = (qx[ 79: 64] * qy[ 79: 64] + qx[ 95: 80] * qy[ 95: 80]) >> SAR[5:0]
  qz[ 95: 80] = (qx[ 95: 80] * qy[ 79: 64] - qx[ 79: 64] * qy[ 95: 80]) >> SAR[5:0]
if sel8 == 5:
  qz[ 79: 64] = (qx[ 79: 64] * qy[ 79: 64] - qx[ 95: 80] * qy[ 95: 80]) >> SAR[5:0]
  qz[ 95: 80] = (qx[ 95: 80] * qy[ 79: 64] + qx[ 79: 64] * qy[ 95: 80]) >> SAR[5:0]

  qu[127:0] = load128({as[31:4],4{0}})
  as[31:0] = as[31:0] + ad[31:0]
\end{lstlisting}

\clearpage
\subsection{EE.FFT.CMUL.S16.ST.XP}\label{instruction:EEFFTCMULS16STXP}
\setlength\tabcolsep{5pt}
\textbf{Instruction Word} \\
\begin{tabular}{|c|c|c|c|c|c|c|c|c|c|c|c|c|}
	\hline
	10101 & sar4[1:0] & upd4[1] & sel8[0] & qy[2:0] & upd4[0] & qv[2:0] & qx[1:0] & sel8[2:1] & ad[3:0] & as[3:0] & 111 & qx[2] \\
	\hline
\end{tabular}

\textbf{Assembler Syntax} \\
EE.FFT.CMUL.S16.ST.XP qx, qy, qv, as, ad, sel8, upd4, sar4 \\
\textbf{Description} \\
This instruction performs a 16-bit signed complex multiplication. It is similar to \hyperref[instruction:EECMULS16]{EE.CMUL.S16} except that the order of registers \lstinline{qx} and \lstinline{qy} is reversed.\\
The result of the operation and the data segments in registers \lstinline{qx} and qv specified by the immediate data \lstinline{upd4} are concatenated into 128 bits, which then are written into memory. After the access is completed, the value in register \lstinline{as} is incremented by the value in register \lstinline{ad}.
\textbf{Operation}
\begin{lstlisting}
if sel8 == 6:
  temp[15:0] = (qx[111: 96] * qy[111: 96] + qx[127:112] * qy[127:112]) >> SAR[5:0]
  temp[31:16] = (qx[127:112] * qy[111: 96] - qx[111: 96] * qy[127:112]) >> SAR[5:0]
if sel8 == 7:
  temp[15:0] = (qx[111: 96] * qy[111: 96] - qx[127:112] * qy[127:112]) >> SAR[5:0]
  temp[31:16] = (qx[127:112] * qy[111: 96] + qx[111: 96] * qy[127:112]) >> SAR[5:0]

if upd4 == 0:  // normal radix 2
  {temp[31:0], qv[ 95:  0]} => store128({as[31:4],4{0}})
if upd4 == 1:  // radix2 last second stage
  {
       temp[31:0],
       qv[ 95: 64],
       qx[ 63: 48] >> sar4,
       qx[ 47: 32] >> sar4,
       qx[ 31: 16] >> sar4,
       qx[ 15:  0] >> sar4
  } => store128({as[31:4],4{0}})
if upd4 == 2:  // radix2 last stage
  {
       temp[31:0],
       qx[ 63: 48] >> sar4,
       qx[ 47: 32] >> sar4,
       qv[ 95: 64],
       qx[ 31: 16] >> sar4,
       qx[ 15:  0] >> sar4
  } => store128({as[31:4],4{0}})

  as[31:0] = as[31:0] + ad[31:0]
\end{lstlisting}



\clearpage
\subsection{EE.FFT.R2BF.S16}\label{instruction:EEFFTR2BFS16}
\setlength\tabcolsep{6pt}
\textbf{Instruction Word} \\
\begin{tabular}{|c|c|c|c|c|c|c|c|c|c|c|c|}
	\hline
	11 & qa0[2:1] & 1100 & qa0[0] & qa1[2:0] & qy[2:1] & 0 & sel2[0] & qx[2:1] & qy[0] & qx[0] & 0100 \\
	\hline
\end{tabular}

\textbf{Assembler Syntax} \\
EE.FFT.R2BF.S16 qa0, qa1, qx, qy, sel2 \\
\textbf{Description} \\
This instruction performs the radix-2 butterfly operation on the 16-byte values in registers \lstinline{qx} and \lstinline{qy}, and the data bit width is 16 bits. Some of the calculation results are written to register \lstinline{qa0}, and others are written to register \lstinline{qa1}.\\
\textbf{Operation}
\begin{lstlisting}
if sel2==0:
  op_a[127:  0] = {qy[ 63:  0], qx[ 63:  0]}
  op_b[127:  0] = {qy[127: 64], qx[127: 64]}
if sel2==1:
  op_a[127:  0] = {qy[ 95: 64], qy[ 31:  0], qx[ 95: 64], qx[ 31:  0]}
  op_b[127:  0] = {qy[127: 96], qy[ 63: 32], qx[127: 96], qx[ 63: 32]}

  qa0[ 15:  0] = op_a[ 15:  0] + op_b[ 15:  0]
  qa0[ 31: 16] = op_a[ 31: 16] + op_b[ 31: 16]
  qa0[ 47: 32] = op_a[ 47: 32] + op_b[ 47: 32]
  qa0[ 63: 48] = op_a[ 63: 48] + op_b[ 63: 48]
  qa0[ 79: 64] = op_a[ 15:  0] - op_b[ 15:  0]
  qa0[ 95: 80] = op_a[ 31: 16] - op_b[ 31: 16]
  qa0[111: 96] = op_a[ 47: 32] - op_b[ 47: 32]
  qa0[127:112] = op_a[ 63: 48] - op_b[ 63: 48]

  qa1[ 15:  0] = op_a[ 79: 64] + op_b[ 79: 64]
  qa1[ 31: 16] = op_a[ 95: 80] + op_b[ 95: 80]
  qa1[ 47: 32] = op_a[111: 96] + op_b[111: 96]
  qa1[ 63: 48] = op_a[127:112] + op_b[127:112]
  qa1[ 79: 64] = op_a[ 79: 64] - op_b[ 79: 64]
  qa1[ 95: 80] = op_a[ 95: 80] - op_b[ 95: 80]
  qa1[111: 96] = op_a[111: 96] - op_b[111: 96]
  qa1[127:112] = op_a[127:112] - op_b[127:112]
\end{lstlisting}

\clearpage
\subsection{EE.FFT.R2BF.S16.ST.INCP}\label{instruction:EEFFTR2BFS16STINCP}
\setlength\tabcolsep{6pt}
\textbf{Instruction Word} \\
\begin{tabular}{|c|c|c|c|c|c|c|c|c|c|c|c|}
	\hline
	11101000 & sar4[0] & qy[2:0] & 1 & qa0[2:0] & qx[1:0] & 0 & sar4[1] & 0100 & as[3:0] & 111 & qx[2] \\
	\hline
\end{tabular}

\textbf{Assembler Syntax} \\
EE.FFT.R2BF.S16.ST.INCP qa0, qx, qy, as, 0..3 \\
\textbf{Description} \\
This instruction performs the radix-2 butterfly operation on the 16-byte values in registers \lstinline{qx} and \lstinline{qy}, and the data bit width is 16 bits. Some of the calculation results are written to register \lstinline{qa0}, and others implement an arithmetic shift and are written into the memory address indicated by \lstinline{as}. After the access is completed, the value in register \lstinline{as} is incremented by 16.\\
\textbf{Operation}
\begin{lstlisting}
  qa0[ 15:  0] = qx[ 15:  0] - qy[ 15:  0]
  qa0[ 31: 16] = qx[ 31: 16] - qy[ 31: 16]
  qa0[ 47: 32] = qx[ 47: 32] - qy[ 47: 32]
  qa0[ 63: 48] = qx[ 63: 48] - qy[ 63: 48]
  qa0[ 79: 64] = qx[ 79: 64] - qy[ 79: 64]
  qa0[ 95: 80] = qx[ 95: 80] - qy[ 95: 80]
  qa0[111: 96] = qx[111: 96] - qy[111: 96]
  qa0[127:112] = qx[127:112] - qy[127:112]

  {
    (qx[127:112] + qy[127:112]) >> sar4,
    (qx[111: 96] + qy[111: 96]) >> sar4,
    (qx[ 95: 80] + qy[ 95: 80]) >> sar4,
    (qx[ 79: 64] + qy[ 79: 64]) >> sar4,
    (qx[ 63: 48] + qy[ 63: 48]) >> sar4,
    (qx[ 47: 32] + qy[ 47: 32]) >> sar4,
    (qx[ 31: 16] + qy[ 31: 16]) >> sar4,
    (qx[ 15:  0] + qy[ 15:  0]) >> sar4
  } => store128({as[31:4],4{0}})
  as[31:0] = as[31:0] + 16
\end{lstlisting}

\clearpage
\subsection{EE.FFT.VST.R32.DECP}\label{instruction:EEFFTVSTR32DECP}
\setlength\tabcolsep{6pt}
\textbf{Instruction Word} \\
\begin{tabular}{|c|c|c|c|c|c|c|c|c|}
	\hline
	11 & qv[2:1] & 1101 & qv[0] & 0110 & sar2[0] & 11 & as[3:0] & 0100 \\
	\hline
\end{tabular}

\textbf{Assembler Syntax} \\
EE.FFT.VST.R32.DECP qv, as, sar2 \\
\textbf{Description} \\
It is a dedicated FFT instruction. This instruction divides data in register \lstinline {qv} into 8 segments of 16-bit data and performs an arithmetic right shift on them by 0 or 1 depending on the immediate number \lstinline{sar2}, and finally writes the result to the memory address indicated by register \lstinline{as} in word big-endian order. After the access is completed, the value in register \lstinline{as} is decreased by 16.\\
\textbf{Operation}
\begin{lstlisting}
  {
    qv[ 31: 16] >> sar2,
    qv[ 15:  0] >> sar2,
    qv[ 63: 48] >> sar2,
    qv[ 47: 32] >> sar2,
    qv[ 95: 80] >> sar2,
    qv[ 79: 64] >> sar2,
    qv[127:112] >> sar2,
    qv[111: 96] >> sar2
  } => store128({as[31:4],4{0}})
  as[31:0] = as[31:0] - 16
\end{lstlisting}

\clearpage
\subsection{EE.GET\_GPIO\_IN}\label{instruction:EEGETGPIOIN}
\setlength\tabcolsep{6pt}
\textbf{Instruction Word} \\
\begin{tabular}{|c|c|c|}
	\hline
	0110010100001000 & au[3:0] & 0100 \\
	\hline
\end{tabular}

\textbf{Assembler Syntax} \\
EE.GET\_GPIO\_IN au \\
\textbf{Description} \\
It is a dedicated CPU GPIO instruction to assign the content of GPIO\_IN to the lower 8 bits of register \lstinline{au}.\\
\textbf{Operation}
\begin{lstlisting}
  au = {24'h0, GPIO_IN[7:0]}
\end{lstlisting}

\clearpage
\subsection{EE.LD.128.USAR.IP}\label{instruction:EELD128USARIP}
\setlength\tabcolsep{6pt}
\textbf{Instruction Word} \\
\begin{tabular}{|c|c|c|c|c|c|c|c|}
	\hline
	1 & imm16[7] & qu[2:1] & 0001 & qu[0] & imm16[6:0] & as[3:0] & 0100 \\
	\hline
\end{tabular}

\textbf{Assembler Syntax} \\
EE.LD.128.USAR.IP qu, as, -2048..2032 \\
\textbf{Description} \\
This instruction forces the lower 4 bits of the access address in register \lstinline{as} to 0 and loads 16-byte data from memory to register \lstinline{qu}. Meanwhile, it saves the value of the lower 4 bits in \lstinline{as} to the special register \lstinline{SAR_BYTE}. After the access is completed, the value in register \lstinline{as} is incremented by 8-bit sign-extended constant in the instruction code segment left-shifted by 4.\\
\textbf{Operation}
\begin{lstlisting}
  qu[127:0] = load128({as[31:4],4{0}})
  SAR_BYTE = as[3:0]
  as[31:0] = as[31:0] + {20{imm16[7]},imm16[7:0],4{0}}
\end{lstlisting}

\clearpage
\subsection{EE.LD.128.USAR.XP}\label{instruction:EELD128USARXP}
\setlength\tabcolsep{6pt}
\textbf{Instruction Word} \\
\begin{tabular}{|c|c|c|c|c|c|c|c|}
	\hline
	10 & qu[2:1] & 1101 & qu[0] & 000 & ad[3:0] & as[3:0] & 0100 \\
	\hline
\end{tabular}

\textbf{Assembler Syntax} \\
EE.LD.128.USAR.XP qu, as, ad \\
\textbf{Description} \\
This instruction forces the lower 4 bits of the access address in register \lstinline{as} to 0 and loads 16-byte data from memory to register \lstinline{qu}. Meanwhile, it saves the value of the lower 4 bits in \lstinline{as} to the special register \lstinline{SAR_BYTE}. After the access is completed, the value in register \lstinline{as} is incremented by the value in register \lstinline{ad}.\\
\textbf{Operation}
\begin{lstlisting}
  qu[127:0] = load128({as[31:4],4{0}})
  SAR_BYTE = as[3:0]
  as[31:0] = as[31:0] + ad[31:0]
\end{lstlisting}

\clearpage
\subsection{EE.LD.ACCX.IP}\label{instruction:EELDACCXIP}
\setlength\tabcolsep{6pt}
\textbf{Instruction Word} \\
\begin{tabular}{|c|c|c|c|c|c|}
	\hline
	0 & imm8[7] & 0011100 & imm8[6:0] & as[3:0] & 0100 \\
	\hline
\end{tabular}

\textbf{Assembler Syntax} \\
EE.LD.ACCX.IP as, -1024..1016 \\
\textbf{Description} \\
This instruction forces the lower 3 bits of the access address in register \lstinline{as} to 0, loads 64-bit data from memory, and save its lower 40 bits to the special register \lstinline{ACCX}. After the access is completed, the value in register \lstinline{as} is incremented by 8-bit sign-extended constant in the instruction code segment left-shifted by 3.\\
\textbf{Operation}
\begin{lstlisting}
  ACCX[39:0] = load64({as[31:3],3{0}})
  as[31:0] = as[31:0] + {21{imm8[7]},imm8[7:0],3{0}}
\end{lstlisting}

\clearpage
\subsection{EE.LD.QACC\_H.H.32.IP}\label{instruction:EELDQACCHH32IP}
\setlength\tabcolsep{6pt}
\textbf{Instruction Word} \\
\begin{tabular}{|c|c|c|c|c|c|}
	\hline
	0 & imm4[7] & 0111100 & imm4[6:0] & as[3:0] & 0100 \\
	\hline
\end{tabular}

\textbf{Assembler Syntax} \\
EE.LD.QACC\_H.H.32.IP as, -512..508 \\
\textbf{Description} \\
This instruction forces the lower 2 bits of the access address in register \lstinline{as} to 0 and loads 32-bit data from memory to the special register \lstinline{QACC_H[159:128]}. After the access is completed, the value in register \lstinline{as} is incremented by 8-bit sign-extended constant in the instruction code segment left-shifted by 2.\\
\textbf{Operation}
\begin{lstlisting}
  QACC_H[159:128] = load32({as[31:2],2{0}})
  as[31:0] = as[31:0] + {22{imm4[7]},imm4[7:0],2{0}}
\end{lstlisting}

\clearpage
\subsection{EE.LD.QACC\_H.L.128.IP}\label{instruction:EELDQACCHL128IP}
\setlength\tabcolsep{6pt}
\textbf{Instruction Word} \\
\begin{tabular}{|c|c|c|c|c|c|}
	\hline
	0 & imm16[7] & 0001100 & imm16[6:0] & as[3:0] & 0100 \\
	\hline
\end{tabular}

\textbf{Assembler Syntax} \\
EE.LD.QACC\_H.L.128.IP as, -2048..2032 \\
\textbf{Description} \\
This instruction forces the lower 4 bits of the access address in register \lstinline{as} to 0 and loads 16-byte data from memory to the special register \lstinline{QACC_H[127:0]}. After the access is completed, the value in register \lstinline{as} is incremented by 8-bit sign-extended constant in the instruction code segment left-shifted by 4.\\
\textbf{Operation}
\begin{lstlisting}
  QACC_H[127:  0] = load128({as[31:4],4{0}})
  as[31:0] = as[31:0] + {20{imm16[7]},imm16[7:0],4{0}}
\end{lstlisting}

\clearpage
\subsection{EE.LD.QACC\_L.H.32.IP}\label{instruction:EELDQACCLH32IP}
\setlength\tabcolsep{6pt}
\textbf{Instruction Word} \\
\begin{tabular}{|c|c|c|c|c|c|}
	\hline
	0 & imm4[7] & 0101100 & imm4[6:0] & as[3:0] & 0100 \\
	\hline
\end{tabular}

\textbf{Assembler Syntax} \\
EE.LD.QACC\_L.H.32.IP as, -512..508 \\
\textbf{Description} \\
This instruction forces the lower 2 bits of the access address in register \lstinline{as} to 0 and loads 32-bit data from memory to the special register \lstinline{QACC_L[159:128]}. After the access is completed, the value in register \lstinline{as} is incremented by 8-bit sign-extended constant in the instruction code segment left-shifted by 2.\\
\textbf{Operation}
\begin{lstlisting}
  QACC_L[159:128] = load32({as[31:2],2{0}})
  as[31:0] = as[31:0] + {22{imm4[7]},imm4[7:0],2{0}}
\end{lstlisting}

\clearpage
\subsection{EE.LD.QACC\_L.L.128.IP}\label{instruction:EELDQACCLL128IP}
\setlength\tabcolsep{6pt}
\textbf{Instruction Word} \\
\begin{tabular}{|c|c|c|c|c|c|}
	\hline
	0 & imm16[7] & 0000000 & imm16[6:0] & as[3:0] & 0100 \\
	\hline
\end{tabular}

\textbf{Assembler Syntax} \\
EE.LD.QACC\_L.L.128.IP as, -2048..2032 \\
\textbf{Description} \\
This instruction forces the lower 4 bits of the access address in register \lstinline{as} to 0 and loads 16-byte data from memory to the special register \lstinline{QACC_L[127:0]}. After the access is completed, the value in register \lstinline{as} is incremented by 8-bit sign-extended constant in the instruction code segment left-shifted by 4.\\
\textbf{Operation}
\begin{lstlisting}
  QACC_L[127:0] = load128({as[31:4],4{0}})
  as[31:0] = as[31:0] + {20{imm16[7]},imm16[7:0],4{0}}
\end{lstlisting}

\clearpage
\subsection{EE.LD.UA\_STATE.IP}\label{instruction:EELDUASTATEIP}
\setlength\tabcolsep{6pt}
\textbf{Instruction Word} \\
\begin{tabular}{|c|c|c|c|c|c|}
	\hline
	0 & imm16[7] & 0100000 & imm16[6:0] & as[3:0] & 0100 \\
	\hline
\end{tabular}

\textbf{Assembler Syntax} \\
EE.LD.UA\_STATE.IP as, -2048..2032 \\
\textbf{Description} \\
This instruction forces the lower 4 bits of the access address in register \lstinline{as} to 0 and loads 16-byte data from memory to the special register \lstinline{UA_STATE}. After the access is completed, the value in register \lstinline{as} is incremented by 8-bit sign-extended constant in the instruction code segment left-shifted by 4.\\
\textbf{Operation}
\begin{lstlisting}
  UA_STATE[127:0] = load128({as[31:4],4{0}})
  as[31:0] = as[31:0] + {20{imm16[7]},imm16[7:0],4{0}}
\end{lstlisting}

\clearpage
\subsection{EE.LDF.128.IP}\label{instruction:EELDF128IP}
\setlength\tabcolsep{6pt}
\textbf{Instruction Word} \\
\begin{tabular}{|c|c|c|c|c|c|c|c|c|c|}
	\hline
	10000 & fu3[3:1] & fu0[3:0] & fu3[0] & fu2[3:1] & fu1[3:0] & imm16f[3:0] & as[3:0] & 111 & fu2[0] \\
	\hline
\end{tabular}

\textbf{Assembler Syntax} \\
EE.LDF.128.IP fu3, fu2, fu1, fu0, as, -128..112 \\
\textbf{Description} \\
This instruction forces the lower 4 bits of the access address in register \lstinline{as} to 0, loads 16-byte data from memory, and stores it in order from low bit to high bit to floating-point registers \lstinline{fu0}, \lstinline{fu1}, \lstinline{fu2}, and \lstinline{fu3}. After the access is completed, the value in register \lstinline{as} is incremented by 4-bit sign-extended constant in the instruction code segment left-shifted by 4.\\
\textbf{Operation}
\begin{lstlisting}
  dataIn[127:0] = load128({as[31:4],4{0}})
  fu3 = dataIn[127: 96]
  fu2 = dataIn[ 95: 64]
  fu1 = dataIn[ 63: 32]
  fu0 = dataIn[ 31:  0]
  as[31:0] = as[31:0] + {24{imm16f[3]},imm16f[3:0],4{0}}
\end{lstlisting}

\clearpage
\subsection{EE.LDF.128.XP}\label{instruction:EELDF128XP}
\setlength\tabcolsep{6pt}
\textbf{Instruction Word} \\
\begin{tabular}{|c|c|c|c|c|c|c|c|c|c|}
	\hline
	10001 & fu3[3:1] & fu0[3:0] & fu3[0] & fu2[3:1] & fu1[3:0] & ad[3:0] & as[3:0] & 111 & fu2[0] \\
	\hline
\end{tabular}

\textbf{Assembler Syntax} \\
EE.LDF.128.XP fu3, fu2, fu1, fu0, as, ad \\
\textbf{Description} \\
This instruction forces the lower 4 bits of the access address in register \lstinline{as} to 0, loads 16-byte data from memory, and stores it in order from low bit to high bit to floating-point registers \lstinline{fu0}, \lstinline{fu1}, \lstinline{fu2}, and \lstinline{fu3}. After the access is completed, the value in register \lstinline{as} is incremented by the value in register \lstinline{ad}.\\
\textbf{Operation}
\begin{lstlisting}
  dataIn[127:0] = load128({as[31:4],4{0}})
  fu3 = dataIn[127: 96]
  fu2 = dataIn[ 95: 64]
  fu1 = dataIn[ 63: 32]
  fu0 = dataIn[ 31:  0]
  as[31:0] = as[31:0] + ad[31:0]
\end{lstlisting}

\clearpage
\subsection{EE.LDF.64.IP}\label{instruction:EELDF64IP}
\setlength\tabcolsep{6pt}
\textbf{Instruction Word} \\
\begin{tabular}{|c|c|c|c|c|c|c|c|c|c|}
	\hline
	111000 & imm8[7:6] & fu0[3:0] & imm8[5:2] & fu1[3:0] & 010 & imm8[0] & as[3:0] & 111 & imm8[1] \\
	\hline
\end{tabular}

\textbf{Assembler Syntax} \\
EE.LDF.64.IP fu1, fu0, as, -1024..1016 \\
\textbf{Description} \\
This instruction forces the lower 3 bits of the access address in register \lstinline{as} to 0, loads 64-bit data from memory, and stores it in order from low bit to high bit to floating-point registers \lstinline{fu0} and \lstinline{fu1}. After the access is completed, the value in register \lstinline{as} is incremented by 8-bit sign-extended constant in the instruction code segment left-shifted by 3.\\
\textbf{Operation}
\begin{lstlisting}
  dataIn[63:0] = load64({as[31:3],3{0}})
  fu1 = dataIn[63:32]
  fu0 = dataIn[31: 0]
  as[31:0] = as[31:0] + {21{imm8[7]},imm8[7:0],3{0}}
\end{lstlisting}

\clearpage
\subsection{EE.LDF.64.XP}\label{instruction:EELDF64XP}
\setlength\tabcolsep{6pt}
\textbf{Instruction Word} \\
\begin{tabular}{|c|c|c|c|c|c|}
	\hline
	fu0[3:0] & 0110 & fu1[3:0] & ad[3:0] & as[3:0] & 0000 \\
	\hline
\end{tabular}

\textbf{Assembler Syntax} \\
EE.LDF.64.XP fu1, fu0, as, ad \\
\textbf{Description} \\
This instruction forces the lower 3 bits of the access address in register \lstinline{as} to 0, loads 64-bit data from memory, and stores it in order from low bit to high bit to floating-point registers \lstinline{fu0} and \lstinline{fu1}. After the access is completed, the value in register \lstinline{as} is incremented by the value in register \lstinline{ad}.\\
\textbf{Operation}
\begin{lstlisting}
  dataIn[63:0] = load64({as[31:3],3{0}})
  fu1 = dataIn[63:32]
  fu0 = dataIn[31: 0]
  as[31:0] = as[31:0] + ad[31:0]
\end{lstlisting}

\clearpage
\subsection{EE.LDQA.S16.128.IP}\label{instruction:EELDQAS16128IP}
\setlength\tabcolsep{6pt}
\textbf{Instruction Word} \\
\begin{tabular}{|c|c|c|c|c|c|}
	\hline
	0 & imm16[7] & 0000010 & imm16[6:0] & as[3:0] & 0100 \\
	\hline
\end{tabular}

\textbf{Assembler Syntax} \\
EE.LDQA.S16.128.IP as, -2048..2032 \\
\textbf{Description} \\
This instruction forces the lower 4 bits of the access address in register \lstinline{as} to 0, loads 16-byte data from memory, divides it into 8 segments of 16 bits, sign-extends each segment to 40 bits, and then store the results to the 160-bit special registers \lstinline{QACC_L} and \lstinline{QACC_H} respectively. After the access is completed, the value in register \lstinline{as} is incremented by 8-bit sign-extended constant in the instruction code segment left-shifted by 4.\\
\textbf{Operation}
\begin{lstlisting}
  dataIn[127:0] = load128({as[31:4],4{0}})
  QACC_L[ 39:  0] = {24{dataIn[15]}, dataIn[ 15:  0]}
  QACC_L[ 79: 40] = {24{dataIn[31]}, dataIn[ 31: 16]}
  ...
  QACC_H[ 39:  0] = {24{dataIn[79]}, dataIn[ 79: 64]}
  QACC_H[ 79: 40] = {24{dataIn[95]}, dataIn[ 95: 80]}

  as[31:0] = as[31:0] + {20{imm16[7]},imm16[7:0],4{0}}
\end{lstlisting}

\clearpage
\subsection{EE.LDQA.S16.128.XP}\label{instruction:EELDQAS16128XP}
\setlength\tabcolsep{6pt}
\textbf{Instruction Word} \\
\begin{tabular}{|c|c|c|c|}
	\hline
	011111100100 & ad[3:0] & as[3:0] & 0100 \\
	\hline
\end{tabular}

\textbf{Assembler Syntax} \\
EE.LDQA.S16.128.XP as, ad \\
\textbf{Description} \\
This instruction forces the lower 4 bits of the access address in register \lstinline{as} to 0, loads 16-byte data from memory, divides it into 8 segments of 16 bits, sign-extends each segment to 40 bits, and then store the results to the 160-bit special registers \lstinline{QACC_L} and \lstinline{QACC_H} respectively. After the access is completed, the value in register \lstinline{as} is incremented by the value in register \lstinline{ad}.\\
\textbf{Operation}
\begin{lstlisting}
  dataIn[127:0] = load128({as[31:4],4{0}})
  QACC_L[ 39:  0] = {24{dataIn[15]}, dataIn[ 15:  0]}
  QACC_L[ 79: 40] = {24{dataIn[31]}, dataIn[ 31: 16]}
  ...
  QACC_H[ 39:  0] = {24{dataIn[79]}, dataIn[ 79: 64]}
  QACC_H[ 79: 40] = {24{dataIn[95]}, dataIn[ 95: 80]}
  as[31:0] = as[31:0] + ad[31:0]
\end{lstlisting}

\clearpage
\subsection{EE.LDQA.S8.128.IP}\label{instruction:EELDQAS8128IP}
\setlength\tabcolsep{6pt}
\textbf{Instruction Word} \\
\begin{tabular}{|c|c|c|c|c|c|}
	\hline
	0 & imm16[7] & 0100010 & imm16[6:0] & as[3:0] & 0100 \\
	\hline
\end{tabular}

\textbf{Assembler Syntax} \\
EE.LDQA.S8.128.IP as, -2048..2032 \\
\textbf{Description} \\
This instruction forces the lower 4 bits of the access address in register \lstinline{as} to 0, loads 16-byte data from memory, divides it into 16 segments of 8 bits, sign-extends each segment to 20 bits, and then store the results to the 160-bit special registers \lstinline{QACC_L} and \lstinline{QACC_H} respectively. After the access is completed, the value in register \lstinline{as} is incremented by 8-bit sign-extended constant in the instruction code segment left-shifted by 4.\\
\textbf{Operation}
\begin{lstlisting}
  dataIn[127:0] = load128({as[31:4],4{0}})
  QACC_L[19:0] = {12{dataIn[7]}, dataIn[7:0]}
  QACC_L[39:20] = {12{dataIn[15]}, dataIn[15:8]}
  ...
  QACC_H[159:140] = {12{dataIn[127]}, dataIn[127:120]}
  as[31:0] = as[31:0] + {20{imm16[7]},imm16[7:0],4{0}}
\end{lstlisting}

\clearpage
\subsection{EE.LDQA.S8.128.XP}\label{instruction:EELDQAS8128XP}
\setlength\tabcolsep{6pt}
\textbf{Instruction Word} \\
\begin{tabular}{|c|c|c|c|}
	\hline
	011100010100 & ad[3:0] & as[3:0] & 0100 \\
	\hline
\end{tabular}

\textbf{Assembler Syntax} \\
EE.LDQA.S8.128.XP as, ad \\
\textbf{Description} \\
This instruction forces the lower 4 bits of the access address in register \lstinline{as} to 0, loads 16-byte data from memory, divides it into 16 segments of 8 bits, sign-extends each segment to 20 bits, and then store the results to the 160-bit special registers \lstinline{QACC_L} and \lstinline{QACC_H} respectively. After the access is completed, the value in register \lstinline{as} is incremented by the value in register \lstinline{ad}.\\
\textbf{Operation}
\begin{lstlisting}
  dataIn[127:0] = load128({as[31:4],4{0}})
  QACC_L[19:0] = {12{dataIn[7]}, dataIn[7:0]}
  QACC_L[39:20] = {12{dataIn[15]}, dataIn[15:8]}
  ...
  QACC_H[159:140] = {12{dataIn[127]}, dataIn[127:120]}
  as[31:0] = as[31:0] + ad[31:0]
\end{lstlisting}

\clearpage
\subsection{EE.LDQA.U16.128.IP}\label{instruction:EELDQAU16128IP}
\setlength\tabcolsep{6pt}
\textbf{Instruction Word} \\
\begin{tabular}{|c|c|c|c|c|c|}
	\hline
	0 & imm16[7] & 0001010 & imm16[6:0] & as[3:0] & 0100 \\
	\hline
\end{tabular}

\textbf{Assembler Syntax} \\
EE.LDQA.U16.128.IP as, -2048..2032 \\
\textbf{Description} \\
This instruction forces the lower 4 bits of the access address in register \lstinline{as} to 0, loads 16-byte data from memory, divides it into 8 segments of 16 bits, zero-extends each segment to 40 bits, and then store the results to the 160-bit special registers \lstinline{QACC_L} and \lstinline{QACC_H} respectively. After the access is completed, the value in register \lstinline{as} is incremented by 8-bit sign-extended constant in the instruction code segment left-shifted by 4.\\
\textbf{Operation}
\begin{lstlisting}
  dataIn[127:0] = load128({as[31:4],4{0}})
  QACC_L[ 39:  0] = {24{0}, dataIn[ 15:  0]}
  QACC_L[ 79: 40] = {24{0}, dataIn[ 31: 16]}
  QACC_L[119: 80] = {24{0}, dataIn[ 47: 32]}
  QACC_L[159:120] = {24{0}, dataIn[ 63: 48]}
  QACC_H[ 39:  0] = {24{0}, dataIn[ 79: 64]}
  QACC_H[ 79: 40] = {24{0}, dataIn[ 95: 80]}
  QACC_H[119: 80] = {24{0}, dataIn[111: 96]}
  QACC_H[159:120] = {24{0}, dataIn[127:112]}
  as[31:0] = as[31:0] + {20{imm16[7]},imm16[7:0],4{0}}
\end{lstlisting}

\clearpage
\subsection{EE.LDQA.U16.128.XP}\label{instruction:EELDQAU16128XP}
\setlength\tabcolsep{6pt}
\textbf{Instruction Word} \\
\begin{tabular}{|c|c|c|c|}
	\hline
	011110100100 & ad[3:0] & as[3:0] & 0100 \\
	\hline
\end{tabular}

\textbf{Assembler Syntax} \\
EE.LDQA.U16.128.XP as, ad \\
\textbf{Description} \\
This instruction forces the lower 4 bits of the access address in register \lstinline{as} to 0, loads 16-byte data from memory, divides it into 8 segments of 16 bits, zero-extends each segment to 40 bits, and then store the results to the 160-bit special registers \lstinline{QACC_L} and \lstinline{QACC_H} respectively. After the access is completed, the value in register \lstinline{as} is incremented by the value in register \lstinline{ad}.\\
\textbf{Operation}
\begin{lstlisting}
  dataIn[127:0] = load128({as[31:4],4{0}})
  QACC_L[ 39:  0] = {24{0}, dataIn[ 15:  0]}
  ...
  QACC_H[159:120] = {24{0}, dataIn[127:112]}
  as[31:0] = as[31:0] + ad[31:0]
\end{lstlisting}

\clearpage
\subsection{EE.LDQA.U8.128.IP}\label{instruction:EELDQAU8128IP}
\setlength\tabcolsep{6pt}
\textbf{Instruction Word} \\
\begin{tabular}{|c|c|c|c|c|c|}
	\hline
	0 & imm16[7] & 0101010 & imm16[6:0] & as[3:0] & 0100 \\
	\hline
\end{tabular}

\textbf{Assembler Syntax} \\
EE.LDQA.U8.128.IP as, -2048..2032 \\
\textbf{Description} \\
This instruction forces the lower 4 bits of the access address in register \lstinline{as} to 0, loads 16-byte data from memory, divides it into 16 segments of 8 bits, zero-extends each segment to 20 bits, and then store the results to the 160-bit special registers \lstinline{QACC_L} and \lstinline{QACC_H} respectively. After the access is completed, the value in register \lstinline{as} is incremented by 8-bit sign-extended constant in the instruction code segment left-shifted by 4.\\
\textbf{Operation}
\begin{lstlisting}
  dataIn[127:0] = load128({as[31:4],4{0}})
  QACC_L[ 19:  0] = {12{0}, dataIn[  7:  0]}
  QACC_L[ 39: 20] = {12{0}, dataIn[ 15:  8]}
  QACC_L[ 59: 40] = {12{0}, dataIn[ 23: 16]}
  ...
  QACC_L[159:140] = {12{0}, dataIn[ 63: 56]}
  QACC_H[ 19:  0] = {12{0}, dataIn[ 71: 64]}
  QACC_H[ 39: 20] = {12{0}, dataIn[ 79: 72]}
  QACC_H[ 59: 40] = {12{0}, dataIn[ 87: 80]}
  ...
  QACC_H[159:140] = {12{0}, dataIn[127:120]}
  as[31:0] = as[31:0] + {20{imm16[7]},imm16[7:0],4{0}}
\end{lstlisting}

\clearpage
\subsection{EE.LDQA.U8.128.XP}\label{instruction:EELDQAU8128XP}
\setlength\tabcolsep{6pt}
\textbf{Instruction Word} \\
\begin{tabular}{|c|c|c|c|}
	\hline
	011100000100 & ad[3:0] & as[3:0] & 0100 \\
	\hline
\end{tabular}

\textbf{Assembler Syntax} \\
EE.LDQA.U8.128.XP as, ad \\
\textbf{Description} \\
This instruction forces the lower 4 bits of the access address in register \lstinline{as} to 0, loads 16-byte data from memory, divides it into 16 segments of 8 bits, zero-extends each segment to 20 bits, and then store the results to the 160-bit special registers \lstinline{QACC_L} and \lstinline{QACC_H} respectively. After the access is completed, the value in register \lstinline{as} is incremented by the value in register \lstinline{ad}.\\
\textbf{Operation}
\begin{lstlisting}
  dataIn[127:0] = load128({as[31:4],4{0}})
  QACC_L[ 19:  0] = {12{0}, dataIn[  7:  0]}
  ...
  QACC_L[159:140] = {12{0}, dataIn[ 63: 56]}
  QACC_H[ 19:  0] = {12{0}, dataIn[ 71: 64]}
  ...
  QACC_H[159:140] = {12{0}, dataIn[127:120]}
  as[31:0] = as[31:0] + ad[31:0]
\end{lstlisting}

\clearpage
\subsection{EE.LDXQ.32}\label{instruction:EELDXQ32}
\setlength\tabcolsep{6pt}
\textbf{Instruction Word} \\
\begin{tabular}{|c|c|c|c|c|c|c|c|c|c|c|c|}
	\hline
	1110000 & sel4[1] & sel8[0] & 111 & sel4[0] & qu[2:0] & qs[1:0] & sel8[2:1] & 1101 & as[3:0] & 111 & qs[2] \\
	\hline
\end{tabular}

\textbf{Assembler Syntax} \\
EE.LDXQ.32 qu, qs, as, 0..3, 0..7 \\
\textbf{Description} \\
This instruction selects one of the 8 segments of 16-bit data in \lstinline{qs} as the addend according to the immediate number \lstinline{sel8}. It adds the addend left-shifted by 2 bits to the value of the address register \lstinline{as}, uses the result as the access address, aligns it to 32 bits (its lower 2 bits are set to 0), and stores loaded data to a 32-bit data segment in register \lstinline{qu} according to the value of the immediate number \lstinline{sel4}. \\
\textbf{Operation}
\begin{lstlisting}
  vaddr0[31:0] = as[31:0] + qs[ 15:  0] * 4
  vaddr1[31:0] = as[31:0] + qs[ 31: 16] * 4
  vaddr2[31:0] = as[31:0] + qs[ 47: 32] * 4
  vaddr3[31:0] = as[31:0] + qs[ 63: 47] * 4
  vaddr4[31:0] = as[31:0] + qs[ 79: 64] * 4
  vaddr5[31:0] = as[31:0] + qs[ 95: 80] * 4
  vaddr6[31:0] = as[31:0] + qs[111: 96] * 4
  vaddr7[31:0] = as[31:0] + qs[127:112] * 4

if sel8 == 0:
  dataIn[31:0] = load32({vaddr0[31:2],2{0}})
if sel8 == 1:
  dataIn[31:0] = load32({vaddr1[31:2],2{0}})
if sel8 == 2:
  dataIn[31:0] = load32({vaddr2[31:2],2{0}})
if sel8 == 3:
  dataIn[31:0] = load32({vaddr3[31:2],2{0}})
if sel8 == 4:
  dataIn[31:0] = load32({vaddr4[31:2],2{0}})
if sel8 == 5:
  dataIn[31:0] = load32({vaddr5[31:2],2{0}})
if sel8 == 6:
  dataIn[31:0] = load32({vaddr6[31:2],2{0}})
if sel8 == 7:
  dataIn[31:0] = load32({vaddr7[31:2],2{0}})

  qu[32*sel4+31:32*sel4] = dataIn[31:0]
\end{lstlisting}

\clearpage
\subsection{EE.MOV.S16.QACC}\label{instruction:EEMOVS16QACC}
\setlength\tabcolsep{6pt}
\textbf{Instruction Word} \\
\begin{tabular}{|c|c|c|c|c|}
	\hline
	11 & qs[2:1] & 1101 & qs[0] & 111111100100100 \\
	\hline
\end{tabular}

\textbf{Assembler Syntax} \\
EE.MOV.S16.QACC qs \\
\textbf{Description} \\
This instruction sign-extends the 8 segments of 16-bit data in register \lstinline{qs} to 40 bits and writes the result to the special registers \lstinline{QACC_H} and \lstinline{QACC_L}.\\
\textbf{Operation}
\begin{lstlisting}
  QACC_L[ 39:  0] = {24{qs[15]}, qs[ 15:  0]}
  QACC_L[ 79: 40] = {24{qs[31]}, qs[ 31: 16]}
  QACC_L[119: 80] = {24{qs[47]}, qs[ 47: 32]}
  QACC_L[159:120] = {24{qs[63]}, qs[ 63: 48]}
  QACC_H[ 39:  0] = {24{qs[79]}, qs[ 79: 64]}
  QACC_H[ 79: 40] = {24{qs[95]}, qs[ 95: 80]}
  QACC_H[119: 80] = {24{qs[79]}, qs[111: 96]}
  QACC_H[159:120] = {24{qs[95]}, qs[127:112]}
\end{lstlisting}

\clearpage
\subsection{EE.MOV.S8.QACC}\label{instruction:EEMOVS8QACC}
\setlength\tabcolsep{6pt}
\textbf{Instruction Word} \\
\begin{tabular}{|c|c|c|c|c|}
	\hline
	11 & qs[2:1] & 1101 & qs[0] & 111111100110100 \\
	\hline
\end{tabular}

\textbf{Assembler Syntax} \\
EE.MOV.S8.QACC qs \\
\textbf{Description} \\
This instruction sign-extends the 16 segments of 8-bit data in the register \lstinline{qs} to 20 bits and writes the result to special registers \lstinline{QACC_H} and \lstinline{QACC_L}.\\
\textbf{Operation}
\begin{lstlisting}
  QACC_L[ 19:  0] = {12{qs[7]}, qs[  7:  0]}
  QACC_L[ 39: 20] = {12{qs[15]}, qs[ 15:  8]}
  ...
  QACC_H[159:140] = {12{qs[127]}, qs[127:120]}
\end{lstlisting}

\clearpage
\subsection{EE.MOV.U16.QACC}\label{instruction:EEMOVU16QACC}
\setlength\tabcolsep{6pt}
\textbf{Instruction Word} \\
\begin{tabular}{|c|c|c|c|c|}
	\hline
	11 & qs[2:1] & 1101 & qs[0] & 111111101100100 \\
	\hline
\end{tabular}

\textbf{Assembler Syntax} \\
EE.MOV.U16.QACC qs \\
\textbf{Description} \\
This instruction zero-extends the 8 segments of 16-bit data in register \lstinline{qs} to 40 bits and writes the result to special registers \lstinline{QACC_H} and \lstinline{QACC_L}.\\
\textbf{Operation}
\begin{lstlisting}
  QACC_L[ 39:  0] = {24{0}, qs[ 15:  0]}
  QACC_L[ 79: 40] = {24{0}, qs[ 31: 16]}
  QACC_L[119: 80] = {24{0}, qs[ 47: 32]}
  QACC_L[159:120] = {24{0}, qs[ 63: 48]}
  QACC_H[ 39:  0] = {24{0}, qs[ 79: 64]}
  QACC_H[ 79: 40] = {24{0}, qs[ 95: 80]}
  QACC_H[119: 80] = {24{0}, qs[111: 96]}
  QACC_H[159:120] = {24{0}, qs[127:112]}
\end{lstlisting}

\clearpage
\subsection{EE.MOV.U8.QACC}\label{instruction:EEMOVU8QACC}
\setlength\tabcolsep{6pt}
\textbf{Instruction Word} \\
\begin{tabular}{|c|c|c|c|c|}
	\hline
	11 & qs[2:1] & 1101 & qs[0] & 111111101110100 \\
	\hline
\end{tabular}

\textbf{Assembler Syntax} \\
EE.MOV.U8.QACC qs \\
\textbf{Description} \\
This instruction zero-extends the 16 segments of 8-bit data in register \lstinline{qs} to 20 bits and writes the result to special registers \lstinline{QACC_H} and \lstinline{QACC_L}.\\
\textbf{Operation}
\begin{lstlisting}
  QACC_L[ 19:  0] = {12{0}, qs[  7:  0]}
  QACC_L[ 39: 20] = {12{0}, qs[ 15:  8]}
  QACC_L[ 59: 40] = {12{0}, qs[ 23: 16]}
  ...
  QACC_L[159:140] = {12{0}, qs[ 63: 56]}
  QACC_H[ 19:  0] = {12{0}, qs[ 71: 64]}
  QACC_H[ 39: 20] = {12{0}, qs[ 79: 72]}
  QACC_H[ 59: 40] = {12{0}, qs[ 87: 80]}
  ...
  QACC_H[159:140] = {12{0}, qs[127:120]}
\end{lstlisting}

\clearpage
\subsection{EE.MOVI.32.A}\label{instruction:EEMOVI32A}
\setlength\tabcolsep{6pt}
\textbf{Instruction Word} \\
\begin{tabular}{|c|c|c|c|c|c|c|c|c|}
	\hline
	11 & qs[2:1] & 1101 & qs[0] & 111 & sel4[1:0] & 01 & au[3:0] & 0100 \\
	\hline
\end{tabular}

\textbf{Assembler Syntax} \\
EE.MOVI.32.A qs, au, 0..3 \\
\textbf{Description} \\
This instruction selects one data segment of 32 bits from register \lstinline{qs} according to immediate number \lstinline{sel4} and assigns it to register \lstinline{au}.\\
\textbf{Operation}
\begin{lstlisting}
if sel4 == 0:
  au = qs[ 31:  0]
if sel4 == 1:
  au = qs[ 63: 32]
if sel4 == 2:
  au = qs[ 95: 64]
if sel4 == 3:
  au = qs[127: 96]
\end{lstlisting}

\clearpage
\subsection{EE.MOVI.32.Q}\label{instruction:EEMOVI32Q}
\setlength\tabcolsep{6pt}
\textbf{Instruction Word} \\
\begin{tabular}{|c|c|c|c|c|c|c|c|c|}
	\hline
	11 & qu[2:1] & 1101 & qu[0] & 011 & sel4[1:0] & 10 & as[3:0] & 0100 \\
	\hline
\end{tabular}

\textbf{Assembler Syntax} \\
EE.MOVI.32.Q qu, as, 0..3 \\
\textbf{Description} \\
This instruction assigns the value in register \lstinline{as} to one data segment of 32 bits in register \lstinline{qu} according to immediate number \lstinline{sel4}.\\
\textbf{Operation}
\begin{lstlisting}
if sel4 == 0:
  qu[ 31:  0] = as
if sel4 == 1:
  qu[ 63: 32] = as
if sel4 == 2:
  qu[ 95: 64] = as
if sel4 == 3:
  qu[127: 96] = as
\end{lstlisting}

\clearpage
\subsection{EE.NOTQ}\label{instruction:EENOTQ}
\setlength\tabcolsep{6pt}
\textbf{Instruction Word} \\
\begin{tabular}{|c|c|c|c|c|c|c|c|c|}
	\hline
	11 & qa[2:1] & 1101 & qa[0] & 1111111 & qx[2:1] & 0 & qx[0] & 0100 \\
	\hline
\end{tabular}

\textbf{Assembler Syntax} \\
EE.NOTQ qa, qx \\
\textbf{Description} \\
This instruction performs a bitwise NOT operation on register \lstinline{qx} and writes the result to register \lstinline{qa}.\\
\textbf{Operation}
\begin{lstlisting}
  qa = ~qx
\end{lstlisting}

\clearpage
\subsection{EE.ORQ}\label{instruction:EEORQ}
\setlength\tabcolsep{6pt}
\textbf{Instruction Word} \\
\begin{tabular}{|c|c|c|c|c|c|c|c|c|c|c|}
	\hline
	11 & qa[2:1] & 1101 & qa[0] & 111 & qy[2:1] & 00 & qx[2:1] & qy[0] & qx[0] & 0100 \\
	\hline
\end{tabular}

\textbf{Assembler Syntax} \\
EE.ORQ qa, qx, qy \\
\textbf{Description} \\
This instruction performs a bitwise OR operation on registers \lstinline{qx} and \lstinline{qy} and writes the result of the logical operation to register \lstinline{qa}.\\
\textbf{Operation}
\begin{lstlisting}
  qa = qx | qy
\end{lstlisting}

\clearpage
\subsection{EE.SET\_BIT\_GPIO\_OUT}\label{instruction:EESETBITGPIOOUT}
\setlength\tabcolsep{6pt}
\textbf{Instruction Word} \\
\begin{tabular}{|c|c|c|}
	\hline
	011101010100 & imm256[7:0] & 0100 \\
	\hline
\end{tabular}

\textbf{Assembler Syntax} \\
EE.SET\_BIT\_GPIO\_OUT 0..255 \\
\textbf{Description} \\
It is a dedicated CPU GPIO instruction to set certain bits of \lstinline{GPIO_OUT}. The assignment content depends on the 8-bit immediate number \lstinline{imm256}.\\
\textbf{Operation}
\begin{lstlisting}
  GPIO_OUT[7:0] =  (GPIO_OUT[7:0] | imm256[7:0])
\end{lstlisting}

\clearpage
\subsection{EE.SLCI.2Q}\label{instruction:EESLCI2Q}
\setlength\tabcolsep{6pt}
\textbf{Instruction Word} \\
\begin{tabular}{|c|c|c|c|c|c|c|c|}
	\hline
	11 & qs1[2:1] & 1100 & qs1[0] & qs0[2:0] & 0110 & sar16[3:0] & 0100 \\
	\hline
\end{tabular}

\textbf{Assembler Syntax} \\
EE.SLCI.2Q qs1, qs0, 0..15 \\
\textbf{Description} \\
This instruction performs a left shift on the 32-byte concatenation of registers \lstinline{qs0} and \lstinline{qs1} and pads the lower bits with 0. The upper 128 bits of the shift result is written to register \lstinline{qs1} and the lower 128 bits is written to \lstinline{qs0}. The left shift amount is 8 times the sum of \lstinline{sar16} and 1.\\
\textbf{Operation}
\begin{lstlisting}
  {qs1[127:  0], qs0[127:  0]} = {qs1[127:  0], qs0[127:  0]} << ((sar16[3:0]+1)*8)
\end{lstlisting}

\clearpage
\subsection{EE.SLCXXP.2Q}\label{instruction:EESLCXXP2Q}
\setlength\tabcolsep{6pt}
\textbf{Instruction Word} \\
\begin{tabular}{|c|c|c|c|c|c|c|c|}
	\hline
	10 & qs1[2:1] & 0110 & qs1[0] & qs0[2:0] & ad[3:0] & as[3:0] & 0100 \\
	\hline
\end{tabular}

\textbf{Assembler Syntax} \\
EE.SLCXXP.2Q qs1, qs0, as, ad \\
\textbf{Description} \\
This instruction performs a left shift on the 32-byte concatenation of registers \lstinline{qs0} and \lstinline{qs1} and pads the lower bits with 0. The upper 128 bits of the shift result is written to register \lstinline{qs1} and the lower 128 bits is written to \lstinline{qs0}. The left shift amount is 8 multiplied by the sum of 1 plus the lower 4-bit value of register \lstinline{as}. After the above operations, the value in \lstinline{as} is incremented by the value in \lstinline{ad}.\\
\textbf{Operation}
\begin{lstlisting}
  {qs1[127:  0], qs0[127:  0]} = {qs1[127:  0], qs0[127:  0]} << ((as[3:0]+1)*8)
  as[31:0] = as[31:0] + ad[31:0]
\end{lstlisting}

\clearpage
\subsection{EE.SRC.Q}\label{instruction:EESRCQ}
\setlength\tabcolsep{6pt}
\textbf{Instruction Word} \\
\begin{tabular}{|c|c|c|c|c|c|c|c|}
	\hline
	11 & qs1[2:1] & 1100 & qs1[0] & qs0[2:0] & 00110 & qa[2:0] & 0100 \\
	\hline
\end{tabular}

\textbf{Assembler Syntax} \\
EE.SRC.Q qa, qs0, qs1 \\
\textbf{Description} \\
This instruction performs an arithmetic right shift on the 32-byte concatenation of registers \lstinline{qs0} and \lstinline{qs1} that hold the loaded data of two consecutive aligned addresses. By this way, you can obtain unaligned 16-byte data, which will be written to register \lstinline{qa}. The right shift amount is \lstinline{SAR_BYTE} multiplied by 8.\\
\textbf{Operation}
\begin{lstlisting}
  qa[127:  0]  = {qs1[127:  0], qs0[127:  0]} >> {SAR_BYTE[3:0] << 3}
\end{lstlisting}

\clearpage
\subsection{EE.SRC.Q.LD.IP}\label{instruction:EESRCQLDIP}
\setlength\tabcolsep{3.5pt}
\textbf{Instruction Word} \\
\begin{tabular}{|c|c|c|c|c|c|c|c|c|c|c|c|c|}
	\hline
	111000 & imm16[7:6] & imm16[2] & qs1[2:0] & imm16[5] & qu[2:0] & qs0[1:0] & imm16[4:3] & 00 & imm16[1:0] & as[3:0] & 111 & qs0[2] \\
	\hline
\end{tabular}

\textbf{Assembler Syntax} \\
EE.SRC.Q.LD.IP qu, as, -2048..2032, qs0, qs1 \\
\textbf{Description} \\
This instruction performs an arithmetic right shift on the 32-byte concatenation of registers \lstinline{qs0} and \lstinline{qs1} that hold the loaded data of two consecutive aligned addresses. By this way, you can obtain unaligned 16-byte data, which will be written to register \lstinline{qs0}. The right shift amount is \lstinline{SAR_BYTE} multiplied by 8.\\
At the same time, the lower 4 bits of the access address in register \lstinline{as} are forced to be 0, and then the 16-byte data is loaded from the memory to register \lstinline{qu}. After the access is completed, the value in register \lstinline{as} is incremented by 8-bit sign-extended constant in the instruction code segment left-shifted by 4.\\
\textbf{Operation}
\begin{lstlisting}
  qs0[127:0]  = {qs1[127:0], qs0[127:0]} >> {SAR_BYTE[3:0] << 3}
  qu[127:0] = load128({as[31:4],4{0}})
  as[31:0] = as[31:0] + {20{imm16[7]},imm16[7:0],4{0}}
\end{lstlisting}

\clearpage
\subsection{EE.SRC.Q.LD.XP}\label{instruction:EESRCQLDXP}
\setlength\tabcolsep{6pt}
\textbf{Instruction Word} \\
\begin{tabular}{|c|c|c|c|c|c|c|c|c|c|}
	\hline
	111010000 & qs1[2:0] & 0 & qu[2:0] & qs0[1:0] & 00 & ad[3:0] & as[3:0] & 111 & qs0[2] \\
	\hline
\end{tabular}

\textbf{Assembler Syntax} \\
EE.SRC.Q.LD.XP qu, as, ad, qs0, qs1 \\
\textbf{Description} \\
This instruction performs an arithmetic right shift on the 32-byte concatenation of registers \lstinline{qs0} and \lstinline{qs1} that hold the loaded data of two consecutive aligned addresses. By this way, you can obtain unaligned 16-byte data, which will be written to register \lstinline{qs0}. The right shift amount is \lstinline{SAR_BYTE} multiplied by 8.\\
At the same time, the lower 4 bits of the access address in register \lstinline{as} are forced to be 0, and then the 16-byte data is loaded from the memory to register \lstinline{qu}. After the access is completed, the value in register \lstinline{as} is incremented by the value in register \lstinline{ad}.\\
\textbf{Operation}
\begin{lstlisting}
  qs0[127:0] = {qs1[127: 0], qs0[127: 0]} >> {SAR_BYTE[3:0] << 3}
  qu[127:0] = load128({as[31:4],4{0}})
  as[31:0] = as[31:0] + ad[31:0]
\end{lstlisting}

\clearpage
\subsection{EE.SRC.Q.QUP}\label{instruction:EESRCQQUP}
\setlength\tabcolsep{6pt}
\textbf{Instruction Word} \\
\begin{tabular}{|c|c|c|c|c|c|c|c|}
	\hline
	11 & qs1[2:1] & 1100 & qs1[0] & qs0[2:0] & 01110 & qa[2:0] & 0100 \\
	\hline
\end{tabular}

\textbf{Assembler Syntax} \\
EE.SRC.Q.QUP qa, qs0, qs1 \\
\textbf{Description} \\
This instruction performs an arithmetic right shift on the 32-byte concatenation of registers \lstinline{qs0} and \lstinline{qs1} that hold the loaded data of two consecutive aligned addresses. In this way, you can obtain unaligned 16-byte data, which will be written to register \lstinline{qa}. The right shift amount is \lstinline{SAR_BYTE} multiplied by 8. At the same time, the value of register \lstinline{qs1} is updated to \lstinline{qs0}.\\
\textbf{Operation}
\begin{lstlisting}
  qa[127:  0] = {qs1[127:  0], qs0[127:  0]} >> {SAR_BYTE[3:0] << 3}
  qs0 = qs1
\end{lstlisting}

\clearpage
\subsection{EE.SRCI.2Q}\label{instruction:EESRCI2Q}
\setlength\tabcolsep{6pt}
\textbf{Instruction Word} \\
\begin{tabular}{|c|c|c|c|c|c|c|c|}
	\hline
	11 & qs1[2:1] & 1100 & qs1[0] & qs0[2:0] & 1010 & sar16[3:0] & 0100 \\
	\hline
\end{tabular}

\textbf{Assembler Syntax} \\
EE.SRCI.2Q qs1, qs0, sar16 \\
\textbf{Description} \\
This instruction performs a logical right shift on the 32-byte concatenation of registers \lstinline{qs0} and \lstinline{qs1} and pads the higher bits with 0. The upper 128 bits of the shift result is written to register \lstinline{qs1} and the lower 128 bits is written to \lstinline{qs0}. The right shift amount is 8 times the sum of \lstinline{sar16} and 1.\\
\textbf{Operation}
\begin{lstlisting}
  {qs1[127:  0], qs0[127:  0]} = {qs1[127:  0], qs0[127:  0]} >> ((sar16+1)*8)
  qs1[127:127-8*sar16] = 0
\end{lstlisting}

\clearpage
\subsection{EE.SRCMB.S16.QACC}\label{instruction:EESRCMBS16QACC}
\setlength\tabcolsep{6pt}
\textbf{Instruction Word} \\
\begin{tabular}{|c|c|c|c|c|c|c|c|}
	\hline
	11 & qu[2:1] & 1101 & qu[0] & 1110 & 010 & as[3:0] & 0100 \\
	\hline
\end{tabular}

\textbf{Assembler Syntax} \\
EE.SRCMB.S16.QACC qu, as, 0 \\
\textbf{Description} \\
This instruction extracts 8 data segments of 40 bits from special registers \lstinline{QACC_H} and \lstinline{QACC_L} and perform arithmetic right shift operations respectively. While writing the shift result back to \lstinline{QACC_H} and \lstinline{QACC_L}, the instruction saturates the result to 16-bit signed numbers and writes the 8 16-bit data obtained after saturation into register \lstinline{qu}.\\
\textbf{Operation}
\begin{lstlisting}
  temp0[39:0] = QACC_L[ 39:  0]
  ...
  temp3[39:0] = QACC_L[159:120]
  temp4[39:0] = QACC_H[ 39:  0]
  ...
  temp7[39:0] = QACC_H[159:120]

  temp_shf0[39:0] = temp0[39:0] >> as[5:0]
  temp_shf1[39:0] = temp1[39:0] >> as[5:0]
  ...
  temp_shf7[39:0] = temp7[39:0] >> as[5:0]

  QACC_L[ 39:  0] = temp_shf0[39:0]
  ...
  QACC_L[159:120] = temp_shf3[39:0]
  QACC_H[ 39:  0] = temp_shf4[39:0]
  ...
  QACC_H[159:120] = temp_shf7[39:0]

  qu[ 15:  0] = min(max(temp_shf0[39:0], -2^{15}), 2^{15}-1)
  ...
  qu[ 63: 48] = min(max(temp_shf3[39:0], -2^{15}), 2^{15}-1)
  qu[ 79: 64] = min(max(temp_shf4[39:0], -2^{15}), 2^{15}-1)
  ...
  qu[127:112] = min(max(temp_shf7[39:0], -2^{15}), 2^{15}-1)
\end{lstlisting}

\clearpage
\subsection{EE.SRCMB.S8.QACC}\label{instruction:EESRCMBS8QACC}
\setlength\tabcolsep{6pt}
\textbf{Instruction Word} \\
\begin{tabular}{|c|c|c|c|c|c|c|c|}
	\hline
	11 & qu[2:1] & 1101 & qu[0] & 1111 & 010 & as[3:0] & 0100 \\
	\hline
\end{tabular}

\textbf{Assembler Syntax} \\
EE.SRCMB.S8.QACC qu, as, 0 \\
\textbf{Description} \\
This instruction extracts 16 data segments of 20 bits from special registers \lstinline{QACC_H} and \lstinline{QACC_L} and perform arithmetic right shift operations respectively. While writing the shift result back to \lstinline{QACC_H} and \lstinline{QACC_L}, the instruction saturates the result to 8-bit signed numbers and writes the 16 8-bit data obtained after saturation into register \lstinline{qu}.\\
\textbf{Operation}
\begin{lstlisting}
  temp0[19:0] = QACC_L[ 19:  0]
  temp1[19:0] = QACC_L[ 39: 20]
  ...
  temp7[19:0] = QACC_L[159:140]
  temp8[19:0] = QACC_H[ 19:  0]
  temp9[19:0] = QACC_H[ 39: 20]
  ...
  temp15[19:0] = QACC_H[159:140]

  temp_shf0[19:0] = temp0[19:0] >> as[4:0]
  temp_shf1[19:0] = temp1[19:0] >> as[4:0]
  ...
  temp_shf15[19:0] = temp15[19:0] >> as[4:0]

  QACC_L[ 19:  0] = temp_shf0[19:0]
  ...
  QACC_L[159:140] = temp_shf7[19:0]
  QACC_H[ 19:  0] = temp_shf8[19:0]
  ...
  QACC_H[159:140] = temp_shf15[19:0]

  qu[  7:  0] = min(max(temp_shf0[19:0], -2^{7}), 2^{7}-1)
  qu[ 15:  8] = min(max(temp_shf1[19:0], -2^{7}), 2^{7}-1)
  ...
  qu[ 63: 56] = min(max(temp_shf7[19:0], -2^{7}), 2^{7}-1)
  qu[ 71: 64] = min(max(temp_shf8[19:0], -2^{7}), 2^{7}-1)
  qu[ 79: 72] = min(max(temp_shf9[19:0], -2^{7}), 2^{7}-1)
  ...
  qu[127:120] = min(max(temp_shf15[19:0], -2^{7}), 2^{7}-1)
\end{lstlisting}

\clearpage
\subsection{EE.SRCQ.128.ST.INCP}\label{instruction:EESRCQ128STINCP}
\setlength\tabcolsep{6pt}
\textbf{Instruction Word} \\
\begin{tabular}{|c|c|c|c|c|c|c|c|}
	\hline
	11 & qs1[2:1] & 1100 & qs1[0] & qs0[2:0] & 1110 & as[3:0] & 0100 \\
	\hline
\end{tabular}

\textbf{Assembler Syntax} \\
EE.SRCQ.128.ST.INCP qs0, qs1, as \\
\textbf{Description} \\
This instruction performs an arithmetic right shift on the 32-byte concatenation of registers \lstinline{qs0} and \lstinline{qs1}. Then, it writes the lower 128 bits of the shift result to memory. After the access, the value in register \lstinline{as} is incremented by 16.\\
\textbf{Operation}
\begin{lstlisting}
  {qs1[127:  0], qs0[127:  0]} >> {SAR_BYTE[3:0] << 3}  => store128({as[31:4],4{0}})
  as[31:0] = as[31:0] + 16
\end{lstlisting}

\clearpage
\subsection{EE.SRCXXP.2Q}\label{instruction:EESRCXXP2Q}
\setlength\tabcolsep{6pt}
\textbf{Instruction Word} \\
\begin{tabular}{|c|c|c|c|c|c|c|c|}
	\hline
	11 & qs1[2:1] & 0110 & qs1[0] & qs0[2:0] & ad[3:0] & as[3:0] & 0100 \\
	\hline
\end{tabular}

\textbf{Assembler Syntax} \\
EE.SRCXXP.2Q qs1, qs0, as, ad \\
\textbf{Description} \\
This instruction performs a logical right shift on the 32-byte concatenation of registers \lstinline{qs0} and \lstinline{qs1} and pads the higher bits with 0. The upper 128 bits of the shift result is written to register \lstinline{qs1} and the lower 128 bits is written to \lstinline{qs0}. The right shift amount is 8 multiplied by the sum of 1 plus the lower 4-bit value of register \lstinline{as}. After the above operations, the value in \lstinline{as} is incremented by the value in \lstinline{ad}.\\
\textbf{Operation}
\begin{lstlisting}
  {qs1[127:  0], qs0[127:  0]} = {qs1[127:  0], qs0[127:  0]} >> ((as[3:0]+1)*8)
  qs1[127:127-8*as[3:  0]] = 0
  as[31:0] = as[31:0] + ad[31:0]
\end{lstlisting}

\clearpage
\subsection{EE.SRS.ACCX}\label{instruction:EESRSACCX}
\setlength\tabcolsep{6pt}
\textbf{Instruction Word} \\
\begin{tabular}{|c|c|c|c|c|}
	\hline
	011111100 & 001 & au[3:0] & as[3:0] & 0100 \\
	\hline
\end{tabular}

\textbf{Assembler Syntax} \\
EE.SRS.ACCX au, as, 0 \\
\textbf{Description} \\
This instruction performs an arithmetic right shift on special register \lstinline{ACCX}. While writing the shift result back to \lstinline{ACCX}, the instruction saturates shift result to a 32-bit signed number and writes the saturated result into register \lstinline{au}.\\
\textbf{Operation}
\begin{lstlisting}
  temp_shf[39:0] = ACCX[39:0] >> as[5:0]
  ACCX = temp_shf[39:0]
  au = min(max(temp_shf[39:0], -2^{31}), 2^{31}-1)
\end{lstlisting}

\clearpage
\subsection{EE.ST.ACCX.IP}\label{instruction:EESTACCXIP}
\setlength\tabcolsep{6pt}
\textbf{Instruction Word} \\
\begin{tabular}{|c|c|c|c|c|c|}
	\hline
	0 & imm8[7] & 0000100 & imm8[6:0] & as[3:0] & 0100 \\
	\hline
\end{tabular}

\textbf{Assembler Syntax} \\
EE.ST.ACCX.IP as, -512..508 \\
\textbf{Description} \\
This instruction forces the lower 3 bits of the access address in register \lstinline{as} to 0, zero-extends special register \lstinline{ACCX} to 64 bits, and stores the result to memory. After the access is completed, the value in register \lstinline{as} is incremented by 8-bit sign-extended constant in the instruction code segment left-shifted by 3.\\
\textbf{Operation}
\begin{lstlisting}
  {24{0},ACCX[39:0]} => store64({as[31:3],3{0}})
  as += imm8
\end{lstlisting}

\clearpage
\subsection{EE.ST.QACC\_H.H.32.IP}\label{instruction:EESTQACCHH32IP}
\setlength\tabcolsep{6pt}
\textbf{Instruction Word} \\
\begin{tabular}{|c|c|c|c|c|c|}
	\hline
	0 & imm4[7] & 0100100 & imm4[6:0] & as[3:0] & 0100 \\
	\hline
\end{tabular}

\textbf{Assembler Syntax} \\
EE.ST.QACC\_H.H.32.IP as, -512..508 \\
\textbf{Description} \\
This instruction forces the lower 2 bits of the access address in register \lstinline{as} to 0 and stores the upper 32 bits in special register \lstinline{QACC_H} to memory. After the access is completed, the value in register \lstinline{as} is incremented by 8-bit sign-extended constant in the instruction code segment left-shifted by 2.\\
\textbf{Operation}
\begin{lstlisting}
  QACC_H[159:128] => store32({as[31:2],2{0}})
  as[31:0] = as[31:0] + {23{imm4[7]},imm4[7:0],2{0}}
\end{lstlisting}

\clearpage
\subsection{EE.ST.QACC\_H.L.128.IP}\label{instruction:EESTQACCHL128IP}
\setlength\tabcolsep{6pt}
\textbf{Instruction Word} \\
\begin{tabular}{|c|c|c|c|c|c|}
	\hline
	0 & imm16[7] & 0011010 & imm16[6:0] & as[3:0] & 0100 \\
	\hline
\end{tabular}

\textbf{Assembler Syntax} \\
EE.ST.QACC\_H.L.128.IP as, -2048..2032 \\
\textbf{Description} \\
This instruction forces the lower 4 bits of the access address in register \lstinline{as} to 0 and stores the lower 128 bits in special register \lstinline{QACC_H} to memory. After the access is completed, the value in register \lstinline{as} is incremented by 8-bit sign-extended constant in the instruction code segment left-shifted by 4.\\
\textbf{Operation}
\begin{lstlisting}
  QACC_H[127:  0] => store128({as[31:4],4{0}})
  as[31:0] = as[31:0] + {23{imm4[7]},imm4[7:0],2{0}}
\end{lstlisting}

\clearpage
\subsection{EE.ST.QACC\_L.H.32.IP}\label{instruction:EESTQACCLH32IP}
\setlength\tabcolsep{6pt}
\textbf{Instruction Word} \\
\begin{tabular}{|c|c|c|c|c|c|}
	\hline
	0 & imm4[7] & 0111010 & imm4[6:0] & as[3:0] & 0100 \\
	\hline
\end{tabular}

\textbf{Assembler Syntax} \\
EE.ST.QACC\_L.H.32.IP as, -512..508 \\
\textbf{Description} \\
This instruction forces the lower 2 bits of the access address in register \lstinline{as} to 0 and stores the upper 32 bits in special register \lstinline{QACC_L} to memory. After the access is completed, the value in register \lstinline{as} is incremented by 8-bit sign-extended constant in the instruction code segment left-shifted by 2.\\
\textbf{Operation}
\begin{lstlisting}
  QACC_L[159:128] => store32({as[31:2],2{0}})
  as[31:0] = as[31:0] + {23{imm4[7]},imm4[7:0],2{0}}
\end{lstlisting}

\clearpage
\subsection{EE.ST.QACC\_L.L.128.IP}\label{instruction:EESTQACCLL128IP}
\setlength\tabcolsep{6pt}
\textbf{Instruction Word} \\
\begin{tabular}{|c|c|c|c|c|c|}
	\hline
	0 & imm16[7] & 0011000 & imm16[6:0] & as[3:0] & 0100 \\
	\hline
\end{tabular}

\textbf{Assembler Syntax} \\
EE.ST.QACC\_L.L.128.IP as, -2048..2032 \\
\textbf{Description} \\
This instruction forces the lower 4 bits of the access address in register \lstinline{as} to 0 and stores the lower 128 bits in special register \lstinline{QACC_L} to memory. After the access is completed, the value in register \lstinline{as} is incremented by 8-bit sign-extended constant in the instruction code segment left-shifted by 4.\\
\textbf{Operation}
\begin{lstlisting}
  QACC_L[127:0] => store128({as[31:4],4{0}})
  as[31:0] = as[31:0] + {23{imm4[7]},imm4[7:0],2{0}}
\end{lstlisting}

\clearpage
\subsection{EE.ST.UA\_STATE.IP}\label{instruction:EESTUASTATEIP}
\setlength\tabcolsep{6pt}
\textbf{Instruction Word} \\
\begin{tabular}{|c|c|c|c|c|c|}
	\hline
	0 & imm16[7] & 0111000 & imm16[6:0] & as[3:0] & 0100 \\
	\hline
\end{tabular}

\textbf{Assembler Syntax} \\
EE.ST.UA\_STATE.IP as, -2048..2032 \\
\textbf{Description} \\
This instruction forces the lower 4 bits of the access address in register \lstinline{as} to 0 and stores the 128 bits data in special register \lstinline{UA_STATE} to memory. After the access is completed, the value in register \lstinline{as} is incremented by 8-bit sign-extended constant in the instruction code segment left-shifted by 4.\\
\textbf{Operation}
\begin{lstlisting}
  UA_STATE[127:0] => store128({as[31:4],4{0}})
  as[31:0] = as[31:0] + {20{imm16[7]},imm16[7:0],4{0}}
\end{lstlisting}

\clearpage
\subsection{EE.STF.128.IP}\label{instruction:EESTF128IP}
\setlength\tabcolsep{6pt}
\textbf{Instruction Word} \\
\begin{tabular}{|c|c|c|c|c|c|c|c|c|c|}
	\hline
	10010 & fv3[3:1] & fv0[3:0] & fv3[0] & fv2[3:1] & fv1[3:0] & imm16f[3:0] & as[3:0] & 111 & fv2[0] \\
	\hline
\end{tabular}

\textbf{Assembler Syntax} \\
EE.STF.128.IP fv3, fv2, fv1, fv0, as, -128..112 \\
\textbf{Description} \\
This instruction forces the lower 4 bits of the access address in register \lstinline{as} to 0 and stores the 16-byte data concatenated from four floating-point registers \lstinline{fu0}, \lstinline{fu1}, \lstinline{fu2}, and \lstinline{fu3} in order from low bit to high bit to memory. After the access is completed, the value in register \lstinline{as} is incremented by 4-bit sign-extended constant in the instruction code segment left-shifted by 4.\\
\textbf{Operation}
\begin{lstlisting}
  {fv3,fv2,fv1,fv0} => store128({as[31:4],4{0}})
  as[31:0] = as[31:0] + {24{imm16f[3]},imm16f[3:0],4{0}}
\end{lstlisting}

\clearpage
\subsection{EE.STF.128.XP}\label{instruction:EESTF128XP}
\setlength\tabcolsep{6pt}
\textbf{Instruction Word} \\
\begin{tabular}{|c|c|c|c|c|c|c|c|c|c|}
	\hline
	10011 & fv3[3:1] & fv0[3:0] & fv3[0] & fv2[3:1] & fv1[3:0] & ad[3:0] & as[3:0] & 111 & fv2[0] \\
	\hline
\end{tabular}

\textbf{Assembler Syntax} \\
EE.STF.128.XP fv3, fv2, fv1, fv0, as, ad \\
\textbf{Description} \\
This instruction forces the lower 4 bits of the access address in register \lstinline{as} to 0 and stores the 16-byte data concatenated from four floating-point registers \lstinline{fu0}, \lstinline{fu1}, \lstinline{fu2}, and \lstinline{fu3} in order from low bit to high bit to memory. After the access is completed, the value in register \lstinline{as} is incremented by the value in register \lstinline{ad}.\\
\textbf{Operation}
\begin{lstlisting}
  {fv3,fv2,fv1,fv0} => store128({as[31:4],4{0}})
  as[31:0] = as[31:0] + ad[31:0]
\end{lstlisting}

\clearpage
\subsection{EE.STF.64.IP}\label{instruction:EESTF64IP}
\setlength\tabcolsep{6pt}
\textbf{Instruction Word} \\
\begin{tabular}{|c|c|c|c|c|c|c|c|c|c|}
	\hline
	111000 & imm8[7:6] & fv0[3:0] & imm8[5:2] & fv1[3:0] & 011 & imm8[0] & as[3:0] & 111 & imm8[1] \\
	\hline
\end{tabular}

\textbf{Assembler Syntax} \\
EE.STF.64.IP fv1, fv0, as, imm8 \\
\textbf{Description} \\
This instruction forces the lower 3 bits of the access address in register \lstinline{as} to 0 and stores the 64-bit data concatenated from two floating-point registers \lstinline{fu0} and \lstinline{fu1} in order from low bit to high bit to memory. After the access is completed, the value in register \lstinline{as} is incremented by 8-bit sign-extended constant in the instruction code segment left-shifted by 3.\\
\textbf{Operation}
\begin{lstlisting}
  {fv1,fv0} => store64({as[31:3],3{0}})
  as[31:0] = as[31:0] + {21{imm8[7]},imm8[7:0],3{0}}
\end{lstlisting}

\clearpage
\subsection{EE.STF.64.XP}\label{instruction:EESTF64XP}
\setlength\tabcolsep{6pt}
\textbf{Instruction Word} \\
\begin{tabular}{|c|c|c|c|c|c|}
	\hline
	fv0[3:0] & 0111 & fv1[3:0] & ad[3:0] & as[3:0] & 0000 \\
	\hline
\end{tabular}

\textbf{Assembler Syntax} \\
EE.STF.64.XP fv1, fv0, as, ad \\
\textbf{Description} \\
This instruction forces the lower 3 bits of the access address in register \lstinline{as} to 0 and stores the 64-bit data concatenated from two floating-point registers \lstinline{fu0} and \lstinline{fu1} in order from low bit to high bit to memory. After the access is completed, the value in register \lstinline{as} is incremented by the value in register \lstinline{ad}.\\
\textbf{Operation}
\begin{lstlisting}
  {fv1,fv0} => store64({as[31:3],3{0}})
  as[31:0] = as[31:0] + ad[31:0]
\end{lstlisting}

\clearpage
\subsection{EE.STXQ.32}\label{instruction:EESTXQ32}
\setlength\tabcolsep{6pt}
\textbf{Instruction Word} \\
\begin{tabular}{|c|c|c|c|c|c|c|c|c|c|c|c|}
	\hline
	1110011 & sel4[1] & sel8[0] & qv[2:0] & sel4[0] & 000 & qs[1:0] & sel8[2:1] & 0000 & as[3:0] & 111 & qs[2] \\
	\hline
\end{tabular}

\textbf{Assembler Syntax} \\
EE.STXQ.32 qv, qs, as, sel4, sel8 \\
\textbf{Description} \\
This instruction selects one of the four data segments of 32 bits in register \lstinline{qv} according to immediate number \lstinline{sel4}. Then, it selects an addend from the 8 data segments of 16 bits in register \lstinline{qs} according to immediate number \lstinline{sel8}, adds the addend left-shifted by 2 bits to the value of the address register \lstinline{as}, aligns the sum to 32 bits (its lower 2 bits are set to 0), and uses the result as the written address.\\
\textbf{Operation}
\begin{lstlisting}
  vaddr0[31:0] = as[31:0] + qs[ 15:  0] * 4
  vaddr1[31:0] = as[31:0] + qs[ 31: 16] * 4
  vaddr2[31:0] = as[31:0] + qs[ 47: 32] * 4
  vaddr3[31:0] = as[31:0] + qs[ 63: 47] * 4
  vaddr4[31:0] = as[31:0] + qs[ 79: 64] * 4
  vaddr5[31:0] = as[31:0] + qs[ 95: 80] * 4
  vaddr6[31:0] = as[31:0] + qs[111: 96] * 4
  vaddr7[31:0] = as[31:0] + qs[127:112] * 4

if sel8 == 0:
  qv[32*sel4+31:32*sel4] =>store32({vaddr0[31:2],2{0}})
if sel8 == 1:
  qv[32*sel4+31:32*sel4] =>store32({vaddr1[31:2],2{0}})
if sel8 == 2:
  qv[32*sel4+31:32*sel4] =>store32({vaddr2[31:2],2{0}})
if sel8 == 3:
  qv[32*sel4+31:32*sel4] =>store32({vaddr3[31:2],2{0}})
if sel8 == 4:
  qv[32*sel4+31:32*sel4] =>store32({vaddr4[31:2],2{0}})
if sel8 == 5:
  qv[32*sel4+31:32*sel4] =>store32({vaddr5[31:2],2{0}})
if sel8 == 6:
  qv[32*sel4+31:32*sel4] =>store32({vaddr6[31:2],2{0}})
if sel8 == 7:
  qv[32*sel4+31:32*sel4] =>store32({vaddr7[31:2],2{0}})
\end{lstlisting}

\clearpage
\subsection{EE.VADDS.S16}\label{instruction:EEVADDSS16}
\setlength\tabcolsep{6pt}
\textbf{Instruction Word} \\
\begin{tabular}{|c|c|c|c|c|c|c|c|c|}
	\hline
	10 & qa[2:1] & 1110 & qa[0] & qy[2] & 0 & qy[1:0] & qx[2:0] & 01100100 \\
	\hline
\end{tabular}

\textbf{Assembler Syntax} \\
EE.VADDS.S16 qa, qx, qy \\
\textbf{Description} \\
This instruction performs a vector addition on 16-bit data in the two registers \lstinline{qx} and \lstinline{qy}. Then, the 8 results obtained from the calculation are saturated, and the saturated results are written to register \lstinline{qa}.\\
\textbf{Operation}
\begin{lstlisting}
  qa[ 15:  0] = min(max(qx[ 15:  0] + qy[ 15:  0], -2^{15}), 2^{15}-1)
  qa[ 31: 16] = min(max(qx[ 31: 16] + qy[ 31: 16], -2^{15}), 2^{15}-1)
  ...
\end{lstlisting}

\clearpage
\subsection{EE.VADDS.S16.LD.INCP}\label{instruction:EEVADDSS16LDINCP}
\setlength\tabcolsep{6pt}
\textbf{Instruction Word} \\
\begin{tabular}{|c|c|c|c|c|c|c|c|c|c|c|c|}
	\hline
	111000 & qu[2:1] & qy[0] & 010 & qu[0] & qa[2:0] & qx[1:0] & qy[2:1] & 1101 & as[3:0] & 111 & qx[2] \\
	\hline
\end{tabular}

\textbf{Assembler Syntax} \\
EE.VADDS.S16.LD.INCP qu, as, qa, qx, qy \\
\textbf{Description} \\
This instruction performs a vector addition on 16-bit data in the two registers \lstinline{qx} and \lstinline{qy}. Then, the 8 results obtained from the calculation are saturated, and the saturated results are written to register \lstinline{qa}. \\ During the operation, the lower 4 bits of the access address in register \lstinline{as} are forced to be 0, and then the 16-byte data is loaded from the memory to register \lstinline{qu}. After the access, the value in register \lstinline{as} is incremented by 16.\\
\textbf{Operation}
\begin{lstlisting}
  qa[ 15:  0] = min(max(qx[ 15:  0] + qy[ 15:  0], -2^{15}), 2^{15}-1)
  qa[ 31: 16] = min(max(qx[ 31: 16] + qy[ 31: 16], -2^{15}), 2^{15}-1)
  ...

  qu[127:0] = load128({as[31:4],4{0}})
  as[31:0] = as[31:0] + 16
\end{lstlisting}

\clearpage
\subsection{EE.VADDS.S16.ST.INCP}\label{instruction:EEVADDSS16STINCP}
\setlength\tabcolsep{6pt}
\textbf{Instruction Word} \\
\begin{tabular}{|c|c|c|c|c|c|c|c|c|c|c|}
	\hline
	11100100 & qy[0] & qv[2:0] & 1 & qa[2:0] & qx[1:0] & qy[2:1] & 0000 & as[3:0] & 111 & qx[2] \\
	\hline
\end{tabular}

\textbf{Assembler Syntax} \\
EE.VADDS.S16.ST.INCP qv, as, qa, qx, qy \\
\textbf{Description} \\
This instruction performs a vector addition on 16-bit data in the two registers \lstinline{qx} and \lstinline{qy}. Then, the 8 results obtained from the calculation are saturated, and the saturated results are written to register \lstinline{qa}. \\ During the operation, the instruction forces the lower 4 bits of the access address in register \lstinline{as} to 0 and stores the value in register \lstinline{qv} to memory. After the access, the value in register \lstinline{as} is incremented by 16.\\
\textbf{Operation}
\begin{lstlisting}
  qa[ 15:  0] = min(max(qx[ 15:  0] + qy[ 15:  0], -2^{15}), 2^{15}-1)
  qa[ 31: 16] = min(max(qx[ 31: 16] + qy[ 31: 16], -2^{15}), 2^{15}-1)
  ...

  qv[127:0] => store128({as[31:4],4{0}})
  as[31:0] = as[31:0] + 16
\end{lstlisting}

\clearpage
\subsection{EE.VADDS.S32}\label{instruction:EEVADDSS32}
\setlength\tabcolsep{6pt}
\textbf{Instruction Word} \\
\begin{tabular}{|c|c|c|c|c|c|c|c|c|}
	\hline
	10 & qa[2:1] & 1110 & qa[0] & qy[2] & 0 & qy[1:0] & qx[2:0] & 01110100 \\
	\hline
\end{tabular}

\textbf{Assembler Syntax} \\
EE.VADDS.S32 qa, qx, qy \\
\textbf{Description} \\
This instruction performs a vector addition on 32-bit data in the two registers \lstinline{qx} and \lstinline{qy}. Then, the 4 results obtained from the calculation are saturated, and the saturated results are written to register \lstinline{qa}.\\
\textbf{Operation}
\begin{lstlisting}
  qa[ 31:  0] = min(max(qx[ 31:  0] + qy[ 31:  0], -2^{31}), 2^{31}-1)
  qa[ 63: 32] = min(max(qx[ 63: 32] + qy[ 63: 32], -2^{31}), 2^{31}-1)
  ...
\end{lstlisting}

\clearpage
\subsection{EE.VADDS.S32.LD.INCP}\label{instruction:EEVADDSS32LDINCP}
\setlength\tabcolsep{6pt}
\textbf{Instruction Word} \\
\begin{tabular}{|c|c|c|c|c|c|c|c|c|c|c|c|}
	\hline
	111000 & qu[2:1] & qy[0] & 011 & qu[0] & qa[2:0] & qx[1:0] & qy[2:1] & 1101 & as[3:0] & 111 & qx[2] \\
	\hline
\end{tabular}

\textbf{Assembler Syntax} \\
EE.VADDS.S32.LD.INCP qu, as, qa, qx, qy \\
\textbf{Description} \\
This instruction performs a vector addition on 32-bit data in the two registers \lstinline{qx} and \lstinline{qy}. Then, the 4 results obtained from the calculation are saturated, and the saturated results are written to register \lstinline{qa}. \\ During the operation, the lower 4 bits of the access address in register \lstinline{as} are forced to be 0, and then the 16-byte data is loaded from the memory to register \lstinline{qu}. After the access, the value in register \lstinline{as} is incremented by 16.\\
\textbf{Operation}
\begin{lstlisting}
  qa[ 31:  0] = min(max(qx[ 31:  0] + qy[ 31:  0], -2^{31}), 2^{31}-1)
  qa[ 63: 32] = min(max(qx[ 63: 32] + qy[ 63: 32], -2^{31}), 2^{31}-1)
  ...

  qu[127:0] = load128({as[31:4],4{0}})
  as[31:0] = as[31:0] + 16
\end{lstlisting}

\clearpage
\subsection{EE.VADDS.S32.ST.INCP}\label{instruction:EEVADDSS32STINCP}
\setlength\tabcolsep{6pt}
\textbf{Instruction Word} \\
\begin{tabular}{|c|c|c|c|c|c|c|c|c|c|c|}
	\hline
	11100100 & qy[0] & qv[2:0] & 1 & qa[2:0] & qx[1:0] & qy[2:1] & 0001 & as[3:0] & 111 & qx[2] \\
	\hline
\end{tabular}

\textbf{Assembler Syntax} \\
EE.VADDS.S32.ST.INCP qv, as, qa, qx, qy \\
\textbf{Description} \\
This instruction performs a vector addition on 32-bit data in the two registers \lstinline{qx} and \lstinline{qy}. Then, the 4 results obtained from the calculation are saturated, and the saturated results are written to register \lstinline{qa}. \\ During the operation, the instruction forces the lower 4 bits of the access address in register \lstinline{as} to 0 and stores the value in register \lstinline{qv} to memory. After the access, the value in register \lstinline{as} is incremented by 16.\\
\textbf{Operation}
\begin{lstlisting}
  qa[ 31:  0] = min(max(qx[ 31:  0] + qy[ 31:  0], -2^{31}), 2^{31}-1)
  qa[ 63: 32] = min(max(qx[ 63: 32] + qy[ 63: 32], -2^{31}), 2^{31}-1)
  ...

  qv[127:0] => store128({as[31:4],4{0}})
  as[31:0] = as[31:0] + 16
\end{lstlisting}

\clearpage
\subsection{EE.VADDS.S8}\label{instruction:EEVADDSS8}
\setlength\tabcolsep{6pt}
\textbf{Instruction Word} \\
\begin{tabular}{|c|c|c|c|c|c|c|c|c|}
	\hline
	10 & qa[2:1] & 1110 & qa[0] & qy[2] & 0 & qy[1:0] & qx[2:0] & 10000100 \\
	\hline
\end{tabular}

\textbf{Assembler Syntax} \\
EE.VADDS.S8 qa, qx, qy \\
\textbf{Description} \\
This instruction performs a vector addition on 8-bit data in the two registers \lstinline{qx} and \lstinline{qy}. Then, the 16 results obtained from the calculation are saturated, and the saturated results are written to register \lstinline{qa}.\\
\textbf{Operation}
\begin{lstlisting}
  qa[  7:  0] = min(max(qx[  7:  0] + qy[  7:  0], -2^{7}), 2^{7}-1)
  qa[ 15:  8] = min(max(qx[ 15:  8] + qy[ 15:  8], -2^{7}), 2^{7}-1)
  ...
  qa[127:120] = min(max(qx[127:120] + qy[127:120], -2^{7}), 2^{7}-1)
\end{lstlisting}

\clearpage
\subsection{EE.VADDS.S8.LD.INCP}\label{instruction:EEVADDSS8LDINCP}
\setlength\tabcolsep{6pt}
\textbf{Instruction Word} \\
\begin{tabular}{|c|c|c|c|c|c|c|c|c|c|c|c|}
	\hline
	111000 & qu[2:1] & qy[0] & 001 & qu[0] & qa[2:0] & qx[1:0] & qy[2:1] & 1100 & as[3:0] & 111 & qx[2] \\
	\hline
\end{tabular}

\textbf{Assembler Syntax} \\
EE.VADDS.S8.LD.INCP qu, as, qa, qx, qy \\
\textbf{Description} \\
This instruction performs a vector addition on 8-bit data in the two registers \lstinline{qx} and \lstinline{qy}. Then, the 16 results obtained from the calculation are saturated, and the saturated results are written to register \lstinline{qa}. \\ During the operation, the lower 4 bits of the access address in register \lstinline{as} are forced to be 0, and then the 16-byte data is loaded from the memory to register \lstinline{qu}. After the access, the value in register \lstinline{as} is incremented by 16.\\
\textbf{Operation}
\begin{lstlisting}
  qa[  7:  0] = min(max(qx[  7:  0] + qy[  7:  0], -2^{7}), 2^{7}-1)
  qa[ 15:  8] = min(max(qx[ 15:  8] + qy[ 15:  8], -2^{7}), 2^{7}-1)
  ...

  qu[127:0] = load128({as[31:4],4{0}})
  as[31:0] = as[31:0] + 16
\end{lstlisting}

\clearpage
\subsection{EE.VADDS.S8.ST.INCP}\label{instruction:EEVADDSS8STINCP}
\setlength\tabcolsep{6pt}
\textbf{Instruction Word} \\
\begin{tabular}{|c|c|c|c|c|c|c|c|c|c|c|}
	\hline
	11100100 & qy[0] & qv[2:0] & 1 & qa[2:0] & qx[1:0] & qy[2:1] & 0010 & as[3:0] & 111 & qx[2] \\
	\hline
\end{tabular}

\textbf{Assembler Syntax} \\
EE.VADDS.S8.ST.INCP qv, as, qa, qx, qy \\
\textbf{Description} \\
This instruction performs a vector addition on 8-bit data in the two registers \lstinline{qx} and \lstinline{qy}. Then, the 16 results obtained from the calculation are saturated, and the saturated results are written to register \lstinline{qa}. \\ During the operation, the instruction forces the lower 4 bits of the access address in register \lstinline{as} to 0 and stores the value in register \lstinline{qv} to memory. After the access, the value in register \lstinline{as} is incremented by 16.\\
\textbf{Operation}
\begin{lstlisting}
  qa[  7:  0] = min(max(qx[  7:  0] + qy[  7:  0], -2^{7}), 2^{7}-1)
  qa[ 15:  8] = min(max(qx[ 15:  8] + qy[ 15:  8], -2^{7}), 2^{7}-1)
  ...

  qv[127:0] => store128({as[31:4],4{0}})
  as[31:0] = as[31:0] + 16
\end{lstlisting}

\clearpage
\subsection{EE.VCMP.EQ.S16}\label{instruction:EEVCMPEQS16}
\setlength\tabcolsep{6pt}
\textbf{Instruction Word} \\
\begin{tabular}{|c|c|c|c|c|c|c|c|c|}
	\hline
	10 & qa[2:1] & 1110 & qa[0] & qy[2] & 0 & qy[1:0] & qx[2:0] & 10010100 \\
	\hline
\end{tabular}

\textbf{Assembler Syntax} \\
EE.VCMP.EQ.S16 qa, qx, qy \\
\textbf{Description} \\
This instruction compares 16-bit vector data. It compares the numerical values of the eight 16-bit data segments in registers \lstinline{qx} and \lstinline{qy}. If the values are equal, it writes 0x{}FFFF into the corresponding 16-bit data segment in register \lstinline{qa}. Otherwise, it writes 0 to the segment.\\
\textbf{Operation}
\begin{lstlisting}[escapeinside=`?]
  qa[ 15:  0] = (qx[ 15:  0]==qy[ 15:  0]) ? 0x{}FFFF : 0
  qa[ 31: 16] = (qx[ 31: 16]==qy[ 31: 16]) ? 0x{}FFFF : 0
  ...
  qa[127:112] = (qx[127:112]==qy[127:112]) ? 0x{}FFFF : 0
\end{lstlisting}

\clearpage
\subsection{EE.VCMP.EQ.S32}\label{instruction:EEVCMPEQS32}
\setlength\tabcolsep{6pt}
\textbf{Instruction Word} \\
\begin{tabular}{|c|c|c|c|c|c|c|c|c|}
	\hline
	10 & qa[2:1] & 1110 & qa[0] & qy[2] & 0 & qy[1:0] & qx[2:0] & 10100100 \\
	\hline
\end{tabular}

\textbf{Assembler Syntax} \\
EE.VCMP.EQ.S32 qa, qx, qy \\
\textbf{Description} \\
This instruction compares 32-bit vector data. It compares the numerical values of the four 32-bit data segments in registers \lstinline{qx} and \lstinline{qy}. If the values are equal, it writes 0x{}FFFFFFFF into the corresponding 32-bit data segment in register \lstinline{qa}. Otherwise, it writes 0 to the segment.\\
\textbf{Operation}
\begin{lstlisting}[escapeinside=`?]
  qa[ 31:  0] = (qx[ 31:  0]==qy[ 31:  0]) ? 0x{}FFFFFFFF : 0
  qa[ 63: 32] = (qx[ 63: 32]==qy[ 63: 32]) ? 0x{}FFFFFFFF : 0
  ...
  qa[127: 96] = (qx[127: 96]==qy[127: 96]) ? 0x{}FFFFFFFF : 0
\end{lstlisting}

\clearpage
\subsection{EE.VCMP.EQ.S8}\label{instruction:EEVCMPEQS8}
\setlength\tabcolsep{6pt}
\textbf{Instruction Word} \\
\begin{tabular}{|c|c|c|c|c|c|c|c|c|}
	\hline
	10 & qa[2:1] & 1110 & qa[0] & qy[2] & 0 & qy[1:0] & qx[2:0] & 10110100 \\
	\hline
\end{tabular}

\textbf{Assembler Syntax} \\
EE.VCMP.EQ.S8 qa, qx, qy \\
\textbf{Description} \\
This instruction compares 8-bit vector data. It compares the numerical values of the 16 8-bit data segments in registers \lstinline{qx} and \lstinline{qy}. If the values are equal, it writes 0x{}FF into the corresponding 8-bit data segment in register \lstinline{qa}. Otherwise, it writes 0 to the segment.\\
\textbf{Operation}
\begin{lstlisting}[escapeinside=`?]
  qa[  7:  0] = (qx[  7:  0]==qy[  7:  0]) ? 0x{}FF : 0
  qa[ 15:  8] = (qx[ 15:  8]==qy[ 15:  8]) ? 0x{}FF : 0
  ...
  qa[127:120] = (qx[127:120]==qy[127:120]) ? 0x{}FF : 0
\end{lstlisting}

\clearpage
\subsection{EE.VCMP.GT.S16}\label{instruction:EEVCMPGTS16}
\setlength\tabcolsep{6pt}
\textbf{Instruction Word} \\
\begin{tabular}{|c|c|c|c|c|c|c|c|c|}
	\hline
	10 & qa[2:1] & 1110 & qa[0] & qy[2] & 0 & qy[1:0] & qx[2:0] & 11000100 \\
	\hline
\end{tabular}

\textbf{Assembler Syntax} \\
EE.VCMP.GT.S16 qa, qx, qy \\
\textbf{Description} \\
This instruction compares 16-bit vector data. It compares the numerical values of the eight 16-bit data segments in registers \lstinline{qx} and \lstinline{qy}. If the former is larger than the latter, it writes 0x{}FFFF into the corresponding 16-bit data segment in register \lstinline{qa}. Otherwise, it writes 0 to the segment.\\
\textbf{Operation}
\begin{lstlisting}[escapeinside=`?]
  qa[ 15:  0] = (qx[ 15:  0]>qy[ 15:  0]) ? 0x{}FFFF : 0
  qa[ 31: 16] = (qx[ 31: 16]>qy[ 31: 16]) ? 0x{}FFFF : 0
  ...
  qa[127:112] = (qx[127:112]>qy[127:112]) ? 0x{}FFFF : 0
\end{lstlisting}

\clearpage
\subsection{EE.VCMP.GT.S32}\label{instruction:EEVCMPGTS32}
\setlength\tabcolsep{6pt}
\textbf{Instruction Word} \\
\begin{tabular}{|c|c|c|c|c|c|c|c|c|}
	\hline
	10 & qa[2:1] & 1110 & qa[0] & qy[2] & 0 & qy[1:0] & qx[2:0] & 11010100 \\
	\hline
\end{tabular}

\textbf{Assembler Syntax} \\
EE.VCMP.GT.S32 qa, qx, qy \\
\textbf{Description} \\
This instruction compares 32-bit vector data. It compares the numerical values of the four 32-bit data segments in registers \lstinline{qx} and \lstinline{qy}. If the former is larger than the latter, it writes 0x{}FFFFFFFF into the corresponding 32-bit data segment in register \lstinline{qa}. Otherwise, it writes 0 to the segment.\\
\textbf{Operation}
\begin{lstlisting}[escapeinside=`?]
  qa[ 31:  0] = (qx[ 31:  0]>qy[ 31:  0]) ? 0x{}FFFFFFFF : 0
  qa[ 63: 32] = (qx[ 63: 32]>qy[ 63: 32]) ? 0x{}FFFFFFFF : 0
  ...
  qa[127: 96] = (qx[127: 96]>qy[127: 96]) ? 0x{}FFFFFFFF : 0
\end{lstlisting}

\clearpage
\subsection{EE.VCMP.GT.S8}\label{instruction:EEVCMPGTS8}
\setlength\tabcolsep{6pt}
\textbf{Instruction Word} \\
\begin{tabular}{|c|c|c|c|c|c|c|c|c|}
	\hline
	10 & qa[2:1] & 1110 & qa[0] & qy[2] & 0 & qy[1:0] & qx[2:0] & 11100100 \\
	\hline
\end{tabular}

\textbf{Assembler Syntax} \\
EE.VCMP.GT.S8 qa, qx, qy \\
\textbf{Description} \\
This instruction compares 8-bit vector data. It compares the numerical values of the 16 8-bit data segments in registers \lstinline{qx} and \lstinline{qy}. If the former is larger than the latter, it writes 0x{}FF into the corresponding 8-bit data segment in register \lstinline{qa}. Otherwise, it writes 0 to the segment.\\
\textbf{Operation}
\begin{lstlisting}[escapeinside=`?]
  qa[  7:  0] = (qx[  7:  0]>qy[  7:  0]) ? 0x{}FF : 0
  qa[ 15:  8] = (qx[ 15:  8]>qy[ 15:  8]) ? 0x{}FF : 0
  ...
  qa[127:120] = (qx[127:120]>qy[127:120]) ? 0x{}FF : 0
\end{lstlisting}

\clearpage
\subsection{EE.VCMP.LT.S16}\label{instruction:EEVCMPLTS16}
\setlength\tabcolsep{6pt}
\textbf{Instruction Word} \\
\begin{tabular}{|c|c|c|c|c|c|c|c|c|}
	\hline
	10 & qa[2:1] & 1110 & qa[0] & qy[2] & 0 & qy[1:0] & qx[2:0] & 11110100 \\
	\hline
\end{tabular}

\textbf{Assembler Syntax} \\
EE.VCMP.LT.S16 qa, qx, qy \\
\textbf{Description} \\
This instruction compares 16-bit vector data. It compares the numerical values of the eight 16-bit data segments in registers \lstinline{qx} and \lstinline{qy}. If the former is smaller than the latter, it writes 0x{}FFFF into the corresponding 16-bit data segment in register \lstinline{qa}. Otherwise, it writes 0 to the segment.\\
\textbf{Operation}
\begin{lstlisting}[escapeinside=`?]
  qa[ 15:  0] = (qx[ 15:  0]<qy[ 15:  0]) ? 0x{}FFFF : 0
  qa[ 31: 16] = (qx[ 31: 16]<qy[ 31: 16]) ? 0x{}FFFF : 0
  ...
  qa[127:112] = (qx[127:112]<qy[127:112]) ? 0x{}FFFF : 0
\end{lstlisting}

\clearpage
\subsection{EE.VCMP.LT.S32}\label{instruction:EEVCMPLTS32}
\setlength\tabcolsep{6pt}
\textbf{Instruction Word} \\
\begin{tabular}{|c|c|c|c|c|c|c|c|c|}
	\hline
	10 & qa[2:1] & 1110 & qa[0] & qy[2] & 1 & qy[1:0] & qx[2:0] & 00000100 \\
	\hline
\end{tabular}

\textbf{Assembler Syntax} \\
EE.VCMP.LT.S32 qa, qx, qy \\
\textbf{Description} \\
This instruction compares 32-bit vector data. It compares the numerical values of the four 32-bit data segments in registers \lstinline{qx} and \lstinline{qy}. If the former is smaller than the latter, it writes 0x{}FFFFFFFF into the corresponding 32-bit data segment in register \lstinline{qa}. Otherwise, it writes 0 to the segment.\\
\textbf{Operation}
\begin{lstlisting}[escapeinside=`?]
  qa[ 31:  0] = (qx[ 31:  0]<qy[ 31:  0]) ? 0x{}FFFFFFFF : 0
  qa[ 63: 32] = (qx[ 63: 32]<qy[ 63: 32]) ? 0x{}FFFFFFFF : 0
  ...
  qa[127: 96] = (qx[127: 96]<qy[127: 96]) ? 0x{}FFFFFFFF : 0
\end{lstlisting}

\clearpage
\subsection{EE.VCMP.LT.S8}\label{instruction:EEVCMPLTS8}
\setlength\tabcolsep{6pt}
\textbf{Instruction Word} \\
\begin{tabular}{|c|c|c|c|c|c|c|c|c|}
	\hline
	10 & qa[2:1] & 1110 & qa[0] & qy[2] & 1 & qy[1:0] & qx[2:0] & 00010100 \\
	\hline
\end{tabular}

\textbf{Assembler Syntax} \\
EE.VCMP.LT.S8 qa, qx, qy \\
\textbf{Description} \\
This instruction compares 8-bit vector data. It compares the numerical values of the 16 8-bit data segments in registers \lstinline{qx} and \lstinline{qy}. If the former is smaller than the latter, it writes 0x{}FF into the corresponding 8-bit data segment in register \lstinline{qa}. Otherwise, it writes 0 to the segment.\\
\textbf{Operation}
\begin{lstlisting}[escapeinside=`?]
  qa[  7:  0] = (qx[  7:  0]<qy[  7:  0]) ? 0x{}FF : 0
  qa[ 15:  8] = (qx[ 15:  8]<qy[ 15:  8]) ? 0x{}FF : 0
  ...
  qa[127:120] = (qx[127:120]<qy[127:120]) ? 0x{}FF : 0
\end{lstlisting}

\clearpage
\subsection{EE.VLD.128.IP}\label{instruction:EEVLD128IP}
\setlength\tabcolsep{6pt}
\textbf{Instruction Word} \\
\begin{tabular}{|c|c|c|c|c|c|c|c|}
	\hline
	1 & imm16[7] & qu[2:1] & 0011 & qu[0] & imm16[6:0] & as[3:0] & 0100 \\
	\hline
\end{tabular}

\textbf{Assembler Syntax} \\
EE.VLD.128.IP qu, as, -2048..2032 \\
\textbf{Description} \\
This instruction forces the lower 4 bits of the access address in register \lstinline{as} to 0 and loads 16-byte data from memory to register \lstinline{qu}. After the access is completed, the value in register \lstinline{as} is incremented by 8-bit sign-extended constant in the instruction code segment left-shifted by 4.\\
\textbf{Operation}
\begin{lstlisting}
  qu[127:0] = load128({as[31:4],4{0}})
  as[31:0] = as[31:0] + {20{imm16[7]},imm16[7:0],4{0}}
\end{lstlisting}

\clearpage
\subsection{EE.VLD.128.XP}\label{instruction:EEVLD128XP}
\setlength\tabcolsep{6pt}
\textbf{Instruction Word} \\
\begin{tabular}{|c|c|c|c|c|c|c|c|}
	\hline
	10 & qu[2:1] & 1101 & qu[0] & 010 & ad[3:0] & as[3:0] & 0100 \\
	\hline
\end{tabular}

\textbf{Assembler Syntax} \\
EE.VLD.128.XP qu, as, ad \\
\textbf{Description} \\
This instruction forces the lower 4 bits of the access address in register \lstinline{as} to 0 and loads 16-byte data from memory to register \lstinline{qu}. After the access is completed, the value in register \lstinline{as} is incremented by the value in register \lstinline{ad}.\\
\textbf{Operation}
\begin{lstlisting}
  qu[127:0] = load128({as[31:4],4{0}})
  as[31:0] = as[31:0] + ad[31:0]
\end{lstlisting}

\clearpage
\subsection{EE.VLD.H.64.IP}\label{instruction:EEVLDH64IP}
\setlength\tabcolsep{6pt}
\textbf{Instruction Word} \\
\begin{tabular}{|c|c|c|c|c|c|c|c|}
	\hline
	1 & imm8[7] & qu[2:1] & 1000 & qu[0] & imm8[6:0] & as[3:0] & 0100 \\
	\hline
\end{tabular}

\textbf{Assembler Syntax} \\
EE.VLD.H.64.IP qu, as, -1024..1016 \\
\textbf{Description} \\
This instruction forces the lower 3 bits of the access address in register \lstinline{as} to 0 and loads 64-bit data from memory to the upper 64-bit segment in register \lstinline{qu}. After the access is completed, the value in register \lstinline{as} is incremented by 8-bit sign-extended constant in the instruction code segment left-shifted by 3.\\
\textbf{Operation}
\begin{lstlisting}
  qu[127: 64] = load64({as[31:  3],3{0}})
  as[31:0] = as[31:0] + {21{imm8[7]},imm8[7:0],3{0}}
\end{lstlisting}

\clearpage
\subsection{EE.VLD.H.64.XP}\label{instruction:EEVLDH64XP}
\setlength\tabcolsep{6pt}
\textbf{Instruction Word} \\
\begin{tabular}{|c|c|c|c|c|c|c|c|}
	\hline
	10 & qu[2:1] & 1101 & qu[0] & 110 & ad[3:0] & as[3:0] & 0100 \\
	\hline
\end{tabular}

\textbf{Assembler Syntax} \\
EE.VLD.H.64.XP qu, as, ad \\
\textbf{Description} \\
This instruction forces the lower 3 bits of the access address in register \lstinline{as} to 0 and loads 64-bit data from memory to the upper 64-bit segment in register \lstinline{qu}. After the access is completed, the value in register \lstinline{as} is incremented by the value in register \lstinline{ad}.\\
\textbf{Operation}
\begin{lstlisting}
  qu[127: 64] = load64({as[31:  3],3{0}})
  as[31:0] = as[31:0] + ad[31:0]
\end{lstlisting}

\clearpage
\subsection{EE.VLD.L.64.IP}\label{instruction:EEVLDL64IP}
\setlength\tabcolsep{6pt}
\textbf{Instruction Word} \\
\begin{tabular}{|c|c|c|c|c|c|c|c|}
	\hline
	1 & imm8[7] & qu[2:1] & 1001 & qu[0] & imm8[6:0] & as[3:0] & 0100 \\
	\hline
\end{tabular}

\textbf{Assembler Syntax} \\
EE.VLD.L.64.IP qu, as, -1024..1016 \\
\textbf{Description} \\
This instruction forces the lower 3 bits of the access address in register \lstinline{as} to 0 and loads 64-bit data from memory to the lower 64-bit segment in register \lstinline{qu}. After the access is completed, the value in register \lstinline{as} is incremented by 8-bit sign-extended constant in the instruction code segment left-shifted by 3.\\
\textbf{Operation}
\begin{lstlisting}
  qu[ 63:  0] = load64({as[31:3],3{0}})
  as[31:0] = as[31:0] + {21{imm8[7]},imm8[7:0],3{0}}
\end{lstlisting}

\clearpage
\subsection{EE.VLD.L.64.XP}\label{instruction:EEVLDL64XP}
\setlength\tabcolsep{6pt}
\textbf{Instruction Word} \\
\begin{tabular}{|c|c|c|c|c|c|c|c|}
	\hline
	10 & qu[2:1] & 1101 & qu[0] & 011 & ad[3:0] & as[3:0] & 0100 \\
	\hline
\end{tabular}

\textbf{Assembler Syntax} \\
EE.VLD.L.64.XP qu, as, ad \\
\textbf{Description} \\
This instruction forces the lower 3 bits of the access address in register \lstinline{as} to 0 and loads 64-bit data from memory to the lower 64-bit segment in register \lstinline{qu}. After the access is completed, the value in register \lstinline{as} is incremented by the value in register \lstinline{ad}.\\
\textbf{Operation}
\begin{lstlisting}
  qu[ 63:  0] = load64({as[31:3],3{0}})
  as[31:0] = as[31:0] + ad[31:0]
\end{lstlisting}

\clearpage
\subsection{EE.VLDBC.16}\label{instruction:EEVLDBC16}
\setlength\tabcolsep{6pt}
\textbf{Instruction Word} \\
\begin{tabular}{|c|c|c|c|c|c|c|}
	\hline
	11 & qu[2:1] & 1101 & qu[0] & 1110011 & as[3:0] & 0100 \\
	\hline
\end{tabular}

\textbf{Assembler Syntax} \\
EE.VLDBC.16 qu, as \\
\textbf{Description} \\
This instruction forces the lower 1 bit of the access address in register \lstinline{as} to 0, loads 16-bit data from memory, and broadcasts it to the eight 16-bit data segments in register \lstinline{qu}.\\
\textbf{Operation}
\begin{lstlisting}
  qu[127:0] = {8{load16({as[31:1],1{0}})}}
\end{lstlisting}

\clearpage
\subsection{EE.VLDBC.16.IP}\label{instruction:EEVLDBC16IP}
\setlength\tabcolsep{6pt}
\textbf{Instruction Word} \\
\begin{tabular}{|c|c|c|c|c|c|c|}
	\hline
	10 & qu[2:1] & 0101 & qu[0] & imm2[6:0] & as[3:0] & 0100 \\
	\hline
\end{tabular}

\textbf{Assembler Syntax} \\
EE.VLDBC.16.IP qu, as, 0..254 \\
\textbf{Description} \\
This instruction forces the lower 1 bit of the access address in register \lstinline{as} to 0, loads 16-bit data from memory, and broadcasts it to the eight 16-bit data segments in register \lstinline{qu}. After the access is completed, the value in register \lstinline{as} is incremented by 7-bit unsign-extended constant in the instruction code segment left-shifted by 1.\\
\textbf{Operation}
\begin{lstlisting}
  qu[127:0] = {8{load16({as[31:1],1{0}})}}
  as[31:0] = as[31:0] + {24{0},imm2[6:0],0}
\end{lstlisting}

\clearpage
\subsection{EE.VLDBC.16.XP}\label{instruction:EEVLDBC16XP}
\setlength\tabcolsep{6pt}
\textbf{Instruction Word} \\
\begin{tabular}{|c|c|c|c|c|c|c|c|}
	\hline
	10 & qu[2:1] & 1101 & qu[0] & 100 & ad[3:0] & as[3:0] & 0100 \\
	\hline
\end{tabular}

\textbf{Assembler Syntax} \\
EE.VLDBC.16.XP qu, as, ad \\
\textbf{Description} \\
This instruction forces the lower 1 bit of the access address in register \lstinline{as} to 0, loads 16-bit data from memory, and broadcasts it to the eight 16-bit data segments in register \lstinline{qu}. After the access is completed, the value in register \lstinline{as} is incremented by the value in register \lstinline{ad}.\\
\textbf{Operation}
\begin{lstlisting}
  qu[127:0] = {8{load16({as[31:1],1{0}})}}
  as  = as + ad[31:0]
\end{lstlisting}

\clearpage
\subsection{EE.VLDBC.32}\label{instruction:EEVLDBC32}
\setlength\tabcolsep{6pt}
\textbf{Instruction Word} \\
\begin{tabular}{|c|c|c|c|c|c|c|}
	\hline
	11 & qu[2:1] & 1101 & qu[0] & 1110111 & as[3:0] & 0100 \\
	\hline
\end{tabular}

\textbf{Assembler Syntax} \\
EE.VLDBC.32 qu, as \\
\textbf{Description} \\
This instruction forces the lower 2 bits of the access address in register \lstinline{as} to 0 and loads 32-bit data from memory and broadcasts it to the four 32-bit data segments in register \lstinline{qu}.\\
\textbf{Operation}
\begin{lstlisting}
  qu[127:0] = {4{load32({as[31:2],2{0}})}}
\end{lstlisting}

\clearpage
\subsection{EE.VLDBC.32.IP}\label{instruction:EEVLDBC32IP}
\setlength\tabcolsep{6pt}
\textbf{Instruction Word} \\
\begin{tabular}{|c|c|c|c|c|c|c|c|}
	\hline
	1 & imm4[7] & qu[2:1] & 0010 & qu[0] & imm4[6:0] & as[3:0] & 0100 \\
	\hline
\end{tabular}

\textbf{Assembler Syntax} \\
EE.VLDBC.32.IP qu, as, -256..252 \\
\textbf{Description} \\
This instruction forces the lower 2 bits of the access address in register \lstinline{as} to 0 and loads 32-bit data from memory and broadcasts it to the four 32-bit data segments in register \lstinline{qu}. After the access is completed, the value in register \lstinline{as} is incremented by 8-bit sign-extended constant in the instruction code segment left-shifted by 2.\\
\textbf{Operation}
\begin{lstlisting}
  qu[127:0] = {4{load32({as[31:2],2{0}})}}
  as[31:0] = as[31:0] + {22{imm4[7]},imm4[7:0],2{0}}
\end{lstlisting}

\clearpage
\subsection{EE.VLDBC.32.XP}\label{instruction:EEVLDBC32XP}
\setlength\tabcolsep{6pt}
\textbf{Instruction Word} \\
\begin{tabular}{|c|c|c|c|c|c|c|c|}
	\hline
	10 & qu[2:1] & 1101 & qu[0] & 001 & ad[3:0] & as[3:0] & 0100 \\
	\hline
\end{tabular}

\textbf{Assembler Syntax} \\
EE.VLDBC.32.XP qu, as, ad \\
\textbf{Description} \\
This instruction forces the lower 2 bits of the access address in register \lstinline{as} to 0, loads 32-bit data from memory and broadcasts it to the four 32-bit data segments in register \lstinline{qu}. After the access is completed, the value in register \lstinline{as} is incremented by the value in register \lstinline{ad}.\\
\textbf{Operation}
\begin{lstlisting}
  qu[127:0] = {4{load32({as[31:2],2{0}})}}
  as  = as + ad[31:0]
\end{lstlisting}

\clearpage
\subsection{EE.VLDBC.8}\label{instruction:EEVLDBC8}
\setlength\tabcolsep{6pt}
\textbf{Instruction Word} \\
\begin{tabular}{|c|c|c|c|c|c|c|}
	\hline
	11 & qu[2:1] & 1101 & qu[0] & 0111011 & as[3:0] & 0100 \\
	\hline
\end{tabular}

\textbf{Assembler Syntax} \\
EE.VLDBC.8 qu, as \\
\textbf{Description} \\
This instruction loads 8-bit data from memory at the address given by the access register \lstinline{as} and broadcasts it to the 16 8-bit data segments in register \lstinline{qu}.\\
\textbf{Operation}
\begin{lstlisting}
  qu[127:0] = {16{load8(as[31:0])}}
\end{lstlisting}

\clearpage
\subsection{EE.VLDBC.8.IP}\label{instruction:EEVLDBC8IP}
\setlength\tabcolsep{6pt}
\textbf{Instruction Word} \\
\begin{tabular}{|c|c|c|c|c|c|c|}
	\hline
	11 & qu[2:1] & 0101 & qu[0] & imm1[6:0] & as[3:0] & 0100 \\
	\hline
\end{tabular}

\textbf{Assembler Syntax} \\
EE.VLDBC.8.IP qu, as, 0..127 \\
\textbf{Description} \\
This instruction loads 8-bit data from memory at the address given by the access register \lstinline{as} and broadcasts it to the 16 8-bit data segments in register \lstinline{qu}. After the access is completed, the value in register \lstinline{as} is incremented by 7-bit unsign-extended constant in the instruction code segment.\\
\textbf{Operation}
\begin{lstlisting}
  qu[127:0] = {16{load8(as[31:0])}}
  as[31:0] = as[31:0] + {25{0},imm1[6:0]}
\end{lstlisting}

\clearpage
\subsection{EE.VLDBC.8.XP}\label{instruction:EEVLDBC8XP}
\setlength\tabcolsep{6pt}
\textbf{Instruction Word} \\
\begin{tabular}{|c|c|c|c|c|c|c|c|}
	\hline
	10 & qu[2:1] & 1101 & qu[0] & 101 & ad[3:0] & as[3:0] & 0100 \\
	\hline
\end{tabular}

\textbf{Assembler Syntax} \\
EE.VLDBC.8.XP qu, as, ad \\
\textbf{Description} \\
This instruction loads 8-bit data from memory at the address given by the access register \lstinline{as} and broadcasts it to the 16 8-bit data segments in register \lstinline{qu}. After the access is completed, the value in register \lstinline{as} is incremented by the value in register \lstinline{ad}.\\
\textbf{Operation}
\begin{lstlisting}
  qu[127:0] = {16{load8(as[31:0])}}
  as  = as + ad[31:0]
\end{lstlisting}

\clearpage
\subsection{EE.VLDHBC.16.INCP}\label{instruction:EEVLDHBC16INCP}
\setlength\tabcolsep{6pt}
\textbf{Instruction Word} \\
\begin{tabular}{|c|c|c|c|c|c|c|c|}
	\hline
	11 & qu[2:1] & 1100 & qu[0] & qu1[2:0] & 0010 & as[3:0] & 0100 \\
	\hline
\end{tabular}

\textbf{Assembler Syntax} \\
EE.VLDHBC.16.INCP qu, qu1, as \\
\textbf{Description} \\
This instruction forces the lower 4 bits of the access address in register \lstinline{as} to 0, loads 16-byte data from memory, extends it to 256 bits according to the following way, and assigns the result to registers \lstinline{qu} and \lstinline{qu1}. After the access, the value in register \lstinline{as} is incremented by 16.\\
\textbf{Operation}
\begin{lstlisting}
  dataIn[127:0] = load128({as[31:4],4{0}})
  qu  = {2{dataIn[ 63: 48]}, 2{dataIn[ 47: 32]}, 2{dataIn[ 31: 16]}, 2{dataIn[ 15:  0]} }
  qu1 = {2{dataIn[127:112]}, 2{dataIn[111: 96]}, 2{dataIn[ 95: 80]}, 2{dataIn[ 79: 64]} }
  as[31:0] = as[31:0] + 16
\end{lstlisting}

\clearpage
\subsection{EE.VMAX.S16}\label{instruction:EEVMAXS16}
\setlength\tabcolsep{6pt}
\textbf{Instruction Word} \\
\begin{tabular}{|c|c|c|c|c|c|c|c|c|}
	\hline
	10 & qa[2:1] & 1110 & qa[0] & qy[2] & 1 & qy[1:0] & qx[2:0] & 00100100 \\
	\hline
\end{tabular}

\textbf{Assembler Syntax} \\
EE.VMAX.S16 qa, qx, qy \\
\textbf{Description} \\
This instruction compares numerical values of the eight 16-bit vector data segments in registers \lstinline{qx} and \lstinline{qy}. The data segment with the larger value is written into the corresponding 16-bit data segment in register \lstinline{qa}.\\
\textbf{Operation}
\begin{lstlisting}[escapeinside=`?]
  qa[ 15:  0] = (qx[ 15:  0]>=qy[ 15:  0]) ? qx[ 15:  0] : qy[ 15:  0]
  qa[ 31: 16] = (qx[ 31: 16]>=qy[ 31: 16]) ? qx[ 31: 16] : qy[ 31: 16]
  ...
  qa[127:112] = (qx[127:112]>=qy[127:112]) ? qx[127:112] : qy[127:112]
\end{lstlisting}

\clearpage
\subsection{EE.VMAX.S16.LD.INCP}\label{instruction:EEVMAXS16LDINCP}
\setlength\tabcolsep{6pt}
\textbf{Instruction Word} \\
\begin{tabular}{|c|c|c|c|c|c|c|c|c|c|c|c|}
	\hline
	111000 & qu[2:1] & qy[0] & 001 & qu[0] & qa[2:0] & qx[1:0] & qy[2:1] & 1101 & as[3:0] & 111 & qx[2] \\
	\hline
\end{tabular}

\textbf{Assembler Syntax} \\
EE.VMAX.S16.LD.INCP qu, as, qa, qx, qy \\
\textbf{Description} \\
This instruction compares numerical values of the eight 16-bit vector data segments in registers \lstinline{qx} and \lstinline{qy}. The data segment with the larger value is written into the corresponding 16-bit data segment in register \lstinline{qa}.\\
During the operation, the lower 4 bits of the access address in register \lstinline{as} are forced to be 0, and then the 16-byte data is loaded from the memory to register \lstinline{qu}. After the access, the value in register \lstinline{as} is incremented by 16.\\
\textbf{Operation}
\begin{lstlisting}[escapeinside=`?]
  qa[ 15:  0] = (qx[ 15:  0]>=qy[ 15:  0]) ? qx[ 15:  0] : qy[ 15:  0]
  qa[ 31: 16] = (qx[ 31: 16]>=qy[ 31: 16]) ? qx[ 31: 16] : qy[ 31: 16]
  ...
  qa[127:112] = (qx[127:112]>=qy[127:112]) ? qx[127:112] : qy[127:112]

  qu[127:0] = load128({as[31:4],4{0}})
  as[31:0] = as[31:0] + 16
\end{lstlisting}

\clearpage
\subsection{EE.VMAX.S16.ST.INCP}\label{instruction:EEVMAXS16STINCP}
\setlength\tabcolsep{6pt}
\textbf{Instruction Word} \\
\begin{tabular}{|c|c|c|c|c|c|c|c|c|c|c|}
	\hline
	11100100 & qy[0] & qv[2:0] & 1 & qa[2:0] & qx[1:0] & qy[2:1] & 0011 & as[3:0] & 111 & qx[2] \\
	\hline
\end{tabular}

\textbf{Assembler Syntax} \\
EE.VMAX.S16.ST.INCP qv, as, qa, qx, qy \\
\textbf{Description} \\
This instruction compares numerical values of the eight 16-bit vector data segments in registers \lstinline{qx} and \lstinline{qy}. The data segment with the larger value is written into the corresponding 16-bit data segment in register \lstinline{qa}.\\
During the operation, the instruction forces the lower 4 bits of the access address in register \lstinline{as} to 0 and stores the value in register \lstinline{qv} to memory. After the access, the value in register \lstinline{as} is incremented by 16.\\
\textbf{Operation}
\begin{lstlisting}[escapeinside=`?]
  qa[ 15:  0] = (qx[ 15:  0]>=qy[ 15:  0]) ? qx[ 15:  0] : qy[ 15:  0]
  qa[ 31: 16] = (qx[ 31: 16]>=qy[ 31: 16]) ? qx[ 31: 16] : qy[ 31: 16]
  ...
  qa[127:112] = (qx[127:112]>=qy[127:112]) ? qx[127:112] : qy[127:112]

  qv[127:0] => store128({as[31:4],4{0}})
  as[31:0] = as[31:0] + 16
\end{lstlisting}

\clearpage
\subsection{EE.VMAX.S32}\label{instruction:EEVMAXS32}
\setlength\tabcolsep{6pt}
\textbf{Instruction Word} \\
\begin{tabular}{|c|c|c|c|c|c|c|c|c|}
	\hline
	10 & qa[2:1] & 1110 & qa[0] & qy[2] & 1 & qy[1:0] & qx[2:0] & 00110100 \\
	\hline
\end{tabular}

\textbf{Assembler Syntax} \\
EE.VMAX.S32 qa, qx, qy \\
\textbf{Description} \\
This instruction compares numerical values of the four 32-bit vector data segments in registers \lstinline{qx} and \lstinline{qy}. The data segment with the larger value is written into the corresponding 32-bit data segment in register \lstinline{qa}.\\
\textbf{Operation}
\begin{lstlisting}[escapeinside=`?]
  qa[ 31:  0] = (qx[ 31:  0]>=qy[ 31:  0]) ? qx[ 31:  0] : qy[ 31:  0]
  qa[ 63: 32] = (qx[ 63: 32]>=qy[ 63: 32]) ? qx[ 63: 32] : qy[ 63: 32]
  ...
  qa[127: 96] = (qx[127: 96]>=qy[127: 96]) ? qx[127: 96] : qy[127: 96]
\end{lstlisting}

\clearpage
\subsection{EE.VMAX.S32.LD.INCP}\label{instruction:EEVMAXS32LDINCP}
\setlength\tabcolsep{6pt}
\textbf{Instruction Word} \\
\begin{tabular}{|c|c|c|c|c|c|c|c|c|c|c|c|}
	\hline
	111000 & qu[2:1] & qy[0] & 001 & qu[0] & qa[2:0] & qx[1:0] & qy[2:1] & 1110 & as[3:0] & 111 & qx[2] \\
	\hline
\end{tabular}

\textbf{Assembler Syntax} \\
EE.VMAX.S32.LD.INCP qu, as, qa, qx, qy \\
\textbf{Description} \\
This instruction compares numerical values of the four 32-bit vector data segments in registers \lstinline{qx} and \lstinline{qy}. The data segment with the larger value is written into the corresponding 32-bit data segment in register \lstinline{qa}.\\
During the operation, the lower 4 bits of the access address in register \lstinline{as} are forced to be 0, and then the 16-byte data is loaded from the memory to register \lstinline{qu}. After the access, the value in register \lstinline{as} is incremented by 16.\\
\textbf{Operation}
\begin{lstlisting}[escapeinside=`?]
  qa[ 31:  0] = (qx[ 31:  0]>=qy[ 31:  0]) ? qx[ 31:  0] : qy[ 31:  0]
  qa[ 63: 32] = (qx[ 63: 32]>=qy[ 63: 32]) ? qx[ 63: 32] : qy[ 63: 32]
  ...
  qa[127: 96] = (qx[127: 96]>=qy[127: 96]) ? qx[127: 96] : qy[127: 96]

  qu[127:0] = load128({as[31:4],4{0}})
  as[31:0] = as[31:0] + 16
\end{lstlisting}

\clearpage
\subsection{EE.VMAX.S32.ST.INCP}\label{instruction:EEVMAXS32STINCP}
\setlength\tabcolsep{6pt}
\textbf{Instruction Word} \\
\begin{tabular}{|c|c|c|c|c|c|c|c|c|c|c|}
	\hline
	11100101 & qy[0] & qv[2:0] & 0 & qa[2:0] & qx[1:0] & qy[2:1] & 0000 & as[3:0] & 111 & qx[2] \\
	\hline
\end{tabular}

\textbf{Assembler Syntax} \\
EE.VMAX.S32.ST.INCP qv, as, qa, qx, qy \\
\textbf{Description} \\
This instruction compares numerical values of the four 32-bit vector data segments in registers \lstinline{qx} and \lstinline{qy}. The data segment with the larger value is written into the corresponding 32-bit data segment in register \lstinline{qa}.\\
During the operation, the instruction forces the lower 4 bits of the access address in register \lstinline{as} to 0 and stores the value in register \lstinline{qv} to memory. After the access, the value in register \lstinline{as} is incremented by 16.\\
\textbf{Operation}
\begin{lstlisting}[escapeinside=`?]
  qa[ 31:  0] = (qx[ 31:  0]>=qy[ 31:  0]) ? qx[ 31:  0] : qy[ 31:  0]
  qa[ 63: 32] = (qx[ 63: 32]>=qy[ 63: 32]) ? qx[ 63: 32] : qy[ 63: 32]
  ...
  qa[127: 96] = (qx[127: 96]>=qy[127: 96]) ? qx[127: 96] : qy[127: 96]

  qv[127:0] => store128({as[31:4],4{0}})
  as[31:0] = as[31:0] + 16
\end{lstlisting}

\clearpage
\subsection{EE.VMAX.S8}\label{instruction:EEVMAXS8}
\setlength\tabcolsep{6pt}
\textbf{Instruction Word} \\
\begin{tabular}{|c|c|c|c|c|c|c|c|c|}
	\hline
	10 & qa[2:1] & 1110 & qa[0] & qy[2] & 1 & qy[1:0] & qx[2:0] & 01000100 \\
	\hline
\end{tabular}

\textbf{Assembler Syntax} \\
EE.VMAX.S8 qa, qx, qy \\
\textbf{Description} \\
This instruction compares numerical values of the 16 8-bit vector data segments in registers \lstinline{qx} and \lstinline{qy}. The data segment with the larger value is written into the corresponding 8-bit data segment in register \lstinline{qa}.\\
\textbf{Operation}
\begin{lstlisting}[escapeinside=`?]
  qa[  7:  0] = (qx[  7:  0]>=qy[  7:  0]) ? qx[  7:  0] : qy[  7:  0]
  qa[ 15:  8] = (qx[ 15:  8]>=qy[ 15:  8]) ? qx[ 15:  8] : qy[ 15:  8]
  ...
  qa[127:120] = (qx[127:120]>=qy[127:120]) ? qx[127:120] : qy[127:120]
\end{lstlisting}

\clearpage
\subsection{EE.VMAX.S8.LD.INCP}\label{instruction:EEVMAXS8LDINCP}
\setlength\tabcolsep{6pt}
\textbf{Instruction Word} \\
\begin{tabular}{|c|c|c|c|c|c|c|c|c|c|c|c|}
	\hline
	111000 & qu[2:1] & qy[0] & 001 & qu[0] & qa[2:0] & qx[1:0] & qy[2:1] & 1111 & as[3:0] & 111 & qx[2] \\
	\hline
\end{tabular}

\textbf{Assembler Syntax} \\
EE.VMAX.S8.LD.INCP qu, as, qa, qx, qy \\
\textbf{Description} \\
This instruction compares numerical values of the 16 8-bit vector data segments in registers \lstinline{qx} and \lstinline{qy}. The data segment with the larger value is written into the corresponding 8-bit data segment in register \lstinline{qa}.\\
During the operation, the lower 4 bits of the access address in register \lstinline{as} are forced to be 0, and then the 16-byte data is loaded from the memory to register \lstinline{qu}. After the access, the value in register \lstinline{as} is incremented by 16.\\
\textbf{Operation}
\begin{lstlisting}[escapeinside=`?]
  qa[  7:  0] = (qx[  7:  0]>=qy[  7:  0]) ? qx[  7:  0] : qy[  7:  0]
  qa[ 15:  8] = (qx[ 15:  8]>=qy[ 15:  8]) ? qx[ 15:  8] : qy[ 15:  8]
  ...
  qa[127:120] = (qx[127:120]>=qy[127:120]) ? qx[127:120] : qy[127:120]

  qu[127:0] = load128({as[31:4],4{0}})
  as[31:0] = as[31:0] + 16
\end{lstlisting}

\clearpage
\subsection{EE.VMAX.S8.ST.INCP}\label{instruction:EEVMAXS8STINCP}
\setlength\tabcolsep{6pt}
\textbf{Instruction Word} \\
\begin{tabular}{|c|c|c|c|c|c|c|c|c|c|c|}
	\hline
	11100101 & qy[0] & qv[2:0] & 1 & qa[2:0] & qx[1:0] & qy[2:1] & 0000 & as[3:0] & 111 & qx[2] \\
	\hline
\end{tabular}

\textbf{Assembler Syntax} \\
EE.VMAX.S8.ST.INCP qv, as, qa, qx, qy \\
\textbf{Description} \\
This instruction compares numerical values of the 16 8-bit vector data segments in registers \lstinline{qx} and \lstinline{qy}. The data segment with the larger value is written into the corresponding 8-bit data segment in register \lstinline{qa}.\\
During the operation, the instruction forces the lower 4 bits of the access address in register \lstinline{as} to 0 and stores the value in register \lstinline{qv} to memory. After the access, the value in register \lstinline{as} is incremented by 16.\\
\textbf{Operation}
\begin{lstlisting}[escapeinside=`?]
  qa[  7:  0] = (qx[  7:  0]>=qy[  7:  0]) ? qx[  7:  0] : qy[  7:  0]
  qa[ 15:  8] = (qx[ 15:  8]>=qy[ 15:  8]) ? qx[ 15:  8] : qy[ 15:  8]
  ...
  qa[127:120] = (qx[127:120]>=qy[127:120]) ? qx[127:120] : qy[127:120]

  qv[127:0] => store128({as[31:4],4{0}})
  as[31:0] = as[31:0] + 16
\end{lstlisting}

\clearpage
\subsection{EE.VMIN.S16}\label{instruction:EEVMINS16}
\setlength\tabcolsep{6pt}
\textbf{Instruction Word} \\
\begin{tabular}{|c|c|c|c|c|c|c|c|c|}
	\hline
	10 & qa[2:1] & 1110 & qa[0] & qy[2] & 1 & qy[1:0] & qx[2:0] & 01010100 \\
	\hline
\end{tabular}

\textbf{Assembler Syntax} \\
EE.VMIN.S16 qa, qx, qy \\
\textbf{Description} \\
This instruction compares numerical values of the eight 16-bit vector data segments in registers \lstinline{qx} and \lstinline{qy}. The data segment with the smaller value is written into the corresponding 16-bit data segment in register \lstinline{qa}.\\
\textbf{Operation}
\begin{lstlisting}[escapeinside=`?]
  qa[ 15:  0] = (qx[ 15:  0]<=qy[ 15:  0]) ? qx[ 15:  0] : qy[ 15:  0]
  qa[ 31: 16] = (qx[ 31: 16]<=qy[ 31: 16]) ? qx[ 31: 16] : qy[ 31: 16]
  ...
  qa[127:112] = (qx[127:112]<=qy[127:112]) ? qx[127:112] : qy[127:112]
\end{lstlisting}

\clearpage
\subsection{EE.VMIN.S16.LD.INCP}\label{instruction:EEVMINS16LDINCP}
\setlength\tabcolsep{6pt}
\textbf{Instruction Word} \\
\begin{tabular}{|c|c|c|c|c|c|c|c|c|c|c|c|}
	\hline
	111000 & qu[2:1] & qy[0] & 010 & qu[0] & qa[2:0] & qx[1:0] & qy[2:1] & 1110 & as[3:0] & 111 & qx[2] \\
	\hline
\end{tabular}

\textbf{Assembler Syntax} \\
EE.VMIN.S16.LD.INCP qu, as, qa, qx, qy \\
\textbf{Description} \\
This instruction compares numerical values of the eight 16-bit vector data segments in registers \lstinline{qx} and \lstinline{qy}. The data segment with the smaller value is written into the corresponding 16-bit data segment in register \lstinline{qa}.\\
During the operation, the lower 4 bits of the access address in register \lstinline{as} are forced to be 0, and then the 16-byte data is loaded from the memory to register \lstinline{qu}. After the access, the value in register \lstinline{as} is incremented by 16.\\
\textbf{Operation}
\begin{lstlisting}[escapeinside=`?]
  qa[ 15:  0] = (qx[ 15:  0]<=qy[ 15:  0]) ? qx[ 15:  0] : qy[ 15:  0]
  qa[ 31: 16] = (qx[ 31: 16]<=qy[ 31: 16]) ? qx[ 31: 16] : qy[ 31: 16]
  ...
  qa[127:112] = (qx[127:112]<=qy[127:112]) ? qx[127:112] : qy[127:112]

  qu[127:0] = load128({as[31:4],4{0}})
  as[31:0] = as[31:0] + 16
\end{lstlisting}

\clearpage
\subsection{EE.VMIN.S16.ST.INCP}\label{instruction:EEVMINS16STINCP}
\setlength\tabcolsep{6pt}
\textbf{Instruction Word} \\
\begin{tabular}{|c|c|c|c|c|c|c|c|c|c|c|}
	\hline
	11100101 & qy[0] & qv[2:0] & 0 & qa[2:0] & qx[1:0] & qy[2:1] & 0001 & as[3:0] & 111 & qx[2] \\
	\hline
\end{tabular}

\textbf{Assembler Syntax} \\
EE.VMIN.S16.ST.INCP qv, as, qa, qx, qy \\
\textbf{Description} \\
This instruction compares numerical values of the eight 16-bit vector data segments in registers \lstinline{qx} and \lstinline{qy}. The data segment with the smaller value is written into the corresponding 16-bit data segment in register \lstinline{qa}.\\
During the operation, the instruction forces the lower 4 bits of the access address in register \lstinline{as} to 0 and stores the value in register \lstinline{qv} to memory. After the access, the value in register \lstinline{as} is incremented by 16.\\
\textbf{Operation}
\begin{lstlisting}[escapeinside=`?]
  qa[ 15:  0] = (qx[ 15:  0]<=qy[ 15:  0]) ? qx[ 15:  0] : qy[ 15:  0]
  qa[ 31: 16] = (qx[ 31: 16]<=qy[ 31: 16]) ? qx[ 31: 16] : qy[ 31: 16]
  ...
  qa[127:112] = (qx[127:112]<=qy[127:112]) ? qx[127:112] : qy[127:112]

  qv[127:0] => store128({as[31:4],4{0}})
  as[31:0] = as[31:0] + 16
\end{lstlisting}

\clearpage
\subsection{EE.VMIN.S32}\label{instruction:EEVMINS32}
\setlength\tabcolsep{6pt}
\textbf{Instruction Word} \\
\begin{tabular}{|c|c|c|c|c|c|c|c|c|}
	\hline
	10 & qa[2:1] & 1110 & qa[0] & qy[2] & 1 & qy[1:0] & qx[2:0] & 01100100 \\
	\hline
\end{tabular}

\textbf{Assembler Syntax} \\
EE.VMIN.S32 qa, qx, qy \\
\textbf{Description} \\
This instruction compares numerical values of the four 32-bit vector data segments in registers \lstinline{qx} and \lstinline{qy}. The data segment with the smaller value is written into the corresponding 32-bit data segment in register \lstinline{qa}.\\
\textbf{Operation}
\begin{lstlisting}[escapeinside=`?]
  qa[ 31:  0] = (qx[ 31:  0]<=qy[ 31:  0]) ? qx[ 31:  0] : qy[ 31:  0]
  qa[ 63: 32] = (qx[ 63: 32]<=qy[ 63: 32]) ? qx[ 63: 32] : qy[ 63: 32]
  ...
  qa[127: 96] = (qx[127: 96]<=qy[127: 96]) ? qx[127: 96] : qy[127: 96]
\end{lstlisting}

\clearpage
\subsection{EE.VMIN.S32.LD.INCP}\label{instruction:EEVMINS32LDINCP}
\setlength\tabcolsep{6pt}
\textbf{Instruction Word} \\
\begin{tabular}{|c|c|c|c|c|c|c|c|c|c|c|c|}
	\hline
	111000 & qu[2:1] & qy[0] & 011 & qu[0] & qa[2:0] & qx[1:0] & qy[2:1] & 1110 & as[3:0] & 111 & qx[2] \\
	\hline
\end{tabular}

\textbf{Assembler Syntax} \\
EE.VMIN.S32.LD.INCP qu, as, qa, qx, qy \\
\textbf{Description} \\
This instruction compares numerical values of the four 32-bit vector data segments in registers \lstinline{qx} and \lstinline{qy}. The data segment with the smaller value is written into the corresponding 32-bit data segment in register \lstinline{qa}.\\
During the operation, the lower 4 bits of the access address in register \lstinline{as} are forced to be 0, and then the 16-byte data is loaded from the memory to register \lstinline{qu}. After the access, the value in register \lstinline{as} is incremented by 16.\\
\textbf{Operation}
\begin{lstlisting}[escapeinside=`?]
  qa[ 31:  0] = (qx[ 31:  0]<=qy[ 31:  0]) ? qx[ 31:  0] : qy[ 31:  0]
  qa[ 63: 32] = (qx[ 63: 32]<=qy[ 63: 32]) ? qx[ 63: 32] : qy[ 63: 32]
  ...
  qa[127: 96] = (qx[127: 96]<=qy[127: 96]) ? qx[127: 96] : qy[127: 96]

  qu[127:0] = load128({as[31:4],4{0}})
  as[31:0] = as[31:0] + 16
\end{lstlisting}

\clearpage
\subsection{EE.VMIN.S32.ST.INCP}\label{instruction:EEVMINS32STINCP}
\setlength\tabcolsep{6pt}
\textbf{Instruction Word} \\
\begin{tabular}{|c|c|c|c|c|c|c|c|c|c|c|}
	\hline
	11100101 & qy[0] & qv[2:0] & 1 & qa[2:0] & qx[1:0] & qy[2:1] & 0001 & as[3:0] & 111 & qx[2] \\
	\hline
\end{tabular}

\textbf{Assembler Syntax} \\
EE.VMIN.S32.ST.INCP qv, as, qa, qx, qy \\
\textbf{Description} \\
This instruction compares numerical values of the four 32-bit vector data segments in registers \lstinline{qx} and \lstinline{qy}. The data segment with the smaller value is written into the corresponding 32-bit data segment in register \lstinline{qa}.\\
During the operation, the instruction forces the lower 4 bits of the access address in register \lstinline{as} to 0 and stores the value in register \lstinline{qv} to memory. After the access, the value in register \lstinline{as} is incremented by 16.\\
\textbf{Operation}
\begin{lstlisting}[escapeinside=`?]
  qa[ 31:  0] = (qx[ 31:  0]<=qy[ 31:  0]) ? qx[ 31:  0] : qy[ 31:  0]
  qa[ 63: 32] = (qx[ 63: 32]<=qy[ 63: 32]) ? qx[ 63: 32] : qy[ 63: 32]
  ...
  qa[127: 96] = (qx[127: 96]<=qy[127: 96]) ? qx[127: 96] : qy[127: 96]

  qv[127:0] => store128({as[31:4],4{0}})
  as[31:0] = as[31:0] + 16
\end{lstlisting}

\clearpage
\subsection{EE.VMIN.S8}\label{instruction:EEVMINS8}
\setlength\tabcolsep{6pt}
\textbf{Instruction Word} \\
\begin{tabular}{|c|c|c|c|c|c|c|c|c|}
	\hline
	10 & qa[2:1] & 1110 & qa[0] & qy[2] & 1 & qy[1:0] & qx[2:0] & 01110100 \\
	\hline
\end{tabular}

\textbf{Assembler Syntax} \\
EE.VMIN.S8 qa, qx, qy \\
\textbf{Description} \\
This instruction compares numerical values of the 16 8-bit vector data segments in registers \lstinline{qx} and \lstinline{qy}. The data segment with the smaller value is written into the corresponding 8-bit data segment in register \lstinline{qa}.\\
\textbf{Operation}
\begin{lstlisting}[escapeinside=`?]
  qa[  7:  0] = (qx[  7:  0]<=qy[  7:  0]) ? qx[  7:  0] : qy[  7:  0]
  qa[ 15:  8] = (qx[ 15:  8]<=qy[ 15:  8]) ? qx[ 15:  8] : qy[ 15:  8]
  ...
  qa[127:120] = (qx[127:120]<=qy[127:120]) ? qx[127:120] : qy[127:120]
\end{lstlisting}

\clearpage
\subsection{EE.VMIN.S8.LD.INCP}\label{instruction:EEVMINS8LDINCP}
\setlength\tabcolsep{6pt}
\textbf{Instruction Word} \\
\begin{tabular}{|c|c|c|c|c|c|c|c|c|c|c|c|}
	\hline
	111000 & qu[2:1] & qy[0] & 010 & qu[0] & qa[2:0] & qx[1:0] & qy[2:1] & 1111 & as[3:0] & 111 & qx[2] \\
	\hline
\end{tabular}

\textbf{Assembler Syntax} \\
EE.VMIN.S8.LD.INCP qu, as, qa, qx, qy \\
\textbf{Description} \\
This instruction compares numerical values of the 16 8-bit vector data segments in registers \lstinline{qx} and \lstinline{qy}. The data segment with the smaller value is written into the corresponding 8-bit data segment in register \lstinline{qa}.\\
During the operation, the lower 4 bits of the access address in register \lstinline{as} are forced to be 0, and then the 16-byte data is loaded from the memory to register \lstinline{qu}. After the access, the value in register \lstinline{as} is incremented by 16.\\
\textbf{Operation}
\begin{lstlisting}[escapeinside=`?]
  qa[  7:  0] = (qx[  7:  0]<=qy[  7:  0]) ? qx[  7:  0] : qy[  7:  0]
  qa[ 15:  8] = (qx[ 15:  8]<=qy[ 15:  8]) ? qx[ 15:  8] : qy[ 15:  8]
  ...
  qa[127:120] = (qx[127:120]<=qy[127:120]) ? qx[127:120] : qy[127:120]

  qu[127:0] = load128({as[31:4],4{0}})
  as[31:0] = as[31:0] + 16
\end{lstlisting}

\clearpage
\subsection{EE.VMIN.S8.ST.INCP}\label{instruction:EEVMINS8STINCP}
\setlength\tabcolsep{6pt}
\textbf{Instruction Word} \\
\begin{tabular}{|c|c|c|c|c|c|c|c|c|c|c|}
	\hline
	11100101 & qy[0] & qv[2:0] & 0 & qa[2:0] & qx[1:0] & qy[2:1] & 0010 & as[3:0] & 111 & qx[2] \\
	\hline
\end{tabular}

\textbf{Assembler Syntax} \\
EE.VMIN.S8.ST.INCP qv, as, qa, qx, qy \\
\textbf{Description} \\
This instruction compares numerical values of the 16 8-bit vector data segments in registers \lstinline{qx} and \lstinline{qy}. The data segment with the smaller value is written into the corresponding 8-bit data segment in register \lstinline{qa}.\\
During the operation, the instruction forces the lower 4 bits of the access address in register \lstinline{as} to 0 and stores the value in register \lstinline{qv} to memory. After the access, the value in register \lstinline{as} is incremented by 16.\\
\textbf{Operation}
\begin{lstlisting}[escapeinside=`?]
  qa[  7:  0] = (qx[  7:  0]<=qy[  7:  0]) ? qx[  7:  0] : qy[  7:  0]
  qa[ 15:  8] = (qx[ 15:  8]<=qy[ 15:  8]) ? qx[ 15:  8] : qy[ 15:  8]
  ...
  qa[127:120] = (qx[127:120]<=qy[127:120]) ? qx[127:120] : qy[127:120]

  qv[127:0] => store128({as[31:4],4{0}})
  as[31:0] = as[31:0] + 16
\end{lstlisting}

\clearpage
\subsection{EE.VMUL.S16}\label{instruction:EEVMULS16}
\setlength\tabcolsep{6pt}
\textbf{Instruction Word} \\
\begin{tabular}{|c|c|c|c|c|c|c|c|c|}
	\hline
	10 & qz[2:1] & 1110 & qz[0] & qy[2] & 1 & qy[1:0] & qx[2:0] & 10000100 \\
	\hline
\end{tabular}

\textbf{Assembler Syntax} \\
EE.VMUL.S16 qz, qx, qy \\
\textbf{Description} \\
This instruction performs a signed vector multiplication on 16-bit data. Registers \lstinline{qx} and \lstinline{qy} are the multiplier and the multiplicand respectively. The eight 32-bit data results obtained from the calculation is arithmetically right-shifted by the value in special register \lstinline{SAR}. Then, the lower 16-bit data of the shift result is written into corresponding segment of register \lstinline{qz}.\\
\textbf{Operation}
\begin{lstlisting}
  qz[ 15:  0] = (qx[ 15:  0] * qy[ 15:  0]) >> SAR[5:0]
  qz[ 31: 16] = (qx[ 31: 16] * qy[ 31: 16]) >> SAR[5:0]
  qz[ 47: 32] = (qx[ 47: 32] * qy[ 47: 32]) >> SAR[5:0]
  qz[ 63: 48] = (qx[ 63: 48] * qy[ 63: 48]) >> SAR[5:0]
  qz[ 79: 64] = (qx[ 79: 64] * qy[ 79: 64]) >> SAR[5:0]
  qz[ 95: 80] = (qx[ 95: 80] * qy[ 95: 80]) >> SAR[5:0]
  qz[111: 96] = (qx[111: 96] * qy[111: 96]) >> SAR[5:0]
  qz[127:112] = (qx[127:112] * qy[127:112]) >> SAR[5:0]
\end{lstlisting}

\clearpage
\subsection{EE.VMUL.S16.LD.INCP}\label{instruction:EEVMULS16LDINCP}
\setlength\tabcolsep{6pt}
\textbf{Instruction Word} \\
\begin{tabular}{|c|c|c|c|c|c|c|c|c|c|c|c|}
	\hline
	111000 & qu[2:1] & qy[0] & 011 & qu[0] & qz[2:0] & qx[1:0] & qy[2:1] & 1111 & as[3:0] & 111 & qx[2] \\
	\hline
\end{tabular}

\textbf{Assembler Syntax} \\
EE.VMUL.S16.LD.INCP qu, as, qz, qx, qy \\
\textbf{Description} \\
This instruction performs a signed vector multiplication on 16-bit data. Registers \lstinline{qx} and \lstinline{qy} are the multiplier and the multiplicand respectively. The eight 32-bit data results obtained from the calculation is arithmetically right-shifted by the value in special register \lstinline{SAR}. Then, the lower 16-bit data of the shift result is written into corresponding segment of register \lstinline{qz}. \\ During the operation, the lower 4 bits of the access address in register \lstinline{as} are forced to be 0, and then the 16-byte data is loaded from the memory to register \lstinline{qu}. After the access, the value in register \lstinline{as} is incremented by 16.\\
\textbf{Operation}
\begin{lstlisting}
  qz[ 15:  0] = (qx[ 15:  0] * qy[ 15:  0]) >> SAR[5:0]
  qz[ 31: 16] = (qx[ 31: 16] * qy[ 31: 16]) >> SAR[5:0]
  qz[ 47: 32] = (qx[ 47: 32] * qy[ 47: 32]) >> SAR[5:0]
  qz[ 63: 48] = (qx[ 63: 48] * qy[ 63: 48]) >> SAR[5:0]
  qz[ 79: 64] = (qx[ 79: 64] * qy[ 79: 64]) >> SAR[5:0]
  qz[ 95: 80] = (qx[ 95: 80] * qy[ 95: 80]) >> SAR[5:0]
  qz[111: 96] = (qx[111: 96] * qy[111: 96]) >> SAR[5:0]
  qz[127:112] = (qx[127:112] * qy[127:112]) >> SAR[5:0]

  qu[127:0] = load128({as[31:4],4{0}})
  as[31:0] = as[31:0] + 16
\end{lstlisting}

\clearpage
\subsection{EE.VMUL.S16.ST.INCP}\label{instruction:EEVMULS16STINCP}
\setlength\tabcolsep{6pt}
\textbf{Instruction Word} \\
\begin{tabular}{|c|c|c|c|c|c|c|c|c|c|c|}
	\hline
	11100101 & qy[0] & qv[2:0] & 1 & qz[2:0] & qx[1:0] & qy[2:1] & 0010 & as[3:0] & 111 & qx[2] \\
	\hline
\end{tabular}

\textbf{Assembler Syntax} \\
EE.VMUL.S16.ST.INCP qv, as, qz, qx, qy \\
\textbf{Description} \\
This instruction performs a signed vector multiplication on 16-bit data. Registers \lstinline{qx} and \lstinline{qy} are the multiplier and the multiplicand respectively. The eight 32-bit data results obtained from the calculation is arithmetically right-shifted by the value in special register \lstinline{SAR}. Then, the lower 16-bit data of the shift result is written into corresponding segment of register \lstinline{qz}. \\ During the operation, the instruction forces the lower 4 bits of the access address in register \lstinline{as} to 0 and stores the value in register \lstinline{qv} to memory. After the access, the value in register \lstinline{as} is incremented by 16.\\
\textbf{Operation}
\begin{lstlisting}
  qz[ 15:  0] = (qx[ 15:  0] * qy[ 15:  0]) >> SAR[5:0]
  qz[ 31: 16] = (qx[ 31: 16] * qy[ 31: 16]) >> SAR[5:0]
  qz[ 47: 32] = (qx[ 47: 32] * qy[ 47: 32]) >> SAR[5:0]
  qz[ 63: 48] = (qx[ 63: 48] * qy[ 63: 48]) >> SAR[5:0]
  qz[ 79: 64] = (qx[ 79: 64] * qy[ 79: 64]) >> SAR[5:0]
  qz[ 95: 80] = (qx[ 95: 80] * qy[ 95: 80]) >> SAR[5:0]
  qz[111: 96] = (qx[111: 96] * qy[111: 96]) >> SAR[5:0]
  qz[127:112] = (qx[127:112] * qy[127:112]) >> SAR[5:0]

  qv[127:0] => store128({as[31:4],4{0}})
  as[31:0] = as[31:0] + 16
\end{lstlisting}

\clearpage
\subsection{EE.VMUL.S8}\label{instruction:EEVMULS8}
\setlength\tabcolsep{6pt}
\textbf{Instruction Word} \\
\begin{tabular}{|c|c|c|c|c|c|c|c|c|}
	\hline
	10 & qz[2:1] & 1110 & qz[0] & qy[2] & 1 & qy[1:0] & qx[2:0] & 10010100 \\
	\hline
\end{tabular}

\textbf{Assembler Syntax} \\
EE.VMUL.S8 qz, qx, qy \\
\textbf{Description} \\
This instruction performs a signed vector multiplication on 8-bit data. Registers \lstinline{qx} and \lstinline{qy} are the multiplier and the multiplicand respectively. The 16 16-bit data results obtained from the calculation is arithmetically right-shifted by the value in special register \lstinline{SAR}. Then, the lower 8-bit data of the shift result is written into corresponding segment of register \lstinline{qz}.\\
\textbf{Operation}
\begin{lstlisting}
  qz[  7:  0] = (qx[  7:  0] * qy[  7:  0]) >> SAR[5:0]
  qz[ 15:  8] = (qx[ 15:  8] * qy[ 15:  8]) >> SAR[5:0]
  qz[ 23: 16] = (qx[ 23: 16] * qy[ 23: 16]) >> SAR[5:0]
  qz[ 31: 24] = (qx[ 31: 24] * qy[ 31: 24]) >> SAR[5:0]
  qz[ 39: 32] = (qx[ 39: 32] * qy[ 39: 32]) >> SAR[5:0]
  qz[ 47: 40] = (qx[ 47: 40] * qy[ 47: 40]) >> SAR[5:0]
  qz[ 55: 48] = (qx[ 55: 48] * qy[ 55: 48]) >> SAR[5:0]
  qz[ 63: 56] = (qx[ 63: 56] * qy[ 63: 56]) >> SAR[5:0]
  qz[ 71: 64] = (qx[ 71: 64] * qy[ 71: 64]) >> SAR[5:0]
  qz[ 79: 72] = (qx[ 79: 72] * qy[ 79: 72]) >> SAR[5:0]
  qz[ 87: 80] = (qx[ 87: 80] * qy[ 87: 80]) >> SAR[5:0]
  qz[ 95: 88] = (qx[ 95: 88] * qy[ 95: 88]) >> SAR[5:0]
  qz[103: 96] = (qx[103: 96] * qy[103: 96]) >> SAR[5:0]
  qz[111:104] = (qx[111:104] * qy[111:104]) >> SAR[5:0]
  qz[119:112] = (qx[119:112] * qy[119:112]) >> SAR[5:0]
  qz[127:120] = (qx[127:120] * qy[127:120]) >> SAR[5:0]
\end{lstlisting}

\clearpage
\subsection{EE.VMUL.S8.LD.INCP}\label{instruction:EEVMULS8LDINCP}
\setlength\tabcolsep{6pt}
\textbf{Instruction Word} \\
\begin{tabular}{|c|c|c|c|c|c|c|c|c|c|c|c|}
	\hline
	111000 & qu[2:1] & qy[0] & 100 & qu[0] & qz[2:0] & qx[1:0] & qy[2:1] & 1100 & as[3:0] & 111 & qx[2] \\
	\hline
\end{tabular}

\textbf{Assembler Syntax} \\
EE.VMUL.S8.LD.INCP qu, as, qz, qx, qy \\
\textbf{Description} \\
This instruction performs a signed vector multiplication on 8-bit data. Registers \lstinline{qx} and \lstinline{qy} are the multiplier and the multiplicand respectively. The 16 16-bit data results obtained from the calculation is arithmetically right-shifted by the value in special register \lstinline{SAR}. Then, the lower 8-bit data of the shift result is written into corresponding segment of register \lstinline{qz}. \\ During the operation, the lower 4 bits of the access address in register \lstinline{as} are forced to be 0, and then the 16-byte data is loaded from the memory to register \lstinline{qu}. After the access, the value in register \lstinline{as} is incremented by 16.\\
\textbf{Operation}
\begin{lstlisting}
  qz[  7:  0] = (qx[  7:  0] * qy[  7:  0]) >> SAR[5:0]
  qz[ 15:  8] = (qx[ 15:  8] * qy[ 15:  8]) >> SAR[5:0]
  qz[ 23: 16] = (qx[ 23: 16] * qy[ 23: 16]) >> SAR[5:0]
  qz[ 31: 24] = (qx[ 31: 24] * qy[ 31: 24]) >> SAR[5:0]
  qz[ 39: 32] = (qx[ 39: 32] * qy[ 39: 32]) >> SAR[5:0]
  qz[ 47: 40] = (qx[ 47: 40] * qy[ 47: 40]) >> SAR[5:0]
  qz[ 55: 48] = (qx[ 55: 48] * qy[ 55: 48]) >> SAR[5:0]
  qz[ 63: 56] = (qx[ 63: 56] * qy[ 63: 56]) >> SAR[5:0]
  qz[ 71: 64] = (qx[ 71: 64] * qy[ 71: 64]) >> SAR[5:0]
  qz[ 79: 72] = (qx[ 79: 72] * qy[ 79: 72]) >> SAR[5:0]
  qz[ 87: 80] = (qx[ 87: 80] * qy[ 87: 80]) >> SAR[5:0]
  qz[ 95: 88] = (qx[ 95: 88] * qy[ 95: 88]) >> SAR[5:0]
  qz[103: 96] = (qx[103: 96] * qy[103: 96]) >> SAR[5:0]
  qz[111:104] = (qx[111:104] * qy[111:104]) >> SAR[5:0]
  qz[119:112] = (qx[119:112] * qy[119:112]) >> SAR[5:0]
  qz[127:120] = (qx[127:120] * qy[127:120]) >> SAR[5:0]

  qu[127:0] = load128({as[31:4],4{0}})
  as[31:0] = as[31:0] + 16
\end{lstlisting}

\clearpage
\subsection{EE.VMUL.S8.ST.INCP}\label{instruction:EEVMULS8STINCP}
\setlength\tabcolsep{6pt}
\textbf{Instruction Word} \\
\begin{tabular}{|c|c|c|c|c|c|c|c|c|c|c|}
	\hline
	11100101 & qy[0] & qv[2:0] & 0 & qz[2:0] & qx[1:0] & qy[2:1] & 0011 & as[3:0] & 111 & qx[2] \\
	\hline
\end{tabular}

\textbf{Assembler Syntax} \\
EE.VMUL.S8.ST.INCP qv, as, qz, qx, qy \\
\textbf{Description} \\
This instruction performs a signed vector multiplication on 8-bit data. Registers \lstinline{qx} and \lstinline{qy} are the multiplier and the multiplicand respectively. The 16 16-bit data results obtained from the calculation is arithmetically right-shifted by the value in special register \lstinline{SAR}. Then, the lower 8-bit data of the shift result is written into corresponding segment of register \lstinline{qz}. \\ During the operation, the instruction forces the lower 4 bits of the access address in register \lstinline{as} to 0 and stores the value in register \lstinline{qv} to memory. After the access, the value in register \lstinline{as} is incremented by 16.\\
\textbf{Operation}
\begin{lstlisting}
  qz[  7:  0] = (qx[  7:  0] * qy[  7:  0]) >> SAR[5:0]
  qz[ 15:  8] = (qx[ 15:  8] * qy[ 15:  8]) >> SAR[5:0]
  qz[ 23: 16] = (qx[ 23: 16] * qy[ 23: 16]) >> SAR[5:0]
  qz[ 31: 24] = (qx[ 31: 24] * qy[ 31: 24]) >> SAR[5:0]
  qz[ 39: 32] = (qx[ 39: 32] * qy[ 39: 32]) >> SAR[5:0]
  qz[ 47: 40] = (qx[ 47: 40] * qy[ 47: 40]) >> SAR[5:0]
  qz[ 55: 48] = (qx[ 55: 48] * qy[ 55: 48]) >> SAR[5:0]
  qz[ 63: 56] = (qx[ 63: 56] * qy[ 63: 56]) >> SAR[5:0]
  qz[ 71: 64] = (qx[ 71: 64] * qy[ 71: 64]) >> SAR[5:0]
  qz[ 79: 72] = (qx[ 79: 72] * qy[ 79: 72]) >> SAR[5:0]
  qz[ 87: 80] = (qx[ 87: 80] * qy[ 87: 80]) >> SAR[5:0]
  qz[ 95: 88] = (qx[ 95: 88] * qy[ 95: 88]) >> SAR[5:0]
  qz[103: 96] = (qx[103: 96] * qy[103: 96]) >> SAR[5:0]
  qz[111:104] = (qx[111:104] * qy[111:104]) >> SAR[5:0]
  qz[119:112] = (qx[119:112] * qy[119:112]) >> SAR[5:0]
  qz[127:120] = (qx[127:120] * qy[127:120]) >> SAR[5:0]

  qv[127:0] => store128({as[31:4],4{0}})
  as[31:0] = as[31:0] + 16
\end{lstlisting}

\clearpage
\subsection{EE.VMUL.U16}\label{instruction:EEVMULU16}
\setlength\tabcolsep{6pt}
\textbf{Instruction Word} \\
\begin{tabular}{|c|c|c|c|c|c|c|c|c|}
	\hline
	10 & qz[2:1] & 1110 & qz[0] & qy[2] & 1 & qy[1:0] & qx[2:0] & 10100100 \\
	\hline
\end{tabular}

\textbf{Assembler Syntax} \\
EE.VMUL.U16 qz, qx, qy \\
\textbf{Description} \\
This instruction performs an unsigned vector multiplication on 16-bit data. Registers \lstinline{qx} and \lstinline{qy} are the multiplier and the multiplicand respectively. The eight 32-bit data results obtained from the calculation is logically right-shifted by the value in special register \lstinline{SAR}. Then, the lower 16-bit data of the shift result is written into corresponding segment of register \lstinline{qz}.\\
\textbf{Operation}
\begin{lstlisting}
  qz[ 15:  0] = (qx[ 15:  0] * qy[ 15:  0]) >> SAR[5:0]
  qz[ 31: 16] = (qx[ 31: 16] * qy[ 31: 16]) >> SAR[5:0]
  qz[ 47: 32] = (qx[ 47: 32] * qy[ 47: 32]) >> SAR[5:0]
  qz[ 63: 48] = (qx[ 63: 48] * qy[ 63: 48]) >> SAR[5:0]
  qz[ 79: 64] = (qx[ 79: 64] * qy[ 79: 64]) >> SAR[5:0]
  qz[ 95: 80] = (qx[ 95: 80] * qy[ 95: 80]) >> SAR[5:0]
  qz[111: 96] = (qx[111: 96] * qy[111: 96]) >> SAR[5:0]
  qz[127:112] = (qx[127:112] * qy[127:112]) >> SAR[5:0]
\end{lstlisting}

\clearpage
\subsection{EE.VMUL.U16.LD.INCP}\label{instruction:EEVMULU16LDINCP}
\setlength\tabcolsep{6pt}
\textbf{Instruction Word} \\
\begin{tabular}{|c|c|c|c|c|c|c|c|c|c|c|c|}
	\hline
	111000 & qu[2:1] & qy[0] & 101 & qu[0] & qz[2:0] & qx[1:0] & qy[2:1] & 1100 & as[3:0] & 111 & qx[2] \\
	\hline
\end{tabular}

\textbf{Assembler Syntax} \\
EE.VMUL.U16.LD.INCP qu, as, qz, qx, qy \\
\textbf{Description} \\
This instruction performs an unsigned vector multiplication on 16-bit data. Registers \lstinline{qx} and \lstinline{qy} are the multiplier and the multiplicand respectively. The eight 32-bit data results obtained from the calculation is logically right-shifted by the value in special register \lstinline{SAR}. Then, the lower 16-bit data of the shift result is written into corresponding segment of register \lstinline{qz}. \\ During the operation, the lower 4 bits of the access address in register \lstinline{as} are forced to be 0, and then the 16-byte data is loaded from the memory to register \lstinline{qu}. After the access, the value in register \lstinline{as} is incremented by 16.\\
\textbf{Operation}
\begin{lstlisting}
  qz[ 15:  0] = (qx[ 15:  0] * qy[ 15:  0]) >> SAR[5:0]
  qz[ 31: 16] = (qx[ 31: 16] * qy[ 31: 16]) >> SAR[5:0]
  qz[ 47: 32] = (qx[ 47: 32] * qy[ 47: 32]) >> SAR[5:0]
  qz[ 63: 48] = (qx[ 63: 48] * qy[ 63: 48]) >> SAR[5:0]
  qz[ 79: 64] = (qx[ 79: 64] * qy[ 79: 64]) >> SAR[5:0]
  qz[ 95: 80] = (qx[ 95: 80] * qy[ 95: 80]) >> SAR[5:0]
  qz[111: 96] = (qx[111: 96] * qy[111: 96]) >> SAR[5:0]
  qz[127:112] = (qx[127:112] * qy[127:112]) >> SAR[5:0]

  qu[127:0] = load128({as[31:4],4{0}})
  as[31:0] = as[31:0] + 16
\end{lstlisting}

\clearpage
\subsection{EE.VMUL.U16.ST.INCP}\label{instruction:EEVMULU16STINCP}
\setlength\tabcolsep{6pt}
\textbf{Instruction Word} \\
\begin{tabular}{|c|c|c|c|c|c|c|c|c|c|c|}
	\hline
	11100101 & qy[0] & qv[2:0] & 1 & qz[2:0] & qx[1:0] & qy[2:1] & 0011 & as[3:0] & 111 & qx[2] \\
	\hline
\end{tabular}

\textbf{Assembler Syntax} \\
EE.VMUL.U16.ST.INCP qv, as, qz, qx, qy \\
\textbf{Description} \\
This instruction performs an unsigned vector multiplication on 16-bit data. Registers \lstinline{qx} and \lstinline{qy} are the multiplier and the multiplicand respectively. The eight 32-bit data results obtained from the calculation is logically right-shifted by the value in special register \lstinline{SAR}. Then, the lower 16-bit data of the shift result is written into corresponding segment of register \lstinline{qz}. \\ During the operation, the instruction forces the lower 4 bits of the access address in register \lstinline{as} to 0 and stores the value in register \lstinline{qv} to memory. After the access, the value in register \lstinline{as} is incremented by 16.\\
\textbf{Operation}
\begin{lstlisting}
  qz[ 15:  0] = (qx[ 15:  0] * qy[ 15:  0]) >> SAR[5:0]
  qz[ 31: 16] = (qx[ 31: 16] * qy[ 31: 16]) >> SAR[5:0]
  qz[ 47: 32] = (qx[ 47: 32] * qy[ 47: 32]) >> SAR[5:0]
  qz[ 63: 48] = (qx[ 63: 48] * qy[ 63: 48]) >> SAR[5:0]
  qz[ 79: 64] = (qx[ 79: 64] * qy[ 79: 64]) >> SAR[5:0]
  qz[ 95: 80] = (qx[ 95: 80] * qy[ 95: 80]) >> SAR[5:0]
  qz[111: 96] = (qx[111: 96] * qy[111: 96]) >> SAR[5:0]
  qz[127:112] = (qx[127:112] * qy[127:112]) >> SAR[5:0]

  qv[127:0] => store128({as[31:4],4{0}})
  as[31:0] = as[31:0] + 16
\end{lstlisting}

\clearpage
\subsection{EE.VMUL.U8}\label{instruction:EEVMULU8}
\setlength\tabcolsep{6pt}
\textbf{Instruction Word} \\
\begin{tabular}{|c|c|c|c|c|c|c|c|c|}
	\hline
	10 & qz[2:1] & 1110 & qz[0] & qy[2] & 1 & qy[1:0] & qx[2:0] & 10110100 \\
	\hline
\end{tabular}

\textbf{Assembler Syntax} \\
EE.VMUL.U8 qz, qx, qy \\
\textbf{Description} \\
This instruction performs an unsigned vector multiplication on 8-bit data. Registers \lstinline{qx} and \lstinline{qy} are the multiplier and the multiplicand respectively. The 16 16-bit data results obtained from the calculation is logically right-shifted by the value in special register \lstinline{SAR}. Then, the lower 8-bit data of the shift result is written into corresponding segment of register \lstinline{qz}.\\
\textbf{Operation}
\begin{lstlisting}
  qz[  7:  0] = (qx[  7:  0] * qy[  7:  0]) >> SAR[5:0]
  qz[ 15:  8] = (qx[ 15:  8] * qy[ 15:  8]) >> SAR[5:0]
  qz[ 23: 16] = (qx[ 23: 16] * qy[ 23: 16]) >> SAR[5:0]
  qz[ 31: 24] = (qx[ 31: 24] * qy[ 31: 24]) >> SAR[5:0]
  qz[ 39: 32] = (qx[ 39: 32] * qy[ 39: 32]) >> SAR[5:0]
  qz[ 47: 40] = (qx[ 47: 40] * qy[ 47: 40]) >> SAR[5:0]
  qz[ 55: 48] = (qx[ 55: 48] * qy[ 55: 48]) >> SAR[5:0]
  qz[ 63: 56] = (qx[ 63: 56] * qy[ 63: 56]) >> SAR[5:0]
  qz[ 71: 64] = (qx[ 71: 64] * qy[ 71: 64]) >> SAR[5:0]
  qz[ 79: 72] = (qx[ 79: 72] * qy[ 79: 72]) >> SAR[5:0]
  qz[ 87: 80] = (qx[ 87: 80] * qy[ 87: 80]) >> SAR[5:0]
  qz[ 95: 88] = (qx[ 95: 88] * qy[ 95: 88]) >> SAR[5:0]
  qz[103: 96] = (qx[103: 96] * qy[103: 96]) >> SAR[5:0]
  qz[111:104] = (qx[111:104] * qy[111:104]) >> SAR[5:0]
  qz[119:112] = (qx[119:112] * qy[119:112]) >> SAR[5:0]
  qz[127:120] = (qx[127:120] * qy[127:120]) >> SAR[5:0]
\end{lstlisting}

\clearpage
\subsection{EE.VMUL.U8.LD.INCP}\label{instruction:EEVMULU8LDINCP}
\setlength\tabcolsep{6pt}
\textbf{Instruction Word} \\
\begin{tabular}{|c|c|c|c|c|c|c|c|c|c|c|c|}
	\hline
	111000 & qu[2:1] & qy[0] & 110 & qu[0] & qz[2:0] & qx[1:0] & qy[2:1] & 1100 & as[3:0] & 111 & qx[2] \\
	\hline
\end{tabular}

\textbf{Assembler Syntax} \\
EE.VMUL.U8.LD.INCP qu, as, qz, qx, qy \\
\textbf{Description} \\
This instruction performs an unsigned vector multiplication on 8-bit data. Registers \lstinline{qx} and \lstinline{qy} are the multiplier and the multiplicand respectively. The 16 32-bit data results obtained from the calculation is logically right-shifted by the value in special register \lstinline{SAR}. Then, the lower 8-bit data of the shift result is written into corresponding segment of register \lstinline{qz}. \\ During the operation, the lower 4 bits of the access address in register \lstinline{as} are forced to be 0, and then the 16-byte data is loaded from the memory to register \lstinline{qu}. After the access, the value in register \lstinline{as} is incremented by 16.\\
\textbf{Operation}
\begin{lstlisting}
  qz[  7:  0] = (qx[  7:  0] * qy[  7:  0]) >> SAR[5:0]
  qz[ 15:  8] = (qx[ 15:  8] * qy[ 15:  8]) >> SAR[5:0]
  qz[ 23: 16] = (qx[ 23: 16] * qy[ 23: 16]) >> SAR[5:0]
  qz[ 31: 24] = (qx[ 31: 24] * qy[ 31: 24]) >> SAR[5:0]
  qz[ 39: 32] = (qx[ 39: 32] * qy[ 39: 32]) >> SAR[5:0]
  qz[ 47: 40] = (qx[ 47: 40] * qy[ 47: 40]) >> SAR[5:0]
  qz[ 55: 48] = (qx[ 55: 48] * qy[ 55: 48]) >> SAR[5:0]
  qz[ 63: 56] = (qx[ 63: 56] * qy[ 63: 56]) >> SAR[5:0]
  qz[ 71: 64] = (qx[ 71: 64] * qy[ 71: 64]) >> SAR[5:0]
  qz[ 79: 72] = (qx[ 79: 72] * qy[ 79: 72]) >> SAR[5:0]
  qz[ 87: 80] = (qx[ 87: 80] * qy[ 87: 80]) >> SAR[5:0]
  qz[ 95: 88] = (qx[ 95: 88] * qy[ 95: 88]) >> SAR[5:0]
  qz[103: 96] = (qx[103: 96] * qy[103: 96]) >> SAR[5:0]
  qz[111:104] = (qx[111:104] * qy[111:104]) >> SAR[5:0]
  qz[119:112] = (qx[119:112] * qy[119:112]) >> SAR[5:0]
  qz[127:120] = (qx[127:120] * qy[127:120]) >> SAR[5:0]

  qu[127:0] = load128({as[31:4],4{0}})
  as[31:0] = as[31:0] + 16
\end{lstlisting}

\clearpage
\subsection{EE.VMUL.U8.ST.INCP}\label{instruction:EEVMULU8STINCP}
\setlength\tabcolsep{6pt}
\textbf{Instruction Word} \\
\begin{tabular}{|c|c|c|c|c|c|c|c|c|c|c|}
	\hline
	11101000 & qy[0] & qv[2:0] & 1 & qz[2:0] & qx[1:0] & qy[2:1] & 0000 & as[3:0] & 111 & qx[2] \\
	\hline
\end{tabular}

\textbf{Assembler Syntax} \\
EE.VMUL.U8.ST.INCP qv, as, qz, qx, qy \\
\textbf{Description} \\
This instruction performs an unsigned vector multiplication on 8-bit data. Registers \lstinline{qx} and \lstinline{qy} are the multiplier and the multiplicand respectively. The 16 16-bit data results obtained from the calculation is logically right-shifted by the value in special register \lstinline{SAR}. Then, the lower 8-bit data of the shift result is written into corresponding segment of register \lstinline{qz}. \\ During the operation, the instruction forces the lower 4 bits of the access address in register \lstinline{as} to 0 and stores the value in register \lstinline{qv} to memory. After the access, the value in register \lstinline{as} is incremented by 16.\\
\textbf{Operation}
\begin{lstlisting}
  qz[  7:  0] = (qx[  7:  0] * qy[  7:  0]) >> SAR[5:0]
  qz[ 15:  8] = (qx[ 15:  8] * qy[ 15:  8]) >> SAR[5:0]
  qz[ 23: 16] = (qx[ 23: 16] * qy[ 23: 16]) >> SAR[5:0]
  qz[ 31: 24] = (qx[ 31: 24] * qy[ 31: 24]) >> SAR[5:0]
  qz[ 39: 32] = (qx[ 39: 32] * qy[ 39: 32]) >> SAR[5:0]
  qz[ 47: 40] = (qx[ 47: 40] * qy[ 47: 40]) >> SAR[5:0]
  qz[ 55: 48] = (qx[ 55: 48] * qy[ 55: 48]) >> SAR[5:0]
  qz[ 63: 56] = (qx[ 63: 56] * qy[ 63: 56]) >> SAR[5:0]
  qz[ 71: 64] = (qx[ 71: 64] * qy[ 71: 64]) >> SAR[5:0]
  qz[ 79: 72] = (qx[ 79: 72] * qy[ 79: 72]) >> SAR[5:0]
  qz[ 87: 80] = (qx[ 87: 80] * qy[ 87: 80]) >> SAR[5:0]
  qz[ 95: 88] = (qx[ 95: 88] * qy[ 95: 88]) >> SAR[5:0]
  qz[103: 96] = (qx[103: 96] * qy[103: 96]) >> SAR[5:0]
  qz[111:104] = (qx[111:104] * qy[111:104]) >> SAR[5:0]
  qz[119:112] = (qx[119:112] * qy[119:112]) >> SAR[5:0]
  qz[127:120] = (qx[127:120] * qy[127:120]) >> SAR[5:0]

  qv[127:0] => store128({as[31:4],4{0}})
  as[31:0] = as[31:0] + 16
\end{lstlisting}

\clearpage
\subsection{EE.VMULAS.S16.ACCX}\label{instruction:EEVMULASS16ACCX}
\setlength\tabcolsep{6pt}
\textbf{Instruction Word} \\
\begin{tabular}{|c|c|c|c|c|c|}
	\hline
	000110100 & qy[2] & 0 & qy[1:0] & qx[2:0] & 10000100 \\
	\hline
\end{tabular}

\textbf{Assembler Syntax} \\
EE.VMULAS.S16.ACCX qx, qy \\
\textbf{Description} \\
This instruction divides registers \lstinline{qx} and \lstinline{qy} into 8 data segments by 16 bits. Then, it performs a signed multiply-accumulate operation on the 8 sets of segments respectively. The accumulated result is added to the value in special register \lstinline{ACCX}. Then, the sum obtained is saturated and then stored in \lstinline{ACCX}.\\
\textbf{Operation}
\begin{lstlisting}
  add0[31:0] = qx[ 15:  0] * qy[ 15:  0]
  add1[31:0] = qx[ 31: 16] * qy[ 31: 16]
  ...
  add7[31:0] = qx[127:112] * qy[127:112]
  sum[40:0] = ACCX[39:0] + ad[31:0]d0[31:0] + ad[31:0]d1[31:0] + ... + ad[31:0]d7[31:0]

  ACCX[39:0] = min(max(sum[40:0], -2^{39}), 2^{39}-1)
\end{lstlisting}

\clearpage
\subsection{EE.VMULAS.S16.ACCX.LD.IP}\label{instruction:EEVMULASS16ACCXLDIP}
\setlength\tabcolsep{6pt}
\textbf{Instruction Word} \\
\begin{tabular}{|c|c|c|c|c|c|c|c|c|c|c|c|c|}
	\hline
	1111 & imm16[5:4] & qu[2:1] & qy[0] & 000 & qu[0] & 000 & qx[1:0] & qy[2:1] & imm16[3:0] & as[3:0] & 111 & qx[2] \\
	\hline
\end{tabular}

\textbf{Assembler Syntax} \\
EE.VMULAS.S16.ACCX.LD.IP qu, as, -512..496, qx, qy \\
\textbf{Description} \\
This instruction divides registers \lstinline{qx} and \lstinline{qy} into 8 data segments by 16 bits. Then, it performs a signed multiply-accumulate operation on the 8 sets of segments respectively. The accumulated result is added to the value in special register \lstinline{ACCX}. Then, the sum obtained is saturated and then stored in \lstinline{ACCX}. \\ During the operation, the lower 4 bits of the access address in register \lstinline{as} are forced to be 0, and then the 16-byte data is loaded from the memory to register \lstinline{qu}. After the access is completed, the value in register \lstinline{as} is incremented by 6-bit sign-extended constant in the instruction code segment left-shifted by 4.\\
\textbf{Operation}
\begin{lstlisting}
  add0[31:0] = qx[ 15:  0] * qy[ 15:  0]
  add1[31:0] = qx[ 31: 16] * qy[ 31: 16]
  ...
  add7[31:0] = qx[127:112] * qy[127:112]
  sum[40:0] = ACCX[39:0] + ad[31:0]d0[31:0] + ad[31:0]d1[31:0] + ... + ad[31:0]d7[31:0]
  ACCX[39:0] = min(max(sum[40:0], -2^{39}), 2^{39}-1)

  qu[127:0] = load128({as[31:4],4{0}})
  as[31:0] = as[31:0] + {20{imm16[7]},imm16[7:0],4{0}}
\end{lstlisting}

\clearpage
\subsection{EE.VMULAS.S16.ACCX.LD.IP.QUP}\label{instruction:EEVMULASS16ACCXLDIPQUP}
\setlength\tabcolsep{5pt}
\textbf{Instruction Word} \\
\begin{tabular}{|c|c|c|c|c|c|c|c|c|c|c|c|c|}
	\hline
	0000 & imm16[5:4] & qu[2:1] & qy[0] & qs0[2:0] & qu[0] & qs1[2:0] & qx[1:0] & qy[2:1] & imm16[3:0] & as[3:0] & 111 & qx[2] \\
	\hline
\end{tabular}

\textbf{Assembler Syntax} \\
EE.VMULAS.S16.ACCX.LD.IP.QUP qu, as, imm16, qx, qy, qs0, qs1 \\
\textbf{Description} \\
This instruction divides registers \lstinline{qx} and \lstinline{qy} into 8 data segments by 16 bits. Then, it performs a signed multiply-accumulate operation on the 8 sets of segments respectively. The accumulated result is added to the value in special register \lstinline{ACCX}. Then, the sum obtained is saturated and then stored in \lstinline{ACCX}.\\
During the operation, the lower 4 bits of the access address in register \lstinline{as} are forced to be 0, and then the 16-byte data is loaded from the memory to register \lstinline{qu}. After the access is completed, the value in register \lstinline{as} is incremented by 6-bit sign-extended constant in the instruction code segment left-shifted by 4.\\
At the same time, this instruction also obtains 16-byte unaligned data by concatenating and shifting consecutive aligned data stored in the two registers \lstinline{qs0} and \lstinline{qs1} and stores it to \lstinline{qs0}. The shift byte is stored in special register \lstinline{SAR_BYTE}.\\
\textbf{Operation}
\begin{lstlisting}
  add0[31:0] = qx[ 15:  0] * qy[ 15:  0]
  add1[31:0] = qx[ 31: 16] * qy[ 31: 16]
  ...
  add7[31:0] = qx[127:112] * qy[127:112]
  sum[40:0] = ACCX[39:0] + ad[31:0]d0[31:0] + ad[31:0]d1[31:0] + ... + ad[31:0]d7[31:0]
  ACCX[39:0] = min(max(sum[40:0], -2^{39}), 2^{39}-1)

  qu[127:0] = load128({as[31:4],4{0}})
  as[31:0] = as[31:0] + {20{imm16[7]},imm16[7:0],4{0}}
  qs0[127:0]  = {qs1[127:0], qs0[127:0]} >> {SAR_BYTE[3:0] << 3}
\end{lstlisting}

\clearpage
\subsection{EE.VMULAS.S16.ACCX.LD.XP}\label{instruction:EEVMULASS16ACCXLDXP}
\setlength\tabcolsep{6pt}
\textbf{Instruction Word} \\
\begin{tabular}{|c|c|c|c|c|c|c|c|c|c|c|c|}
	\hline
	111100 & qu[2:1] & qy[0] & 001 & qu[0] & 000 & qx[1:0] & qy[2:1] & ad[3:0] & as[3:0] & 111 & qx[2] \\
	\hline
\end{tabular}

\textbf{Assembler Syntax} \\
EE.VMULAS.S16.ACCX.LD.XP qu, as, ad, qx, qy \\
\textbf{Description} \\
This instruction divides registers \lstinline{qx} and \lstinline{qy} into 8 data segments by 16 bits. Then, it performs a signed multiply-accumulate operation on the 8 sets of segments respectively. The accumulated result is added to the value in special register \lstinline{ACCX}. Then, the sum obtained is saturated and then stored in \lstinline{ACCX}.\\
During the operation, the lower 4 bits of the access address in register \lstinline{as} are forced to be 0, and then the 16-byte data is loaded from the memory to register \lstinline{qu}. After the access is completed, the value in register \lstinline{as} is incremented by the value in register \lstinline{ad}.\\
\textbf{Operation}
\begin{lstlisting}
  add0[31:0] = qx[ 15:  0] * qy[ 15:  0]
  add1[31:0] = qx[ 31: 16] * qy[ 31: 16]
  ...
  add7[31:0] = qx[127:112] * qy[127:112]
  sum[40:0] = ACCX[39:0] + ad[31:0]d0[31:0] + ad[31:0]d1[31:0] + ... + ad[31:0]d7[31:0]
  ACCX[39:0] = min(max(sum[40:0], -2^{39}), 2^{39}-1)

  qu[127:0] = load128({as[31:4],4{0}})
  as[31:0] = as[31:0] + ad[31:0]
\end{lstlisting}

\clearpage
\subsection{EE.VMULAS.S16.ACCX.LD.XP.QUP}\label{instruction:EEVMULASS16ACCXLDXPQUP}
\setlength\tabcolsep{6pt}
\textbf{Instruction Word} \\
\begin{tabular}{|c|c|c|c|c|c|c|c|c|c|c|c|}
	\hline
	101100 & qu[2:1] & qy[0] & qs0[2:0] & qu[0] & qs1[2:0] & qx[1:0] & qy[2:1] & ad[3:0] & as[3:0] & 111 & qx[2] \\
	\hline
\end{tabular}

\textbf{Assembler Syntax} \\
EE.VMULAS.S16.ACCX.LD.XP.QUP qu, as, ad, qx, qy, qs0, qs1 \\
\textbf{Description} \\
This instruction divides registers \lstinline{qx} and \lstinline{qy} into 8 data segments by 16 bits. Then, it performs a signed multiply-accumulate operation on the 8 sets of segments respectively. The accumulated result is added to the value in special register \lstinline{ACCX}. Then, the sum obtained is saturated and then stored in \lstinline{ACCX}.\\
During the operation, the lower 4 bits of the access address in register \lstinline{as} are forced to be 0, and then the 16-byte data is loaded from the memory to register \lstinline{qu}. After the access is completed, the value in register \lstinline{as} is incremented by the value in register \lstinline{ad}.\\
At the same time, this instruction also obtains 16-byte unaligned data by concatenating and shifting consecutive aligned data stored in the two registers \lstinline{qs0} and \lstinline{qs1} and stores it to \lstinline{qs0}. The shift byte is stored in special register \lstinline{SAR_BYTE}.\\
\textbf{Operation}
\begin{lstlisting}
  add0[31:0] = qx[ 15:  0] * qy[ 15:  0]
  add1[31:0] = qx[ 31: 16] * qy[ 31: 16]
  ...
  add7[31:0] = qx[127:112] * qy[127:112]
  sum[40:0] = ACCX[39:0] + ad[31:0]d0[31:0] + ad[31:0]d1[31:0] + ... + ad[31:0]d7[31:0]
  ACCX[39:0] = min(max(sum[40:0], -2^{39}), 2^{39}-1)

  qu[127:0] = load128({as[31:4],4{0}})
  as[31:0] = as[31:0] + ad[31:0]
  qs0[127:0] = {qs1[127: 0], qs0[127: 0]} >> {SAR_BYTE[3:0] << 3}
\end{lstlisting}

\clearpage
\subsection{EE.VMULAS.S16.QACC}\label{instruction:EEVMULASS16QACC}
\setlength\tabcolsep{6pt}
\textbf{Instruction Word} \\
\begin{tabular}{|c|c|c|c|c|c|}
	\hline
	000110100 & qy[2] & 1 & qy[1:0] & qx[2:0] & 10000100 \\
	\hline
\end{tabular}

\textbf{Assembler Syntax} \\
EE.VMULAS.S16.QACC qx, qy \\
\textbf{Description} \\
This instruction divides registers \lstinline{qx} and \lstinline{qy} into 8 data segments by 16 bits. Then, the signed multiplication result of 8 sets of segments is added to the corresponding 40-bit data segment in special registers \lstinline{QACC_H} and \lstinline{QACC_L} respectively. The calculated result is saturated to a 40-bit signed number and then stored to the corresponding 40-bit data segment in \lstinline{QACC_H} and \lstinline{QACC_L}.\\
\textbf{Operation}
\begin{lstlisting}
  QACC_L[ 39:  0] = min(max(QACC_L[ 39:  0] + qx[ 15:  0] * qy[ 15:  0], -2^{39}), 2^{39}-1)
  QACC_L[ 79: 40] = min(max(QACC_L[ 79: 40] + qx[ 31: 16] * qy[ 31: 16], -2^{39}), 2^{39}-1)
  QACC_L[119: 80] = min(max(QACC_L[119: 80] + qx[ 47: 32] * qy[ 47: 32], -2^{39}), 2^{39}-1)
  QACC_L[159:120] = min(max(QACC_L[159:120] + qx[ 63: 48] * qy[ 63: 48], -2^{39}), 2^{39}-1)
  QACC_H[ 39:  0] = min(max(QACC_H[ 39:  0] + qx[ 79: 64] * qy[ 79: 64], -2^{39}), 2^{39}-1)
  QACC_H[ 79: 40] = min(max(QACC_H[ 79: 40] + qx[ 95: 80] * qy[ 95: 80], -2^{39}), 2^{39}-1)
  QACC_H[119: 80] = min(max(QACC_H[119: 80] + qx[111: 96] * qy[111: 96], -2^{39}), 2^{39}-1)
  QACC_H[159:120] = min(max(QACC_H[159:120] + qx[127:112] * qy[127:112], -2^{39}), 2^{39}-1)
\end{lstlisting}

\clearpage
\subsection{EE.VMULAS.S16.QACC.LD.IP}\label{instruction:EEVMULASS16QACCLDIP}
\setlength\tabcolsep{6pt}
\textbf{Instruction Word} \\
\begin{tabular}{|c|c|c|c|c|c|c|c|c|c|c|c|c|}
	\hline
	1111 & imm16[5:4] & qu[2:1] & qy[0] & 000 & qu[0] & 001 & qx[1:0] & qy[2:1] & imm16[3:0] & as[3:0] & 111 & qx[2] \\
	\hline
\end{tabular}

\textbf{Assembler Syntax} \\
EE.VMULAS.S16.QACC.LD.IP qu, as, imm16, qx, qy \\
\textbf{Description} \\
This instruction divides registers \lstinline{qx} and \lstinline{qy} into 8 data segments by 16 bits. Then, the signed multiplication result of 8 sets of segments is added to the corresponding 40-bit data segment in special registers \lstinline{QACC_H} and \lstinline{QACC_L} respectively. The calculated result is saturated to a 40-bit signed number and then stored to the corresponding 40-bit data segment in \lstinline{QACC_H} and \lstinline{QACC_L}.\\
During the operation, the lower 4 bits of the access address in register \lstinline{as} are forced to be 0, and then the 16-byte data is loaded from the memory to register \lstinline{qu}. After the access is completed, the value in register \lstinline{as} is incremented by 6-bit sign-extended constant in the instruction code segment left-shifted by 4.\\
\textbf{Operation}
\begin{lstlisting}
  QACC_L[ 39:  0] = min(max(QACC_L[ 39:  0] + qx[ 15:  0] * qy[ 15:  0], -2^{39}), 2^{39}-1)
  QACC_L[ 79: 40] = min(max(QACC_L[ 79: 40] + qx[ 31: 16] * qy[ 31: 16], -2^{39}), 2^{39}-1)
  QACC_L[119: 80] = min(max(QACC_L[119: 80] + qx[ 47: 32] * qy[ 47: 32], -2^{39}), 2^{39}-1)
  QACC_L[159:120] = min(max(QACC_L[159:120] + qx[ 63: 48] * qy[ 63: 48], -2^{39}), 2^{39}-1)
  QACC_H[ 39:  0] = min(max(QACC_H[ 39:  0] + qx[ 79: 64] * qy[ 79: 64], -2^{39}), 2^{39}-1)
  QACC_H[ 79: 40] = min(max(QACC_H[ 79: 40] + qx[ 95: 80] * qy[ 95: 80], -2^{39}), 2^{39}-1)
  QACC_H[119: 80] = min(max(QACC_H[119: 80] + qx[111: 96] * qy[111: 96], -2^{39}), 2^{39}-1)
  QACC_H[159:120] = min(max(QACC_H[159:120] + qx[127:112] * qy[127:112], -2^{39}), 2^{39}-1)

  qu[127:0] = load128({as[31:4],4{0}})
  as[31:0] = as[31:0] + {22{imm16[5]},imm16[5:0],4{0}}
\end{lstlisting}

\clearpage
\subsection{EE.VMULAS.S16.QACC.LD.IP.QUP}\label{instruction:EEVMULASS16QACCLDIPQUP}
\setlength\tabcolsep{5pt}
\textbf{Instruction Word} \\
\begin{tabular}{|c|c|c|c|c|c|c|c|c|c|c|c|c|}
	\hline
	0001 & imm16[5:4] & qu[2:1] & qy[0] & qs0[2:0] & qu[0] & qs1[2:0] & qx[1:0] & qy[2:1] & imm16[3:0] & as[3:0] & 111 & qx[2] \\
	\hline
\end{tabular}

\textbf{Assembler Syntax} \\
EE.VMULAS.S16.QACC.LD.IP.QUP qu, as, imm16, qx, qy, qs0, qs1 \\
\textbf{Description} \\
This instruction divides registers \lstinline{qx} and \lstinline{qy} into 8 data segments by 16 bits. Then, the signed multiplication result of 8 sets of segments is added to the corresponding 40-bit data segment in special registers \lstinline{QACC_H} and \lstinline{QACC_L} respectively. The calculated result is saturated to a 40-bit signed number and then stored to the corresponding 40-bit data segment in \lstinline{QACC_H} and \lstinline{QACC_L}.\\
During the operation, the lower 4 bits of the access address in register \lstinline{as} are forced to be 0, and then the 16-byte data is loaded from the memory to register \lstinline{qu}. After the access is completed, the value in register \lstinline{as} is incremented by 6-bit sign-extended constant in the instruction code segment left-shifted by 4.\\
At the same time, this instruction also obtains 16-byte unaligned data by concatenating and shifting consecutive aligned data stored in the two registers \lstinline{qs0} and \lstinline{qs1} and stores it to \lstinline{qs0}. The shift byte is stored in special register \lstinline{SAR_BYTE}.\\
\textbf{Operation}
\begin{lstlisting}
  QACC_L[ 39:  0] = min(max(QACC_L[ 39:  0] + qx[ 15:  0] * qy[ 15:  0], -2^{39}), 2^{39}-1)
  QACC_L[ 79: 40] = min(max(QACC_L[ 79: 40] + qx[ 31: 16] * qy[ 31: 16], -2^{39}), 2^{39}-1)
  QACC_L[119: 80] = min(max(QACC_L[119: 80] + qx[ 47: 32] * qy[ 47: 32], -2^{39}), 2^{39}-1)
  QACC_L[159:120] = min(max(QACC_L[159:120] + qx[ 63: 48] * qy[ 63: 48], -2^{39}), 2^{39}-1)
  QACC_H[ 39:  0] = min(max(QACC_H[ 39:  0] + qx[ 79: 64] * qy[ 79: 64], -2^{39}), 2^{39}-1)
  QACC_H[ 79: 40] = min(max(QACC_H[ 79: 40] + qx[ 95: 80] * qy[ 95: 80], -2^{39}), 2^{39}-1)
  QACC_H[119: 80] = min(max(QACC_H[119: 80] + qx[111: 96] * qy[111: 96], -2^{39}), 2^{39}-1)
  QACC_H[159:120] = min(max(QACC_H[159:120] + qx[127:112] * qy[127:112], -2^{39}), 2^{39}-1)

  qu[127:0] = load128({as[31:4],4{0}})
  as[31:0] = as[31:0] + {22{imm16[5]},imm16[5:0],4{0}}
  qs0[127:0]  = {qs1[127:0], qs0[127:0]} >> {SAR_BYTE[3:0] << 3}
\end{lstlisting}

\clearpage
\subsection{EE.VMULAS.S16.QACC.LD.XP}\label{instruction:EEVMULASS16QACCLDXP}
\setlength\tabcolsep{6pt}
\textbf{Instruction Word} \\
\begin{tabular}{|c|c|c|c|c|c|c|c|c|c|c|c|}
	\hline
	111100 & qu[2:1] & qy[0] & 001 & qu[0] & 001 & qx[1:0] & qy[2:1] & ad[3:0] & as[3:0] & 111 & qx[2] \\
	\hline
\end{tabular}

\textbf{Assembler Syntax} \\
EE.VMULAS.S16.QACC.LD.XP qu, as, ad, qx, qy \\
\textbf{Description} \\
This instruction divides registers \lstinline{qx} and \lstinline{qy} into 8 data segments by 16 bits. Then, the signed multiplication result of 8 sets of segments is added to the corresponding 40-bit data segment in special registers \lstinline{QACC_H} and \lstinline{QACC_L} respectively. The calculated result is saturated to a 40-bit signed number and then stored to the corresponding 40-bit data segment in \lstinline{QACC_H} and \lstinline{QACC_L}.\\
During the operation, the lower 4 bits of the access address in register \lstinline{as} are forced to be 0, and then the 16-byte data is loaded from the memory to register \lstinline{qu}. After the access is completed, the value in register \lstinline{as} is incremented by the value in register \lstinline{ad}.\\
\textbf{Operation}
\begin{lstlisting}
  QACC_L[ 39:  0] = min(max(QACC_L[ 39:  0] + qx[ 15:  0] * qy[ 15:  0], -2^{39}), 2^{39}-1)
  QACC_L[ 79: 40] = min(max(QACC_L[ 79: 40] + qx[ 31: 16] * qy[ 31: 16], -2^{39}), 2^{39}-1)
  QACC_L[119: 80] = min(max(QACC_L[119: 80] + qx[ 47: 32] * qy[ 47: 32], -2^{39}), 2^{39}-1)
  QACC_L[159:120] = min(max(QACC_L[159:120] + qx[ 63: 48] * qy[ 63: 48], -2^{39}), 2^{39}-1)
  QACC_H[ 39:  0] = min(max(QACC_H[ 39:  0] + qx[ 79: 64] * qy[ 79: 64], -2^{39}), 2^{39}-1)
  QACC_H[ 79: 40] = min(max(QACC_H[ 79: 40] + qx[ 95: 80] * qy[ 95: 80], -2^{39}), 2^{39}-1)
  QACC_H[119: 80] = min(max(QACC_H[119: 80] + qx[111: 96] * qy[111: 96], -2^{39}), 2^{39}-1)
  QACC_H[159:120] = min(max(QACC_H[159:120] + qx[127:112] * qy[127:112], -2^{39}), 2^{39}-1)

  qu[127:0] = load128({as[31:4],4{0}})
  as[31:0] = as[31:0] + ad[31:0]
\end{lstlisting}

\clearpage
\subsection{EE.VMULAS.S16.QACC.LD.XP.QUP}\label{instruction:EEVMULASS16QACCLDXPQUP}
\setlength\tabcolsep{6pt}
\textbf{Instruction Word} \\
\begin{tabular}{|c|c|c|c|c|c|c|c|c|c|c|c|}
	\hline
	101101 & qu[2:1] & qy[0] & qs0[2:0] & qu[0] & qs1[2:0] & qx[1:0] & qy[2:1] & ad[3:0] & as[3:0] & 111 & qx[2] \\
	\hline
\end{tabular}

\textbf{Assembler Syntax} \\
EE.VMULAS.S16.QACC.LD.XP.QUP qu, as, ad, qx, qy, qs0, qs1 \\
\textbf{Description} \\
This instruction divides registers \lstinline{qx} and \lstinline{qy} into 8 data segments by 16 bits. Then, the signed multiplication result of 8 sets of segments is added to the corresponding 40-bit data segment in special registers \lstinline{QACC_H} and \lstinline{QACC_L} respectively. The calculated result is saturated to a 40-bit signed number and then stored to the corresponding 40-bit data segment in \lstinline{QACC_H} and \lstinline{QACC_L}.\\
During the operation, the lower 4 bits of the access address in register \lstinline{as} are forced to be 0, and then the 16-byte data is loaded from the memory to register \lstinline{qu}. After the access is completed, the value in register \lstinline{as} is incremented by the value in register \lstinline{ad}.\\
At the same time, this instruction also obtains 16-byte unaligned data by concatenating and shifting consecutive aligned data stored in the two registers \lstinline{qs0} and \lstinline{qs1} and stores it to \lstinline{qs0}. The shift byte is stored in special register \lstinline{SAR_BYTE}.\\
\textbf{Operation}
\begin{lstlisting}
  QACC_L[ 39:  0] = min(max(QACC_L[ 39:  0] + qx[ 15:  0] * qy[ 15:  0], -2^{39}), 2^{39}-1)
  QACC_L[ 79: 40] = min(max(QACC_L[ 79: 40] + qx[ 31: 16] * qy[ 31: 16], -2^{39}), 2^{39}-1)
  QACC_L[119: 80] = min(max(QACC_L[119: 80] + qx[ 47: 32] * qy[ 47: 32], -2^{39}), 2^{39}-1)
  QACC_L[159:120] = min(max(QACC_L[159:120] + qx[ 63: 48] * qy[ 63: 48], -2^{39}), 2^{39}-1)
  QACC_H[ 39:  0] = min(max(QACC_H[ 39:  0] + qx[ 79: 64] * qy[ 79: 64], -2^{39}), 2^{39}-1)
  QACC_H[ 79: 40] = min(max(QACC_H[ 79: 40] + qx[ 95: 80] * qy[ 95: 80], -2^{39}), 2^{39}-1)
  QACC_H[119: 80] = min(max(QACC_H[119: 80] + qx[111: 96] * qy[111: 96], -2^{39}), 2^{39}-1)
  QACC_H[159:120] = min(max(QACC_H[159:120] + qx[127:112] * qy[127:112], -2^{39}), 2^{39}-1)

  qu[127:0] = load128({as[31:4],4{0}})
  as[31:0] = as[31:0] + ad[31:0]
  qs0[127:0]  = {qs1[127:0], qs0[127:0]} >> {SAR_BYTE[3:0] << 3}
\end{lstlisting}

\clearpage
\subsection{EE.VMULAS.S16.QACC.LDBC.INCP}\label{instruction:EEVMULASS16QACCLDBCINCP}
\setlength\tabcolsep{6pt}
\textbf{Instruction Word} \\
\begin{tabular}{|c|c|c|c|c|c|c|c|c|c|}
	\hline
	100 & qu[2] & 0111 & qu[1] & qy[2] & qu[0] & qy[1:0] & qx[2:0] & as[3:0] & 0100 \\
	\hline
\end{tabular}

\textbf{Assembler Syntax} \\
EE.VMULAS.S16.QACC.LDBC.INCP qu, as, qx, qy \\
\textbf{Description} \\
This instruction divides registers \lstinline{qx} and \lstinline{qy} into 8 data segments by 16 bits. Then, the signed multiplication result of 8 sets of segments is added to the corresponding 40-bit data segment in special registers \lstinline{QACC_H} and \lstinline{QACC_L} respectively. The calculated result is saturated to a 40-bit signed number and then stored to the corresponding 40-bit data segment in \lstinline{QACC_H} and \lstinline{QACC_L}.\\
At the same time, this instruction forces the lower 1 bit of the access address in register \lstinline{as} to 0, loads 16-bit data from memory, and broadcasts it to the eight 16-bit data segments in register \lstinline{qu}. After the access, the value in register \lstinline{as} is incremented by 2.\\
\textbf{Operation}
\begin{lstlisting}
  QACC_L[ 39:  0] = min(max(QACC_L[ 39:  0] + qx[ 15:  0] * qy[ 15:  0], -2^{39}), 2^{39}-1)
  QACC_L[ 79: 40] = min(max(QACC_L[ 79: 40] + qx[ 31: 16] * qy[ 31: 16], -2^{39}), 2^{39}-1)
  QACC_L[119: 80] = min(max(QACC_L[119: 80] + qx[ 47: 32] * qy[ 47: 32], -2^{39}), 2^{39}-1)
  QACC_L[159:120] = min(max(QACC_L[159:120] + qx[ 63: 48] * qy[ 63: 48], -2^{39}), 2^{39}-1)
  QACC_H[ 39:  0] = min(max(QACC_H[ 39:  0] + qx[ 79: 64] * qy[ 79: 64], -2^{39}), 2^{39}-1)
  QACC_H[ 79: 40] = min(max(QACC_H[ 79: 40] + qx[ 95: 80] * qy[ 95: 80], -2^{39}), 2^{39}-1)
  QACC_H[119: 80] = min(max(QACC_H[119: 80] + qx[111: 96] * qy[111: 96], -2^{39}), 2^{39}-1)
  QACC_H[159:120] = min(max(QACC_H[159:120] + qx[127:112] * qy[127:112], -2^{39}), 2^{39}-1)

  qu[127:0] = {8{load16({as[31:1],1{0}})}}
  as[31:0] = as[31:0] + 2
\end{lstlisting}

\clearpage
\subsection{EE.VMULAS.S16.QACC.LDBC.INCP.QUP}\label{instruction:EEVMULASS16QACCLDBCINCPQUP}
\setlength\tabcolsep{6pt}
\textbf{Instruction Word} \\
\begin{tabular}{|c|c|c|c|c|c|c|c|c|c|c|c|}
	\hline
	111000 & qu[2:1] & qy[0] & qs0[2:0] & qu[0] & qs1[2:0] & qx[1:0] & qy[2:1] & 1000 & as[3:0] & 111 & qx[2] \\
	\hline
\end{tabular}

\textbf{Assembler Syntax} \\
EE.VMULAS.S16.QACC.LDBC.INCP.QUP qu, as, qx, qy, qs0, qs1 \\
\textbf{Description} \\
This instruction divides registers \lstinline{qx} and \lstinline{qy} into 8 data segments by 16 bits. Then, the signed multiplication result of 8 sets of segments is added to the corresponding 40-bit data segment in special registers \lstinline{QACC_H} and \lstinline{QACC_L} respectively. The calculated result is saturated to a 40-bit signed number and then stored to the corresponding 40-bit data segment in \lstinline{QACC_H} and \lstinline{QACC_L}.\\
At the same time, this instruction forces the lower 1 bit of the access address in register \lstinline{as} to 0, loads 16-bit data from memory, and broadcasts it to the eight 16-bit data segments in register \lstinline{qu}. After the access, the value in register \lstinline{as} is incremented by 2.\\
At the same time, this instruction also obtains 16-byte unaligned data by concatenating and shifting consecutive aligned data stored in the two registers \lstinline{qs0} and \lstinline{qs1} and stores it to \lstinline{qs0}. The shift byte is stored in special register \lstinline{SAR_BYTE}.\\
\textbf{Operation}
\begin{lstlisting}
  QACC_L[ 39:  0] = min(max(QACC_L[ 39:  0] + qx[ 15:  0] * qy[ 15:  0], -2^{39}), 2^{39}-1)
  QACC_L[ 79: 40] = min(max(QACC_L[ 79: 40] + qx[ 31: 16] * qy[ 31: 16], -2^{39}), 2^{39}-1)
  QACC_L[119: 80] = min(max(QACC_L[119: 80] + qx[ 47: 32] * qy[ 47: 32], -2^{39}), 2^{39}-1)
  QACC_L[159:120] = min(max(QACC_L[159:120] + qx[ 63: 48] * qy[ 63: 48], -2^{39}), 2^{39}-1)
  QACC_H[ 39:  0] = min(max(QACC_H[ 39:  0] + qx[ 79: 64] * qy[ 79: 64], -2^{39}), 2^{39}-1)
  QACC_H[ 79: 40] = min(max(QACC_H[ 79: 40] + qx[ 95: 80] * qy[ 95: 80], -2^{39}), 2^{39}-1)
  QACC_H[119: 80] = min(max(QACC_H[119: 80] + qx[111: 96] * qy[111: 96], -2^{39}), 2^{39}-1)
  QACC_H[159:120] = min(max(QACC_H[159:120] + qx[127:112] * qy[127:112], -2^{39}), 2^{39}-1)

  qu[127:0] = {8{load16({as[31:1],1{0}})}}
  as[31:0] = as[31:0] + 2
  qs0[127:0]  = {qs1[127:0], qs0[127:0]} >> {SAR_BYTE[3:0] << 3}
\end{lstlisting}

\clearpage
\subsection{EE.VMULAS.S8.ACCX}\label{instruction:EEVMULASS8ACCX}
\setlength\tabcolsep{6pt}
\textbf{Instruction Word} \\
\begin{tabular}{|c|c|c|c|c|c|}
	\hline
	000110100 & qy[2] & 0 & qy[1:0] & qx[2:0] & 11000100 \\
	\hline
\end{tabular}

\textbf{Assembler Syntax} \\
EE.VMULAS.S8.ACCX qx, qy \\
\textbf{Description} \\
This instruction divides registers \lstinline{qx} and \lstinline{qy} into 16 data segments by 8 bits. Then, it performs a signed multiply-accumulate operation on the 16 sets of segments respectively. The accumulated result is added to the value in special register \lstinline{ACCX}. Then, the sum obtained is saturated and then stored in \lstinline{ACCX}.\\
\textbf{Operation}
\begin{lstlisting}
  add0[15:0] = qx[  7:  0] * qy[  7:  0]
  add1[15:0] = qx[ 15:  8] * qy[ 15:  8]
  ...
  add15[15:0] = qx[127:120] * qy[127:120]
  sum[40:0] = ACCX[39:0] + ad[31:0]d0[15:0] + ad[31:0]d1[15:0] + ... + ad[31:0]d15[15:0]

  ACCX[39:0] = min(max(sum[40:0], -2^{39}), 2^{39}-1)
\end{lstlisting}

\clearpage
\subsection{EE.VMULAS.S8.ACCX.LD.IP}\label{instruction:EEVMULASS8ACCXLDIP}
\setlength\tabcolsep{6pt}
\textbf{Instruction Word} \\
\begin{tabular}{|c|c|c|c|c|c|c|c|c|c|c|c|c|}
	\hline
	1111 & imm16[5:4] & qu[2:1] & qy[0] & 000 & qu[0] & 010 & qx[1:0] & qy[2:1] & imm16[3:0] & as[3:0] & 111 & qx[2] \\
	\hline
\end{tabular}

\textbf{Assembler Syntax} \\
EE.VMULAS.S8.ACCX.LD.IP qu, as, -512..496, qx, qy \\
\textbf{Description} \\
This instruction divides registers \lstinline{qx} and \lstinline{qy} into 16 data segments by 8 bits. Then, it performs a signed multiply-accumulate operation on the 16 sets of segments respectively. The accumulated result is added to the value in special register \lstinline{ACCX}. Then, the sum obtained is saturated and then stored in \lstinline{ACCX}. \\ During the operation, the lower 4 bits of the access address in register \lstinline{as} are forced to be 0, and then the 16-byte data is loaded from the memory to register \lstinline{qu}. After the access is completed, the value in register \lstinline{as} is incremented by 6-bit sign-extended constant in the instruction code segment left-shifted by 4.\\
\textbf{Operation}
\begin{lstlisting}
  add0[15:0] = qx[  7:  0] * qy[  7:  0]
  add1[15:0] = qx[ 15:  8] * qy[ 15:  8]
  ...
  add15[15:0] = qx[127:120] * qy[127:120]
  sum[40:0] = ACCX[39:0] + ad[31:0]d0[15:0] + ad[31:0]d1[15:0] + ... + ad[31:0]d15[15:0]

  qu[127:0] = load128({as[31:4],4{0}})
  as[31:0] = as[31:0] + {20{imm16[7]},imm16[7:0],4{0}}
\end{lstlisting}

\clearpage
\subsection{EE.VMULAS.S8.ACCX.LD.IP.QUP}\label{instruction:EEVMULASS8ACCXLDIPQUP}
\setlength\tabcolsep{5pt}
\textbf{Instruction Word} \\
\begin{tabular}{|c|c|c|c|c|c|c|c|c|c|c|c|c|}
	\hline
	0010 & imm16[5:4] & qu[2:1] & qy[0] & qs0[2:0] & qu[0] & qs1[2:0] & qx[1:0] & qy[2:1] & imm16[3:0] & as[3:0] & 111 & qx[2] \\
	\hline
\end{tabular}

\textbf{Assembler Syntax} \\
EE.VMULAS.S8.ACCX.LD.IP.QUP qu, as, imm16, qx, qy, qs0, qs1 \\
\textbf{Description} \\
This instruction divides registers \lstinline{qx} and \lstinline{qy} into 16 data segments by 8 bits. Then, it performs a signed multiply-accumulate operation on the 16 sets of segments respectively. The accumulated result is added to the value in special register \lstinline{ACCX}. Then, the sum obtained is saturated and then stored in \lstinline{ACCX}.\\
During the operation, the lower 4 bits of the access address in register \lstinline{as} are forced to be 0, and then the 16-byte data is loaded from the memory to register \lstinline{qu}. After the access is completed, the value in register \lstinline{as} is incremented by 6-bit sign-extended constant in the instruction code segment left-shifted by 4.\\
At the same time, this instruction also obtains 16-byte unaligned data by concatenating and shifting consecutive aligned data stored in the two registers \lstinline{qs0} and \lstinline{qs1} and stores it to \lstinline{qs0}. The shift byte is stored in special register \lstinline{SAR_BYTE}.\\
\textbf{Operation}
\begin{lstlisting}
  add0[15:0] = qx[  7:  0] * qy[  7:  0]
  add1[15:0] = qx[ 15:  8] * qy[ 15:  8]
  ...
  add15[15:0] = qx[127:120] * qy[127:120]
  sum[40:0] = ACCX[39:0] + ad[31:0]d0[15:0] + ad[31:0]d1[15:0] + ... + ad[31:0]d15[15:0]

  qu[127:0] = load128({as[31:4],4{0}})
  as[31:0] = as[31:0] + {20{imm16[7]},imm16[7:0],4{0}}
  qs0[127:0]  = {qs1[127:0], qs0[127:0]} >> {SAR_BYTE[3:0] << 3}
\end{lstlisting}

\clearpage
\subsection{EE.VMULAS.S8.ACCX.LD.XP}\label{instruction:EEVMULASS8ACCXLDXP}
\setlength\tabcolsep{6pt}
\textbf{Instruction Word} \\
\begin{tabular}{|c|c|c|c|c|c|c|c|c|c|c|c|}
	\hline
	111100 & qu[2:1] & qy[0] & 001 & qu[0] & 010 & qx[1:0] & qy[2:1] & ad[3:0] & as[3:0] & 111 & qx[2] \\
	\hline
\end{tabular}

\textbf{Assembler Syntax} \\
EE.VMULAS.S8.ACCX.LD.XP qu, as, ad, qx, qy \\
\textbf{Description} \\
This instruction divides registers \lstinline{qx} and \lstinline{qy} into 16 data segments by 8 bits. Then, it performs a signed multiply-accumulate operation on the 16 sets of segments respectively. The accumulated result is added to the value in special register \lstinline{ACCX}. Then, the sum obtained is saturated and then stored in \lstinline{ACCX}.\\
During the operation, the lower 4 bits of the access address in register \lstinline{as} are forced to be 0, and then the 16-byte data is loaded from the memory to register \lstinline{qu}. After the access is completed, the value in register \lstinline{as} is incremented by the value in register \lstinline{ad}.\\
\textbf{Operation}
\begin{lstlisting}
  add0[15:0] = qx[  7:  0] * qy[  7:  0]
  add1[15:0] = qx[ 15:  8] * qy[ 15:  8]
  ...
  add15[15:0] = qx[127:120] * qy[127:120]
  sum[40:0] = ACCX[39:0] + ad[31:0]d0[15:0] + ad[31:0]d1[15:0] + ... + ad[31:0]d15[15:0]
  ACCX[39:0] = min(max(sum[40:0], -2^{39}), 2^{39}-1)

  qu[127:0] = load128({as[31:4],4{0}})
  as[31:0] = as[31:0] + ad[31:0]
\end{lstlisting}

\clearpage
\subsection{EE.VMULAS.S8.ACCX.LD.XP.QUP}\label{instruction:EEVMULASS8ACCXLDXPQUP}
\setlength\tabcolsep{6pt}
\textbf{Instruction Word} \\
\begin{tabular}{|c|c|c|c|c|c|c|c|c|c|c|c|}
	\hline
	101110 & qu[2:1] & qy[0] & qs0[2:0] & qu[0] & qs1[2:0] & qx[1:0] & qy[2:1] & ad[3:0] & as[3:0] & 111 & qx[2] \\
	\hline
\end{tabular}

\textbf{Assembler Syntax} \\
EE.VMULAS.S8.ACCX.LD.XP.QUP qu, as, ad, qx, qy, qs0, qs1 \\
\textbf{Description} \\
This instruction divides registers \lstinline{qx} and \lstinline{qy} into 16 data segments by 8 bits. Then, it performs a signed multiply-accumulate operation on the 16 sets of segments respectively. The accumulated result is added to the value in special register \lstinline{ACCX}. Then, the sum obtained is saturated and then stored in \lstinline{ACCX}.\\
During the operation, the lower 4 bits of the access address in register \lstinline{as} are forced to be 0, and then the 16-byte data is loaded from the memory to register \lstinline{qu}. After the access is completed, the value in register \lstinline{as} is incremented by the value in register \lstinline{ad}.\\
At the same time, this instruction also obtains 16-byte unaligned data by concatenating and shifting consecutive aligned data stored in the two registers \lstinline{qs0} and \lstinline{qs1} and stores it to \lstinline{qs0}. The shift byte is stored in special register \lstinline{SAR_BYTE}.\\
\textbf{Operation}
\begin{lstlisting}
  add0[15:0] = qx[  7:  0] * qy[  7:  0]
  add1[15:0] = qx[ 15:  8] * qy[ 15:  8]
  ...
  add15[15:0] = qx[127:120] * qy[127:120]
  sum[40:0] = ACCX[39:0] + ad[31:0]d0[15:0] + ad[31:0]d1[15:0] + ... + ad[31:0]d15[15:0]
  ACCX[39:0] = min(max(sum[40:0], -2^{39}), 2^{39}-1)

  qu[127:0] = load128({as[31:4],4{0}})
  as[31:0] = as[31:0] + ad[31:0]
  qs0[127:0] = {qs1[127: 0], qs0[127: 0]} >> {SAR_BYTE[3:0] << 3}
\end{lstlisting}

\clearpage
\subsection{EE.VMULAS.S8.QACC}\label{instruction:EEVMULASS8QACC}
\setlength\tabcolsep{6pt}
\textbf{Instruction Word} \\
\begin{tabular}{|c|c|c|c|c|c|}
	\hline
	000110100 & qy[2] & 1 & qy[1:0] & qx[2:0] & 11000100 \\
	\hline
\end{tabular}

\textbf{Assembler Syntax} \\
EE.VMULAS.S8.QACC qx, qy \\
\textbf{Description} \\
This instruction divides registers \lstinline{qx} and \lstinline{qy} into 16 data segments by 8 bits. Then, the signed multiplication result of the 16 sets of segments is added to the corresponding 20-bit data segment in special registers \lstinline{QACC_H} and \lstinline{QACC_L} respectively. The calculated result is saturated to a 20-bit signed number and then stored to the corresponding 20-bit data segment in \lstinline{QACC_H} and \lstinline{QACC_L}.\\
\textbf{Operation}
\begin{lstlisting}
  QACC_L[ 19:  0] = min(max(QACC_L[ 19:  0] + qx[  7:  0] * qy[  7:  0], -2^{19}), 2^{19}-1)
  QACC_L[ 39: 20] = min(max(QACC_L[ 39: 20] + qx[ 15:  8] * qy[ 15:  8], -2^{19}), 2^{19}-1)
  QACC_L[ 59: 40] = min(max(QACC_L[ 59: 40] + qx[ 23: 16] * qy[ 23: 16], -2^{19}), 2^{19}-1)
  QACC_L[ 79: 60] = min(max(QACC_L[ 79: 60] + qx[ 31: 24] * qy[ 31: 24], -2^{19}), 2^{19}-1)
  QACC_L[ 99: 80] = min(max(QACC_L[ 99: 80] + qx[ 39: 32] * qy[ 39: 32], -2^{19}), 2^{19}-1)
  QACC_L[119:100] = min(max(QACC_L[119:100] + qx[ 47: 40] * qy[ 47: 40], -2^{19}), 2^{19}-1)
  QACC_L[139:120] = min(max(QACC_L[139:120] + qx[ 55: 48] * qy[ 55: 48], -2^{19}), 2^{19}-1)
  QACC_L[159:140] = min(max(QACC_L[159:140] + qx[ 63: 56] * qy[ 63: 56], -2^{19}), 2^{19}-1)
  QACC_H[ 19:  0] = min(max(QACC_H[ 19:  0] + qx[ 71: 64] * qy[ 71: 64], -2^{19}), 2^{19}-1)
  QACC_H[ 39: 20] = min(max(QACC_H[ 39: 20] + qx[ 79: 72] * qy[ 79: 72], -2^{19}), 2^{19}-1)
  QACC_H[ 59: 40] = min(max(QACC_H[ 59: 40] + qx[ 87: 80] * qy[ 87: 80], -2^{19}), 2^{19}-1)
  QACC_H[ 79: 60] = min(max(QACC_H[ 79: 60] + qx[ 95: 88] * qy[ 95: 88], -2^{19}), 2^{19}-1)
  QACC_H[ 99: 80] = min(max(QACC_H[ 99: 80] + qx[103: 96] * qy[103: 96], -2^{19}), 2^{19}-1)
  QACC_H[119:100] = min(max(QACC_H[119:100] + qx[111:104] * qy[111:104], -2^{19}), 2^{19}-1)
  QACC_H[139:120] = min(max(QACC_H[139:120] + qx[119:112] * qy[119:112], -2^{19}), 2^{19}-1)
  QACC_H[159:140] = min(max(QACC_H[159:140] + qx[127:120] * qy[127:120], -2^{19}), 2^{19}-1)
\end{lstlisting}

\clearpage
\subsection{EE.VMULAS.S8.QACC.LD.IP}\label{instruction:EEVMULASS8QACCLDIP}
\setlength\tabcolsep{6pt}
\textbf{Instruction Word} \\
\begin{tabular}{|c|c|c|c|c|c|c|c|c|c|c|c|c|}
	\hline
	1111 & imm16[5:4] & qu[2:1] & qy[0] & 000 & qu[0] & 011 & qx[1:0] & qy[2:1] & imm16[3:0] & as[3:0] & 111 & qx[2] \\
	\hline
\end{tabular}
This instruction divides registers \lstinline{qx} and \lstinline{qy} into 16 data segments by 8 bits. Then, the signed multiplication result of the 16 sets of segments is added to the corresponding 20-bit data segment in special registers \lstinline{QACC_H} and \lstinline{QACC_L} respectively. The calculated result is saturated to a 20-bit signed number and then stored to the corresponding 20-bit data segment in \lstinline{QACC_H} and \lstinline{QACC_L}.\\
During the operation, the lower 4 bits of the access address in register \lstinline{as} are forced to be 0, and then the 16-byte data is loaded from the memory to register \lstinline{qu}. After the access is completed, the value in register \lstinline{as} is incremented by 6-bit sign-extended constant in the instruction code segment left-shifted by 4.\\
\textbf{Assembler Syntax} \\
EE.VMULAS.S8.QACC.LD.IP qu, as, imm16, qx, qy \\
\textbf{Description} \\
\textbf{Operation}
\begin{lstlisting}
  QACC_L[ 19:  0] = min(max(QACC_L[ 19:  0] + qx[  7:  0] * qy[  7:  0], -2^{19}), 2^{19}-1)
  QACC_L[ 39: 20] = min(max(QACC_L[ 39: 20] + qx[ 15:  8] * qy[ 15:  8], -2^{19}), 2^{19}-1)
  QACC_L[ 59: 40] = min(max(QACC_L[ 59: 40] + qx[ 23: 16] * qy[ 23: 16], -2^{19}), 2^{19}-1)
  QACC_L[ 79: 60] = min(max(QACC_L[ 79: 60] + qx[ 31: 24] * qy[ 31: 24], -2^{19}), 2^{19}-1)
  QACC_L[ 99: 80] = min(max(QACC_L[ 99: 80] + qx[ 39: 32] * qy[ 39: 32], -2^{19}), 2^{19}-1)
  QACC_L[119:100] = min(max(QACC_L[119:100] + qx[ 47: 40] * qy[ 47: 40], -2^{19}), 2^{19}-1)
  QACC_L[139:120] = min(max(QACC_L[139:120] + qx[ 55: 48] * qy[ 55: 48], -2^{19}), 2^{19}-1)
  QACC_L[159:140] = min(max(QACC_L[159:140] + qx[ 63: 56] * qy[ 63: 56], -2^{19}), 2^{19}-1)
  QACC_H[ 19:  0] = min(max(QACC_H[ 19:  0] + qx[ 71: 64] * qy[ 71: 64], -2^{19}), 2^{19}-1)
  QACC_H[ 39: 20] = min(max(QACC_H[ 39: 20] + qx[ 79: 72] * qy[ 79: 72], -2^{19}), 2^{19}-1)
  QACC_H[ 59: 40] = min(max(QACC_H[ 59: 40] + qx[ 87: 80] * qy[ 87: 80], -2^{19}), 2^{19}-1)
  QACC_H[ 79: 60] = min(max(QACC_H[ 79: 60] + qx[ 95: 88] * qy[ 95: 88], -2^{19}), 2^{19}-1)
  QACC_H[ 99: 80] = min(max(QACC_H[ 99: 80] + qx[103: 96] * qy[103: 96], -2^{19}), 2^{19}-1)
  QACC_H[119:100] = min(max(QACC_H[119:100] + qx[111:104] * qy[111:104], -2^{19}), 2^{19}-1)
  QACC_H[139:120] = min(max(QACC_H[139:120] + qx[119:112] * qy[119:112], -2^{19}), 2^{19}-1)
  QACC_H[159:140] = min(max(QACC_H[159:140] + qx[127:120] * qy[127:120], -2^{19}), 2^{19}-1)

  qu[127:0] = load128({as[31:4],4{0}})
  as[31:0] = as[31:0] + {22{imm16[5]},imm16[5:0],4{0}}
\end{lstlisting}

\clearpage
\subsection{EE.VMULAS.S8.QACC.LD.IP.QUP}\label{instruction:EEVMULASS8QACCLDIPQUP}
\setlength\tabcolsep{5pt}
\textbf{Instruction Word} \\
\begin{tabular}{|c|c|c|c|c|c|c|c|c|c|c|c|c|}
	\hline
	0011 & imm16[5:4] & qu[2:1] & qy[0] & qs0[2:0] & qu[0] & qs1[2:0] & qx[1:0] & qy[2:1] & imm16[3:0] & as[3:0] & 111 & qx[2] \\
	\hline
\end{tabular}

\textbf{Assembler Syntax} \\
EE.VMULAS.S8.QACC.LD.IP.QUP qu, as, imm16, qx, qy, qs0, qs1 \\
\textbf{Description} \\
This instruction divides registers \lstinline{qx} and \lstinline{qy} into 16 data segments by 8 bits. Then, the signed multiplication result of the 16 sets of segments is added to the corresponding 20-bit data segment in special registers \lstinline{QACC_H} and \lstinline{QACC_L} respectively. The calculated result is saturated to a 20-bit signed number and then stored to the corresponding 20-bit data segment in \lstinline{QACC_H} and \lstinline{QACC_L}.\\
During the operation, the lower 4 bits of the access address in register \lstinline{as} are forced to be 0, and then the 16-byte data is loaded from the memory to register \lstinline{qu}. After the access is completed, the value in register \lstinline{as} is incremented by 6-bit sign-extended constant in the instruction code segment left-shifted by 4.\\
At the same time, this instruction also obtains 16-byte unaligned data by concatenating and shifting the two registers \lstinline{qs0} and \lstinline{qs1} that store consecutive aligned data and stores it to \lstinline{qs0}. The shift byte is stored in special register \lstinline{SAR_BYTE}.\\
\textbf{Operation}
\begin{lstlisting}
  QACC_L[ 19:  0] = min(max(QACC_L[ 19:  0] + qx[  7:  0] * qy[  7:  0], -2^{19}), 2^{19}-1)
  QACC_L[ 39: 20] = min(max(QACC_L[ 39: 20] + qx[ 15:  8] * qy[ 15:  8], -2^{19}), 2^{19}-1)
  QACC_L[ 59: 40] = min(max(QACC_L[ 59: 40] + qx[ 23: 16] * qy[ 23: 16], -2^{19}), 2^{19}-1)
  QACC_L[ 79: 60] = min(max(QACC_L[ 79: 60] + qx[ 31: 24] * qy[ 31: 24], -2^{19}), 2^{19}-1)
  QACC_L[ 99: 80] = min(max(QACC_L[ 99: 80] + qx[ 39: 32] * qy[ 39: 32], -2^{19}), 2^{19}-1)
  QACC_L[119:100] = min(max(QACC_L[119:100] + qx[ 47: 40] * qy[ 47: 40], -2^{19}), 2^{19}-1)
  QACC_L[139:120] = min(max(QACC_L[139:120] + qx[ 55: 48] * qy[ 55: 48], -2^{19}), 2^{19}-1)
  QACC_L[159:140] = min(max(QACC_L[159:140] + qx[ 63: 56] * qy[ 63: 56], -2^{19}), 2^{19}-1)
  QACC_H[ 19:  0] = min(max(QACC_H[ 19:  0] + qx[ 71: 64] * qy[ 71: 64], -2^{19}), 2^{19}-1)
  QACC_H[ 39: 20] = min(max(QACC_H[ 39: 20] + qx[ 79: 72] * qy[ 79: 72], -2^{19}), 2^{19}-1)
  QACC_H[ 59: 40] = min(max(QACC_H[ 59: 40] + qx[ 87: 80] * qy[ 87: 80], -2^{19}), 2^{19}-1)
  QACC_H[ 79: 60] = min(max(QACC_H[ 79: 60] + qx[ 95: 88] * qy[ 95: 88], -2^{19}), 2^{19}-1)
  QACC_H[ 99: 80] = min(max(QACC_H[ 99: 80] + qx[103: 96] * qy[103: 96], -2^{19}), 2^{19}-1)
  QACC_H[119:100] = min(max(QACC_H[119:100] + qx[111:104] * qy[111:104], -2^{19}), 2^{19}-1)
  QACC_H[139:120] = min(max(QACC_H[139:120] + qx[119:112] * qy[119:112], -2^{19}), 2^{19}-1)
  QACC_H[159:140] = min(max(QACC_H[159:140] + qx[127:120] * qy[127:120], -2^{19}), 2^{19}-1)

  qu[127:0] = load128({as[31:4],4{0}})
  as[31:0] = as[31:0] + {22{imm16[5]},imm16[5:0],4{0}}
  qs0[127:0]  = {qs1[127:0], qs0[127:0]} >> {SAR_BYTE[3:0] << 3}
\end{lstlisting}

\clearpage
\subsection{EE.VMULAS.S8.QACC.LD.XP}\label{instruction:EEVMULASS8QACCLDXP}
\setlength\tabcolsep{6pt}
\textbf{Instruction Word} \\
\begin{tabular}{|c|c|c|c|c|c|c|c|c|c|c|c|}
	\hline
	111100 & qu[2:1] & qy[0] & 001 & qu[0] & 011 & qx[1:0] & qy[2:1] & ad[3:0] & as[3:0] & 111 & qx[2] \\
	\hline
\end{tabular}

\textbf{Assembler Syntax} \\
EE.VMULAS.S8.QACC.LD.XP qu, as, ad, qx, qy \\
\textbf{Description} \\
This instruction divides registers \lstinline{qx} and \lstinline{qy} into 16 data segments by 8 bits. Then, the signed multiplication result of the 16 sets of segments is added to the corresponding 20-bit data segment in special registers \lstinline{QACC_H} and \lstinline{QACC_L} respectively. The calculated result is saturated to a 20-bit signed number and then stored to the corresponding 20-bit data segment in \lstinline{QACC_H} and \lstinline{QACC_L}.\\
During the operation, the lower 4 bits of the access address in register \lstinline{as} are forced to be 0, and then the 16-byte data is loaded from the memory to register \lstinline{qu}. After the access is completed, the value in register \lstinline{as} is incremented by the value in register \lstinline{ad}.\\
\textbf{Operation}
\begin{lstlisting}
  QACC_L[ 19:  0] = min(max(QACC_L[ 19:  0] + qx[  7:  0] * qy[  7:  0], -2^{19}), 2^{19}-1)
  QACC_L[ 39: 20] = min(max(QACC_L[ 39: 20] + qx[ 15:  8] * qy[ 15:  8], -2^{19}), 2^{19}-1)
  QACC_L[ 59: 40] = min(max(QACC_L[ 59: 40] + qx[ 23: 16] * qy[ 23: 16], -2^{19}), 2^{19}-1)
  QACC_L[ 79: 60] = min(max(QACC_L[ 79: 60] + qx[ 31: 24] * qy[ 31: 24], -2^{19}), 2^{19}-1)
  QACC_L[ 99: 80] = min(max(QACC_L[ 99: 80] + qx[ 39: 32] * qy[ 39: 32], -2^{19}), 2^{19}-1)
  QACC_L[119:100] = min(max(QACC_L[119:100] + qx[ 47: 40] * qy[ 47: 40], -2^{19}), 2^{19}-1)
  QACC_L[139:120] = min(max(QACC_L[139:120] + qx[ 55: 48] * qy[ 55: 48], -2^{19}), 2^{19}-1)
  QACC_L[159:140] = min(max(QACC_L[159:140] + qx[ 63: 56] * qy[ 63: 56], -2^{19}), 2^{19}-1)
  QACC_H[ 19:  0] = min(max(QACC_H[ 19:  0] + qx[ 71: 64] * qy[ 71: 64], -2^{19}), 2^{19}-1)
  QACC_H[ 39: 20] = min(max(QACC_H[ 39: 20] + qx[ 79: 72] * qy[ 79: 72], -2^{19}), 2^{19}-1)
  QACC_H[ 59: 40] = min(max(QACC_H[ 59: 40] + qx[ 87: 80] * qy[ 87: 80], -2^{19}), 2^{19}-1)
  QACC_H[ 79: 60] = min(max(QACC_H[ 79: 60] + qx[ 95: 88] * qy[ 95: 88], -2^{19}), 2^{19}-1)
  QACC_H[ 99: 80] = min(max(QACC_H[ 99: 80] + qx[103: 96] * qy[103: 96], -2^{19}), 2^{19}-1)
  QACC_H[119:100] = min(max(QACC_H[119:100] + qx[111:104] * qy[111:104], -2^{19}), 2^{19}-1)
  QACC_H[139:120] = min(max(QACC_H[139:120] + qx[119:112] * qy[119:112], -2^{19}), 2^{19}-1)
  QACC_H[159:140] = min(max(QACC_H[159:140] + qx[127:120] * qy[127:120], -2^{19}), 2^{19}-1)

  qu[127:0] = load128({as[31:4],4{0}})
  as[31:0] = as[31:0] + ad[31:0]
\end{lstlisting}

\clearpage
\subsection{EE.VMULAS.S8.QACC.LD.XP.QUP}\label{instruction:EEVMULASS8QACCLDXPQUP}
\setlength\tabcolsep{6pt}
\textbf{Instruction Word} \\
\begin{tabular}{|c|c|c|c|c|c|c|c|c|c|c|c|}
	\hline
	101111 & qu[2:1] & qy[0] & qs0[2:0] & qu[0] & qs1[2:0] & qx[1:0] & qy[2:1] & ad[3:0] & as[3:0] & 111 & qx[2] \\
	\hline
\end{tabular}

\textbf{Assembler Syntax} \\
EE.VMULAS.S8.QACC.LD.XP.QUP qu, as, ad, qx, qy, qs0, qs1 \\
\textbf{Description} \\
This instruction divides registers \lstinline{qx} and \lstinline{qy} into 16 data segments by 8 bits. Then, the signed multiplication result of the 16 sets of segments is added to the corresponding 20-bit data segment in special registers \lstinline{QACC_H} and \lstinline{QACC_L} respectively. The calculated result is saturated to a 20-bit signed number and then stored to the corresponding 20-bit data segment in \lstinline{QACC_H} and \lstinline{QACC_L}.\\
During the operation, the lower 4 bits of the access address in register \lstinline{as} are forced to be 0, and then the 16-byte data is loaded from the memory to register \lstinline{qu}. After the access is completed, the value in register \lstinline{as} is incremented by the value in register \lstinline{ad}.\\
At the same time, this instruction also obtains 16-byte unaligned data by concatenating and shifting consecutive aligned data stored in the two registers \lstinline{qs0} and \lstinline{qs1} and stores it to \lstinline{qs0}. The shift byte is stored in special register \lstinline{SAR_BYTE}.\\
\textbf{Operation}
\begin{lstlisting}
  QACC_L[ 19:  0] = min(max(QACC_L[ 19:  0] + qx[  7:  0] * qy[  7:  0], -2^{19}), 2^{19}-1)
  QACC_L[ 39: 20] = min(max(QACC_L[ 39: 20] + qx[ 15:  8] * qy[ 15:  8], -2^{19}), 2^{19}-1)
  QACC_L[ 59: 40] = min(max(QACC_L[ 59: 40] + qx[ 23: 16] * qy[ 23: 16], -2^{19}), 2^{19}-1)
  QACC_L[ 79: 60] = min(max(QACC_L[ 79: 60] + qx[ 31: 24] * qy[ 31: 24], -2^{19}), 2^{19}-1)
  QACC_L[ 99: 80] = min(max(QACC_L[ 99: 80] + qx[ 39: 32] * qy[ 39: 32], -2^{19}), 2^{19}-1)
  QACC_L[119:100] = min(max(QACC_L[119:100] + qx[ 47: 40] * qy[ 47: 40], -2^{19}), 2^{19}-1)
  QACC_L[139:120] = min(max(QACC_L[139:120] + qx[ 55: 48] * qy[ 55: 48], -2^{19}), 2^{19}-1)
  QACC_L[159:140] = min(max(QACC_L[159:140] + qx[ 63: 56] * qy[ 63: 56], -2^{19}), 2^{19}-1)
  QACC_H[ 19:  0] = min(max(QACC_H[ 19:  0] + qx[ 71: 64] * qy[ 71: 64], -2^{19}), 2^{19}-1)
  QACC_H[ 39: 20] = min(max(QACC_H[ 39: 20] + qx[ 79: 72] * qy[ 79: 72], -2^{19}), 2^{19}-1)
  QACC_H[ 59: 40] = min(max(QACC_H[ 59: 40] + qx[ 87: 80] * qy[ 87: 80], -2^{19}), 2^{19}-1)
  QACC_H[ 79: 60] = min(max(QACC_H[ 79: 60] + qx[ 95: 88] * qy[ 95: 88], -2^{19}), 2^{19}-1)
  QACC_H[ 99: 80] = min(max(QACC_H[ 99: 80] + qx[103: 96] * qy[103: 96], -2^{19}), 2^{19}-1)
  QACC_H[119:100] = min(max(QACC_H[119:100] + qx[111:104] * qy[111:104], -2^{19}), 2^{19}-1)
  QACC_H[139:120] = min(max(QACC_H[139:120] + qx[119:112] * qy[119:112], -2^{19}), 2^{19}-1)
  QACC_H[159:140] = min(max(QACC_H[159:140] + qx[127:120] * qy[127:120], -2^{19}), 2^{19}-1)

  qu[127:0] = load128({as[31:4],4{0}})
  as[31:0] = as[31:0] + ad[31:0]
  qs0[127:0]  = {qs1[127:0], qs0[127:0]} >> {SAR_BYTE[3:0] << 3}
\end{lstlisting}

\clearpage
\subsection{EE.VMULAS.S8.QACC.LDBC.INCP}\label{instruction:EEVMULASS8QACCLDBCINCP}
\setlength\tabcolsep{6pt}
\textbf{Instruction Word} \\
\begin{tabular}{|c|c|c|c|c|c|c|c|c|c|}
	\hline
	101 & qu[2] & 0111 & qu[1] & qy[2] & qu[0] & qy[1:0] & qx[2:0] & as[3:0] & 0100 \\
	\hline
\end{tabular}

\textbf{Assembler Syntax} \\
EE.VMULAS.S8.QACC.LDBC.INCP qu, as, qx, qy \\
\textbf{Description} \\
This instruction divides registers \lstinline{qx} and \lstinline{qy} into 16 data segments by 8 bits. Then, the signed multiplication result of the 16 sets of segments is added to the corresponding 20-bit data segment in special registers \lstinline{QACC_H} and \lstinline{QACC_L} respectively. The calculated result is saturated to a 20-bit signed number and then stored to the corresponding 20-bit data segment in \lstinline{QACC_H} and \lstinline{QACC_L}.\\
At the same time, this instruction loads 8-bit data from memory at the address given by the access register \lstinline{as} and broadcasts it to the 16 8-bit data segments in register \lstinline{qu}. After the access, the value in register \lstinline{as} is incremented by 1.\\
\textbf{Operation}
\begin{lstlisting}
  QACC_L[ 19:  0] = min(max(QACC_L[ 19:  0] + qx[  7:  0] * qy[  7:  0], -2^{19}), 2^{19}-1)
  QACC_L[ 39: 20] = min(max(QACC_L[ 39: 20] + qx[ 15:  8] * qy[ 15:  8], -2^{19}), 2^{19}-1)
  QACC_L[ 59: 40] = min(max(QACC_L[ 59: 40] + qx[ 23: 16] * qy[ 23: 16], -2^{19}), 2^{19}-1)
  QACC_L[ 79: 60] = min(max(QACC_L[ 79: 60] + qx[ 31: 24] * qy[ 31: 24], -2^{19}), 2^{19}-1)
  QACC_L[ 99: 80] = min(max(QACC_L[ 99: 80] + qx[ 39: 32] * qy[ 39: 32], -2^{19}), 2^{19}-1)
  QACC_L[119:100] = min(max(QACC_L[119:100] + qx[ 47: 40] * qy[ 47: 40], -2^{19}), 2^{19}-1)
  QACC_L[139:120] = min(max(QACC_L[139:120] + qx[ 55: 48] * qy[ 55: 48], -2^{19}), 2^{19}-1)
  QACC_L[159:140] = min(max(QACC_L[159:140] + qx[ 63: 56] * qy[ 63: 56], -2^{19}), 2^{19}-1)
  QACC_H[ 19:  0] = min(max(QACC_H[ 19:  0] + qx[ 71: 64] * qy[ 71: 64], -2^{19}), 2^{19}-1)
  QACC_H[ 39: 20] = min(max(QACC_H[ 39: 20] + qx[ 79: 72] * qy[ 79: 72], -2^{19}), 2^{19}-1)
  QACC_H[ 59: 40] = min(max(QACC_H[ 59: 40] + qx[ 87: 80] * qy[ 87: 80], -2^{19}), 2^{19}-1)
  QACC_H[ 79: 60] = min(max(QACC_H[ 79: 60] + qx[ 95: 88] * qy[ 95: 88], -2^{19}), 2^{19}-1)
  QACC_H[ 99: 80] = min(max(QACC_H[ 99: 80] + qx[103: 96] * qy[103: 96], -2^{19}), 2^{19}-1)
  QACC_H[119:100] = min(max(QACC_H[119:100] + qx[111:104] * qy[111:104], -2^{19}), 2^{19}-1)
  QACC_H[139:120] = min(max(QACC_H[139:120] + qx[119:112] * qy[119:112], -2^{19}), 2^{19}-1)
  QACC_H[159:140] = min(max(QACC_H[159:140] + qx[127:120] * qy[127:120], -2^{19}), 2^{19}-1)

  qu[127:0] = {16{load8(as[31:0])}}
  as[31:0] = as[31:0] + 1
\end{lstlisting}

\clearpage
\subsection{EE.VMULAS.S8.QACC.LDBC.INCP.QUP}\label{instruction:EEVMULASS8QACCLDBCINCPQUP}
\setlength\tabcolsep{6pt}
\textbf{Instruction Word} \\
\begin{tabular}{|c|c|c|c|c|c|c|c|c|c|c|c|}
	\hline
	111000 & qu[2:1] & qy[0] & qs0[2:0] & qu[0] & qs1[2:0] & qx[1:0] & qy[2:1] & 1001 & as[3:0] & 111 & qx[2] \\
	\hline
\end{tabular}

\textbf{Assembler Syntax} \\
EE.VMULAS.S8.QACC.LDBC.INCP.QUP qu, as, qx, qy, qs0, qs1 \\
\textbf{Description} \\
This instruction divides registers \lstinline{qx} and \lstinline{qy} into 16 data segments by 8 bits. Then, the signed multiplication result of the 16 sets of segments is added to the corresponding 20-bit data segment in special registers \lstinline{QACC_H} and \lstinline{QACC_L} respectively. The calculated result is saturated to a 20-bit signed number and then stored to the corresponding 20-bit data segment in \lstinline{QACC_H} and \lstinline{QACC_L}.\\
At the same time, this instruction loads 8-bit data from memory at the address given by the access register \lstinline{as} and broadcasts it to the 16 8-bit data segments in register \lstinline{qu}. After the access, the value in register \lstinline{as} is incremented by 1.\\
At the same time, this instruction also obtains 16-byte unaligned data by concatenating and shifting consecutive aligned data stored in the two registers \lstinline{qs0} and \lstinline{qs1} and stores it to \lstinline{qs0}. The shift byte is stored in special register \lstinline{SAR_BYTE}.\\
\textbf{Operation}
\begin{lstlisting}
  QACC_L[ 19:  0] = min(max(QACC_L[ 19:  0] + qx[  7:  0] * qy[  7:  0], -2^{19}), 2^{19}-1)
  QACC_L[ 39: 20] = min(max(QACC_L[ 39: 20] + qx[ 15:  8] * qy[ 15:  8], -2^{19}), 2^{19}-1)
  QACC_L[ 59: 40] = min(max(QACC_L[ 59: 40] + qx[ 23: 16] * qy[ 23: 16], -2^{19}), 2^{19}-1)
  QACC_L[ 79: 60] = min(max(QACC_L[ 79: 60] + qx[ 31: 24] * qy[ 31: 24], -2^{19}), 2^{19}-1)
  QACC_L[ 99: 80] = min(max(QACC_L[ 99: 80] + qx[ 39: 32] * qy[ 39: 32], -2^{19}), 2^{19}-1)
  QACC_L[119:100] = min(max(QACC_L[119:100] + qx[ 47: 40] * qy[ 47: 40], -2^{19}), 2^{19}-1)
  QACC_L[139:120] = min(max(QACC_L[139:120] + qx[ 55: 48] * qy[ 55: 48], -2^{19}), 2^{19}-1)
  QACC_L[159:140] = min(max(QACC_L[159:140] + qx[ 63: 56] * qy[ 63: 56], -2^{19}), 2^{19}-1)
  QACC_H[ 19:  0] = min(max(QACC_H[ 19:  0] + qx[ 71: 64] * qy[ 71: 64], -2^{19}), 2^{19}-1)
  QACC_H[ 39: 20] = min(max(QACC_H[ 39: 20] + qx[ 79: 72] * qy[ 79: 72], -2^{19}), 2^{19}-1)
  QACC_H[ 59: 40] = min(max(QACC_H[ 59: 40] + qx[ 87: 80] * qy[ 87: 80], -2^{19}), 2^{19}-1)
  QACC_H[ 79: 60] = min(max(QACC_H[ 79: 60] + qx[ 95: 88] * qy[ 95: 88], -2^{19}), 2^{19}-1)
  QACC_H[ 99: 80] = min(max(QACC_H[ 99: 80] + qx[103: 96] * qy[103: 96], -2^{19}), 2^{19}-1)
  QACC_H[119:100] = min(max(QACC_H[119:100] + qx[111:104] * qy[111:104], -2^{19}), 2^{19}-1)
  QACC_H[139:120] = min(max(QACC_H[139:120] + qx[119:112] * qy[119:112], -2^{19}), 2^{19}-1)
  QACC_H[159:140] = min(max(QACC_H[159:140] + qx[127:120] * qy[127:120], -2^{19}), 2^{19}-1)

  qu[127:0] = {16{load8(as[31:0])}}
  as[31:0] = as[31:0] + 1
  qs0[127:0]  = {qs1[127:0], qs0[127:0]} >> {SAR_BYTE[3:0] << 3}
\end{lstlisting}

\clearpage
\subsection{EE.VMULAS.U16.ACCX}\label{instruction:EEVMULASU16ACCX}
\setlength\tabcolsep{6pt}
\textbf{Instruction Word} \\
\begin{tabular}{|c|c|c|c|c|c|}
	\hline
	000010100 & qy[2] & 0 & qy[1:0] & qx[2:0] & 10000100 \\
	\hline
\end{tabular}

\textbf{Assembler Syntax} \\
EE.VMULAS.U16.ACCX qx, qy \\
\textbf{Description} \\
This instruction divides registers \lstinline{qx} and \lstinline{qy} into 8 data segments by 16 bits. Then, it performs an unsigned multiply-accumulate operation on the 8 sets of segments respectively. The accumulated result is added to the value in special register \lstinline{ACCX}. Then, the sum obtained is saturated and then stored in \lstinline{ACCX}.\\
\textbf{Operation}
\begin{lstlisting}
  add0[31:0] = qx[ 15:  0] * qy[ 15:  0]
  add1[31:0] = qx[ 31: 16] * qy[ 31: 16]
  ...
  add7[31:0] = qx[127:112] * qy[127:112]
  sum[40:0] = ACCX[39:0] + ad[31:0]d0[31:0] + ad[31:0]d1[31:0] + ... + ad[31:0]d7[31:0]

  ACCX[39:0] = min(max(sum[40:0], 0), 2^{40}-1)
\end{lstlisting}

\clearpage
\subsection{EE.VMULAS.U16.ACCX.LD.IP}\label{instruction:EEVMULASU16ACCXLDIP}
\setlength\tabcolsep{6pt}
\textbf{Instruction Word} \\
\begin{tabular}{|c|c|c|c|c|c|c|c|c|c|c|c|c|}
	\hline
	1111 & imm16[5:4] & qu[2:1] & qy[0] & 000 & qu[0] & 100 & qx[1:0] & qy[2:1] & imm16[3:0] & as[3:0] & 111 & qx[2] \\
	\hline
\end{tabular}

\textbf{Assembler Syntax} \\
EE.VMULAS.U16.ACCX.LD.IP qu, as, -512..496, qx, qy \\
\textbf{Description} \\
This instruction divides registers \lstinline{qx} and \lstinline{qy} into 8 data segments by 16 bits. Then, it performs an unsigned multiply-accumulate operation on the 8 sets of segments respectively. The accumulated result is added to the value in special register \lstinline{ACCX}. Then, the sum obtained is saturated and then stored in \lstinline{ACCX}. \\ During the operation, the lower 4 bits of the access address in register \lstinline{as} are forced to be 0, and then the 16-byte data is loaded from the memory to register \lstinline{qu}. After the access is completed, the value in register \lstinline{as} is incremented by 6-bit sign-extended constant in the instruction code segment left-shifted by 4.\\
\textbf{Operation}
\begin{lstlisting}
  add0[31:0] = qx[ 15:  0] * qy[ 15:  0]
  add1[31:0] = qx[ 31: 16] * qy[ 31: 16]
  ...
  add7[31:0] = qx[127:112] * qy[127:112]
  sum[40:0] = ACCX[39:0] + ad[31:0]d0[31:0] + ad[31:0]d1[31:0] + ... + ad[31:0]d7[31:0]
  ACCX[40:0] = min(max(sum[40:0], 0), 2^{40}-1)

  qu[127:0] = load128({as[31:4],4{0}})
  as[31:0] = as[31:0] + {20{imm16[7]},imm16[7:0],4{0}}
\end{lstlisting}

\clearpage
\subsection{EE.VMULAS.U16.ACCX.LD.IP.QUP}\label{instruction:EEVMULASU16ACCXLDIPQUP}
\setlength\tabcolsep{5pt}
\textbf{Instruction Word} \\
\begin{tabular}{|c|c|c|c|c|c|c|c|c|c|c|c|c|}
	\hline
	0100 & imm16[5:4] & qu[2:1] & qy[0] & qs0[2:0] & qu[0] & qs1[2:0] & qx[1:0] & qy[2:1] & imm16[3:0] & as[3:0] & 111 & qx[2] \\
	\hline
\end{tabular}

\textbf{Assembler Syntax} \\
EE.VMULAS.U16.ACCX.LD.IP.QUP qu, as, imm16, qx, qy, qs0, qs1 \\
\textbf{Description} \\
This instruction divides registers \lstinline{qx} and \lstinline{qy} into 8 data segments by 16 bits. Then, it performs a signed multiply-accumulate operation on the 8 sets of segments respectively. The accumulated result is added to the value in special register \lstinline{ACCX}. Then, the sum obtained is saturated and then stored in \lstinline{ACCX}.\\
During the operation, the lower 4 bits of the access address in register \lstinline{as} are forced to be 0, and then the 16-byte data is loaded from the memory to register \lstinline{qu}. After the access is completed, the value in register \lstinline{as} is incremented by 6-bit sign-extended constant in the instruction code segment left-shifted by 4.\\
At the same time, this instruction also obtains 16-byte unaligned data by concatenating and shifting consecutive aligned data stored in the two registers \lstinline{qs0} and \lstinline{qs1} and stores it to \lstinline{qs0}. The shift byte is stored in special register \lstinline{SAR_BYTE}.\\
\textbf{Operation}
\begin{lstlisting}
  add0[31:0] = qx[ 15:  0] * qy[ 15:  0]
  add1[31:0] = qx[ 31: 16] * qy[ 31: 16]
  ...
  add7[31:0] = qx[127:112] * qy[127:112]
  sum[40:0] = ACCX[39:0] + ad[31:0]d0[31:0] + ad[31:0]d1[31:0] + ... + ad[31:0]d7[31:0]
  ACCX[40:0] = min(max(sum[40:0], 0), 2^{40}-1)

  qu[127:0] = load128({as[31:4],4{0}})
  as[31:0] = as[31:0] + {20{imm16[7]},imm16[7:0],4{0}}
  qs0[127:0]  = {qs1[127:0], qs0[127:0]} >> {SAR_BYTE[3:0] << 3}
\end{lstlisting}

\clearpage
\subsection{EE.VMULAS.U16.ACCX.LD.XP}\label{instruction:EEVMULASU16ACCXLDXP}
\setlength\tabcolsep{6pt}
\textbf{Instruction Word} \\
\begin{tabular}{|c|c|c|c|c|c|c|c|c|c|c|c|}
	\hline
	111100 & qu[2:1] & qy[0] & 001 & qu[0] & 100 & qx[1:0] & qy[2:1] & ad[3:0] & as[3:0] & 111 & qx[2] \\
	\hline
\end{tabular}

\textbf{Assembler Syntax} \\
EE.VMULAS.U16.ACCX.LD.XP qu, as, ad, qx, qy \\
\textbf{Description} \\
This instruction divides registers \lstinline{qx} and \lstinline{qy} into 8 data segments by 16 bits. Then, it performs an unsigned multiply-accumulate operation on the 8 sets of segments respectively. The accumulated result is added to the value in special register \lstinline{ACCX}. Then, the sum obtained is saturated and then stored in \lstinline{ACCX}.\\
During the operation, the lower 4 bits of the access address in register \lstinline{as} are forced to be 0, and then the 16-byte data is loaded from the memory to register \lstinline{qu}. After the access is completed, the value in register \lstinline{as} is incremented by the value in register \lstinline{ad}.\\
\textbf{Operation}
\begin{lstlisting}
  add0[31:0] = qx[ 15:  0] * qy[ 15:  0]
  add1[31:0] = qx[ 31: 16] * qy[ 31: 16]
  ...
  add7[31:0] = qx[127:112] * qy[127:112]
  sum[40:0] = ACCX[39:0] + ad[31:0]d0[31:0] + ad[31:0]d1[31:0] + ... + ad[31:0]d7[31:0]
  ACCX[40:0] = min(max(sum[40:0], 0), 2^{40}-1)

  qu[127:0] = load128({as[31:4],4{0}})
  as[31:0] = as[31:0] + ad[31:0]
\end{lstlisting}

\clearpage
\subsection{EE.VMULAS.U16.ACCX.LD.XP.QUP}\label{instruction:EEVMULASU16ACCXLDXPQUP}
\setlength\tabcolsep{6pt}
\textbf{Instruction Word} \\
\begin{tabular}{|c|c|c|c|c|c|c|c|c|c|c|c|}
	\hline
	110000 & qu[2:1] & qy[0] & qs0[2:0] & qu[0] & qs1[2:0] & qx[1:0] & qy[2:1] & ad[3:0] & as[3:0] & 111 & qx[2] \\
	\hline
\end{tabular}

\textbf{Assembler Syntax} \\
EE.VMULAS.U16.ACCX.LD.XP.QUP qu, as, ad, qx, qy, qs0, qs1 \\
\textbf{Description} \\
This instruction divides registers \lstinline{qx} and \lstinline{qy} into 8 data segments by 16 bits. Then, it performs an unsigned multiply-accumulate operation on the 8 sets of segments respectively. The accumulated result is added to the value in special register \lstinline{ACCX}. Then, the sum obtained is saturated and then stored in \lstinline{ACCX}.\\
During the operation, the lower 4 bits of the access address in register \lstinline{as} are forced to be 0, and then the 16-byte data is loaded from the memory to register \lstinline{qu}. After the access is completed, the value in register \lstinline{as} is incremented by the value in register \lstinline{ad}.\\
At the same time, this instruction also obtains 16-byte unaligned data by concatenating and shifting consecutive aligned data stored in the two registers \lstinline{qs0} and \lstinline{qs1} and stores it to \lstinline{qs0}. The shift byte is stored in special register \lstinline{SAR_BYTE}.\\
\textbf{Operation}
\begin{lstlisting}
  add0[31:0] = qx[ 15:  0] * qy[ 15:  0]
  add1[31:0] = qx[ 31: 16] * qy[ 31: 16]
  ...
  add7[31:0] = qx[127:112] * qy[127:112]
  sum[40:0] = ACCX[39:0] + ad[31:0]d0[31:0] + ad[31:0]d1[31:0] + ... + ad[31:0]d7[31:0]
  ACCX[40:0] = min(max(sum[40:0], 0), 2^{40}-1)

  qu[127:0] = load128({as[31:4],4{0}})
  as[31:0] = as[31:0] + ad[31:0]
  qs0[127:0] = {qs1[127:0], qs0[127:0]} >> {SAR_BYTE[3:0] << 3}
\end{lstlisting}

\clearpage
\subsection{EE.VMULAS.U16.QACC}\label{instruction:EEVMULASU16QACC}
\setlength\tabcolsep{6pt}
\textbf{Instruction Word} \\
\begin{tabular}{|c|c|c|c|c|c|}
	\hline
	000010100 & qy[2] & 1 & qy[1:0] & qx[2:0] & 10000100 \\
	\hline
\end{tabular}

\textbf{Assembler Syntax} \\
EE.VMULAS.U16.QACC qx, qy \\
\textbf{Description} \\
This instruction divides registers \lstinline{qx} and \lstinline{qy} into 8 data segments by 16 bits. Then, the unsigned multiplication result of the 8 sets of segments is added to the corresponding 40-bit data segment in special registers \lstinline{QACC_H} and \lstinline{QACC_L} respectively. The calculated result is saturated to a 40-bit unsigned number and then stored to the corresponding 40-bit data segment in \lstinline{QACC_H} and \lstinline{QACC_L}.\\
\textbf{Operation}
\begin{lstlisting}
  QACC_L[ 39:  0] = min(QACC_L[ 39:  0] + qx[ 15:  0] * qy[ 15:  0], 2^{40}-1)
  QACC_L[ 79: 40] = min(QACC_L[ 79: 40] + qx[ 31: 16] * qy[ 31: 16], 2^{40}-1)
  ...
  QACC_L[159:120] = min(QACC_L[159:120] + qx[ 63: 48] * qy[ 63: 48], 2^{40}-1)
  QACC_H[ 39:  0] = min(QACC_H[ 39:  0] + qx[ 79: 64] * qy[ 79: 64], 2^{40}-1)
  QACC_H[ 79: 40] = min(QACC_H[ 79: 40] + qx[ 95: 80] * qy[ 95: 80], 2^{40}-1)
  ...
  QACC_H[159:120] = min(QACC_H[159:120] + qx[127:112] * qy[127:112], 2^{40}-1)
\end{lstlisting}

\clearpage
\subsection{EE.VMULAS.U16.QACC.LD.IP}\label{instruction:EEVMULASU16QACCLDIP}
\setlength\tabcolsep{6pt}
\textbf{Instruction Word} \\
\begin{tabular}{|c|c|c|c|c|c|c|c|c|c|c|c|c|}
	\hline
	1111 & imm16[5:4] & qu[2:1] & qy[0] & 000 & qu[0] & 101 & qx[1:0] & qy[2:1] & imm16[3:0] & as[3:0] & 111 & qx[2] \\
	\hline
\end{tabular}

\textbf{Assembler Syntax} \\
EE.VMULAS.U16.QACC.LD.IP qu, as, imm16, qx, qy \\
\textbf{Description} \\
This instruction divides registers \lstinline{qx} and \lstinline{qy} into 8 data segments by 16 bits. Then, the unsigned multiplication result of the 8 sets of segments is added to the corresponding 40-bit data segment in special registers \lstinline{QACC_H} and \lstinline{QACC_L} respectively. The calculated result is saturated to a 40-bit unsigned number and then stored to the corresponding 40-bit data segment in \lstinline{QACC_H} and \lstinline{QACC_L}.\\
During the operation, the lower 4 bits of the access address in register \lstinline{as} are forced to be 0, and then the 16-byte data is loaded from the memory to register \lstinline{qu}. After the access is completed, the value in register \lstinline{as} is incremented by 6-bit sign-extended constant in the instruction code segment left-shifted by 4.\\
\textbf{Operation}
\begin{lstlisting}
  QACC_L[ 39:  0] = min(QACC_L[ 39:  0] + qx[ 15:  0] * qy[ 15:  0], 2^{40}-1)
  QACC_L[ 79: 40] = min(QACC_L[ 79: 40] + qx[ 31: 16] * qy[ 31: 16], 2^{40}-1)
  ...
  QACC_L[159:120] = min(QACC_L[159:120] + qx[ 63: 48] * qy[ 63: 48], 2^{40}-1)
  QACC_H[ 39:  0] = min(QACC_H[ 39:  0] + qx[ 79: 64] * qy[ 79: 64], 2^{40}-1)
  QACC_H[ 79: 40] = min(QACC_H[ 79: 40] + qx[ 95: 80] * qy[ 95: 80], 2^{40}-1)
  ...
  QACC_H[159:120] = min(QACC_H[159:120] + qx[127:112] * qy[127:112], 2^{40}-1)

  qu[127:0] = load128({as[31:4],4{0}})
  as[31:0] = as[31:0] + {22{imm16[5]},imm16[5:0],4{0}}
\end{lstlisting}

\clearpage
\subsection{EE.VMULAS.U16.QACC.LD.IP.QUP}\label{instruction:EEVMULASU16QACCLDIPQUP}
\setlength\tabcolsep{5pt}
\textbf{Instruction Word} \\
\begin{tabular}{|c|c|c|c|c|c|c|c|c|c|c|c|c|}
	\hline
	0101 & imm16[5:4] & qu[2:1] & qy[0] & qs0[2:0] & qu[0] & qs1[2:0] & qx[1:0] & qy[2:1] & imm16[3:0] & as[3:0] & 111 & qx[2] \\
	\hline
\end{tabular}

\textbf{Assembler Syntax} \\
EE.VMULAS.U16.QACC.LD.IP.QUP qu, as, imm16, qx, qy, qs0, qs1 \\
\textbf{Description} \\
This instruction divides registers \lstinline{qx} and \lstinline{qy} into 8 data segments by 16 bits. Then, the unsigned multiplication result of the 8 sets of segments is added to the corresponding 40-bit data segment in special registers \lstinline{QACC_H} and \lstinline{QACC_L} respectively. The calculated result is saturated to a 40-bit unsigned number and then stored to the corresponding 40-bit data segment in \lstinline{QACC_H} and \lstinline{QACC_L}.\\
During the operation, the lower 4 bits of the access address in register \lstinline{as} are forced to be 0, and then the 16-byte data is loaded from the memory to register \lstinline{qu}. After the access is completed, the value in register \lstinline{as} is incremented by 6-bit sign-extended constant in the instruction code segment left-shifted by 4.\\
At the same time, this instruction also obtains 16-byte unaligned data by concatenating and shifting consecutive aligned data stored in the two registers \lstinline{qs0} and \lstinline{qs1} and stores it to \lstinline{qs0}. The shift byte is stored in special register \lstinline{SAR_BYTE}.\\
\textbf{Operation}
\begin{lstlisting}
  QACC_L[ 39:  0] = min(QACC_L[ 39:  0] + qx[ 15:  0] * qy[ 15:  0], 2^{40}-1)
  QACC_L[ 79: 40] = min(QACC_L[ 79: 40] + qx[ 31: 16] * qy[ 31: 16], 2^{40}-1)
  ...
  QACC_L[159:120] = min(QACC_L[159:120] + qx[ 63: 48] * qy[ 63: 48], 2^{40}-1)
  QACC_H[ 39:  0] = min(QACC_H[ 39:  0] + qx[ 79: 64] * qy[ 79: 64], 2^{40}-1)
  QACC_H[ 79: 40] = min(QACC_H[ 79: 40] + qx[ 95: 80] * qy[ 95: 80], 2^{40}-1)
  ...
  QACC_H[159:120] = min(QACC_H[159:120] + qx[127:112] * qy[127:112], 2^{40}-1)

  qu[127:0] = load128({as[31:4],4{0}})
  as[31:0] = as[31:0] + {22{imm16[5]},imm16[5:0],4{0}}
  qs0[127:0]  = {qs1[127:0], qs0[127:0]} >> {SAR_BYTE[3:0] << 3}
\end{lstlisting}

\clearpage
\subsection{EE.VMULAS.U16.QACC.LD.XP}\label{instruction:EEVMULASU16QACCLDXP}
\setlength\tabcolsep{6pt}
\textbf{Instruction Word} \\
\begin{tabular}{|c|c|c|c|c|c|c|c|c|c|c|c|}
	\hline
	111100 & qu[2:1] & qy[0] & 001 & qu[0] & 101 & qx[1:0] & qy[2:1] & ad[3:0] & as[3:0] & 111 & qx[2] \\
	\hline
\end{tabular}

\textbf{Assembler Syntax} \\
EE.VMULAS.U16.QACC.LD.XP qu, as, ad, qx, qy \\
\textbf{Description} \\
This instruction divides registers \lstinline{qx} and \lstinline{qy} into 8 data segments by 16 bits. Then, the unsigned multiplication result of the 8 sets of segments is added to the corresponding 40-bit data segment in special registers \lstinline{QACC_H} and \lstinline{QACC_L} respectively. The calculated result is saturated to a 40-bit unsigned number and then stored to the corresponding 40-bit data segment in \lstinline{QACC_H} and \lstinline{QACC_L}.\\
During the operation, the lower 4 bits of the access address in register \lstinline{as} are forced to be 0, and then the 16-byte data is loaded from the memory to register \lstinline{qu}. After the access is completed, the value in register \lstinline{as} is incremented by the value in register \lstinline{ad}.\\
\textbf{Operation}
\begin{lstlisting}
  QACC_L[ 39:  0] = min(QACC_L[ 39:  0] + qx[ 15:  0] * qy[ 15:  0], 2^{40}-1)
  QACC_L[ 79: 40] = min(QACC_L[ 79: 40] + qx[ 31: 16] * qy[ 31: 16], 2^{40}-1)
  ...
  QACC_L[159:120] = min(QACC_L[159:120] + qx[ 63: 48] * qy[ 63: 48], 2^{40}-1)
  QACC_H[ 39:  0] = min(QACC_H[ 39:  0] + qx[ 79: 64] * qy[ 79: 64], 2^{40}-1)
  QACC_H[ 79: 40] = min(QACC_H[ 79: 40] + qx[ 95: 80] * qy[ 95: 80], 2^{40}-1)
  ...
  QACC_H[159:120] = min(QACC_H[159:120] + qx[127:112] * qy[127:112], 2^{40}-1)

  qu[127:0] = load128({as[31:4],4{0}})
  as[31:0] = as[31:0] + ad[31:0]
\end{lstlisting}

\clearpage
\subsection{EE.VMULAS.U16.QACC.LD.XP.QUP}\label{instruction:EEVMULASU16QACCLDXPQUP}
\setlength\tabcolsep{6pt}
\textbf{Instruction Word} \\
\begin{tabular}{|c|c|c|c|c|c|c|c|c|c|c|c|}
	\hline
	110001 & qu[2:1] & qy[0] & qs0[2:0] & qu[0] & qs1[2:0] & qx[1:0] & qy[2:1] & ad[3:0] & as[3:0] & 111 & qx[2] \\
	\hline
\end{tabular}

\textbf{Assembler Syntax} \\
EE.VMULAS.U16.QACC.LD.XP.QUP qu, as, ad, qx, qy, qs0, qs1 \\
\textbf{Description} \\
This instruction divides registers \lstinline{qx} and \lstinline{qy} into 8 data segments by 16 bits. Then, the unsigned multiplication result of the 8 sets of segments is added to the corresponding 40-bit data segment in special registers \lstinline{QACC_H} and \lstinline{QACC_L} respectively. The calculated result is saturated to a 40-bit unsigned number and then stored to the corresponding 40-bit data segment in \lstinline{QACC_H} and \lstinline{QACC_L}.\\
During the operation, the lower 4 bits of the access address in register \lstinline{as} are forced to be 0, and then the 16-byte data is loaded from the memory to register \lstinline{qu}. After the access is completed, the value in register \lstinline{as} is incremented by the value in register \lstinline{ad}.\\
At the same time, this instruction also obtains 16-byte unaligned data by concatenating and shifting consecutive aligned data stored in the two registers \lstinline{qs0} and \lstinline{qs1} and stores it to \lstinline{qs0}. The shift byte is stored in special register \lstinline{SAR_BYTE}.\\
\textbf{Operation}
\begin{lstlisting}
  QACC_L[ 39:  0] = min(QACC_L[ 39:  0] + qx[ 15:  0] * qy[ 15:  0], 2^{40}-1)
  QACC_L[ 79: 40] = min(QACC_L[ 79: 40] + qx[ 31: 16] * qy[ 31: 16], 2^{40}-1)
  ...
  QACC_L[159:120] = min(QACC_L[159:120] + qx[ 63: 48] * qy[ 63: 48], 2^{40}-1)
  QACC_H[ 39:  0] = min(QACC_H[ 39:  0] + qx[ 79: 64] * qy[ 79: 64], 2^{40}-1)
  QACC_H[ 79: 40] = min(QACC_H[ 79: 40] + qx[ 95: 80] * qy[ 95: 80], 2^{40}-1)
  ...
  QACC_H[159:120] = min(QACC_H[159:120] + qx[127:112] * qy[127:112], 2^{40}-1)

  qu[127:0] = load128({as[31:4],4{0}})
  as[31:0] = as[31:0] + ad[31:0]
  qs0[127:0]  = {qs1[127:0], qs0[127:0]} >> {SAR_BYTE[3:0] << 3}
\end{lstlisting}

\clearpage
\subsection{EE.VMULAS.U16.QACC.LDBC.INCP}\label{instruction:EEVMULASU16QACCLDBCINCP}
\setlength\tabcolsep{6pt}
\textbf{Instruction Word} \\
\begin{tabular}{|c|c|c|c|c|c|c|c|c|c|}
	\hline
	110 & qu[2] & 0111 & qu[1] & qy[2] & qu[0] & qy[1:0] & qx[2:0] & as[3:0] & 0100 \\
	\hline
\end{tabular}

\textbf{Assembler Syntax} \\
EE.VMULAS.U16.QACC.LDBC.INCP qu, as, qx, qy \\
\textbf{Description} \\
This instruction divides registers \lstinline{qx} and \lstinline{qy} into 8 data segments by 16 bits. Then, the unsigned multiplication result of the 8 sets of segments is added to the corresponding 40-bit data segment in special registers \lstinline{QACC_H} and \lstinline{QACC_L} respectively. The calculated result is saturated to a 40-bit unsigned number and then stored to the corresponding 40-bit data segment in \lstinline{QACC_H} and \lstinline{QACC_L}.\\
At the same time, this instruction forces the lower 1 bit of the access address in register \lstinline{as} to 0, loads 16-bit data from memory, and broadcasts it to the eight 16-bit data segments in register \lstinline{qu}. After the access, the value in register \lstinline{as} is incremented by 2.\\
\textbf{Operation}
\begin{lstlisting}
  QACC_L[ 39:  0] = min(QACC_L[ 39:  0] + qx[ 15:  0] * qy[ 15:  0], 2^{40}-1)
  QACC_L[ 79: 40] = min(QACC_L[ 79: 40] + qx[ 31: 16] * qy[ 31: 16], 2^{40}-1)
  ...
  QACC_L[159:120] = min(QACC_L[159:120] + qx[ 63: 48] * qy[ 63: 48], 2^{40}-1)
  QACC_H[ 39:  0] = min(QACC_H[ 39:  0] + qx[ 79: 64] * qy[ 79: 64], 2^{40}-1)
  QACC_H[ 79: 40] = min(QACC_H[ 79: 40] + qx[ 95: 80] * qy[ 95: 80], 2^{40}-1)
  ...
  QACC_H[159:120] = min(QACC_H[159:120] + qx[127:112] * qy[127:112], 2^{40}-1)

  qu[127:0] = {8{load16({as[31:1],1{0}})}}
  as[31:0] = as[31:0] + 2
\end{lstlisting}

\clearpage
\subsection{EE.VMULAS.U16.QACC.LDBC.INCP.QUP}\label{instruction:EEVMULASU16QACCLDBCINCPQUP}
\setlength\tabcolsep{6pt}
\textbf{Instruction Word} \\
\begin{tabular}{|c|c|c|c|c|c|c|c|c|c|c|c|}
	\hline
	111000 & qu[2:1] & qy[0] & qs0[2:0] & qu[0] & qs1[2:0] & qx[1:0] & qy[2:1] & 1010 & as[3:0] & 111 & qx[2] \\
	\hline
\end{tabular}

\textbf{Assembler Syntax} \\
EE.VMULAS.U16.QACC.LDBC.INCP.QUP qu, as, qx, qy, qs0, qs1 \\
\textbf{Description} \\
This instruction divides registers \lstinline{qx} and \lstinline{qy} into 8 data segments by 16 bits. Then, the unsigned multiplication result of the 8 sets of segments is added to the corresponding 40-bit data segment in special registers \lstinline{QACC_H} and \lstinline{QACC_L} respectively. The calculated result is saturated to a 40-bit unsigned number and then stored to the corresponding 40-bit data segment in \lstinline{QACC_H} and \lstinline{QACC_L}.\\
At the same time, this instruction forces the lower 1 bit of the access address in register \lstinline{as} to 0, loads 16-bit data from memory, and broadcasts it to the eight 16-bit data segments in register \lstinline{qu}. After the access, the value in register \lstinline{as} is incremented by 2.\\
At the same time, this instruction also obtains 16-byte unaligned data by concatenating and shifting consecutive aligned data stored in the two registers \lstinline{qs0} and \lstinline{qs1} and stores it to \lstinline{qs0}. The shift byte is stored in special register \lstinline{SAR_BYTE}.\\
\textbf{Operation}
\begin{lstlisting}
  QACC_L[ 39:  0] = min(QACC_L[ 39:  0] + qx[ 15:  0] * qy[ 15:  0], 2^{40}-1)
  QACC_L[ 79: 40] = min(QACC_L[ 79: 40] + qx[ 31: 16] * qy[ 31: 16], 2^{40}-1)
  ...
  QACC_L[159:120] = min(QACC_L[159:120] + qx[ 63: 48] * qy[ 63: 48], 2^{40}-1)
  QACC_H[ 39:  0] = min(QACC_H[ 39:  0] + qx[ 79: 64] * qy[ 79: 64], 2^{40}-1)
  QACC_H[ 79: 40] = min(QACC_H[ 79: 40] + qx[ 95: 80] * qy[ 95: 80], 2^{40}-1)
  ...
  QACC_H[159:120] = min(QACC_H[159:120] + qx[127:112] * qy[127:112], 2^{40}-1)

  qu[127:0] = {8{load16({as[31:1],1{0}})}}
  as[31:0] = as[31:0] + 2
  qs0[127:0]  = {qs1[127:0], qs0[127:0]} >> {SAR_BYTE[3:0] << 3}
\end{lstlisting}

\clearpage
\subsection{EE.VMULAS.U8.ACCX}\label{instruction:EEVMULASU8ACCX}
\setlength\tabcolsep{6pt}
\textbf{Instruction Word} \\
\begin{tabular}{|c|c|c|c|c|c|}
	\hline
	000010100 & qy[2] & 0 & qy[1:0] & qx[2:0] & 11000100 \\
	\hline
\end{tabular}

\textbf{Assembler Syntax} \\
EE.VMULAS.U8.ACCX qx, qy \\
\textbf{Description} \\
This instruction divides registers \lstinline{qx} and \lstinline{qy} into 16 data segments by 8 bits. Then, it performs an unsigned multiply-accumulate operation on the 16 sets of segments respectively. The accumulated result is added to the value in special register \lstinline{ACCX}. Then, the sum obtained is saturated and then stored in \lstinline{ACCX}.\\
\textbf{Operation}
\begin{lstlisting}
  add0[15:0] = qx[  7:  0] * qy[  7:  0]
  add1[15:0] = qx[ 15:  8] * qy[ 15:  8]
  ...
  add15[15:0] = qx[127:120] * qy[127:120]
  sum[40:0] = ACCX[39:0] + ad[31:0]d0[15:0] + ad[31:0]d1[15:0] + ... + ad[31:0]d15[15:0]

  ACCX[40:0] = min(max(sum[40:0], 0), 2^{40}-1)
\end{lstlisting}

\clearpage
\subsection{EE.VMULAS.U8.ACCX.LD.IP}\label{instruction:EEVMULASU8ACCXLDIP}
\setlength\tabcolsep{6pt}
\textbf{Instruction Word} \\
\begin{tabular}{|c|c|c|c|c|c|c|c|c|c|c|c|c|}
	\hline
	1111 & imm16[5:4] & qu[2:1] & qy[0] & 000 & qu[0] & 110 & qx[1:0] & qy[2:1] & imm16[3:0] & as[3:0] & 111 & qx[2] \\
	\hline
\end{tabular}

\textbf{Assembler Syntax} \\
EE.VMULAS.U8.ACCX.LD.IP qu, as, -512..496, qx, qy \\
\textbf{Description} \\
This instruction divides registers \lstinline{qx} and \lstinline{qy} into 16 data segments by 8 bits. Then, it performs an unsigned multiply-accumulate operation on the 16 sets of segments respectively. The accumulated result is added to the value in special register \lstinline{ACCX}. Then, the sum obtained is saturated and then stored in \lstinline{ACCX}. \\ During the operation, the lower 4 bits of the access address in register \lstinline{as} are forced to be 0, and then the 16-byte data is loaded from the memory to register \lstinline{qu}. After the access is completed, the value in register \lstinline{as} is incremented by 6-bit sign-extended constant in the instruction code segment left-shifted by 4.\\
\textbf{Operation}
\begin{lstlisting}
  add0[15:0] = qx[  7:  0] * qy[  7:  0]
  add1[15:0] = qx[ 15:  8] * qy[ 15:  8]
  ...
  add15[15:0] = qx[127:120] * qy[127:120]
  sum[40:0] = ACCX[39:0] + ad[31:0]d0[15:0] + ad[31:0]d1[15:0] + ... + ad[31:0]d15[15:0]
  ACCX[40:0] = min(max(sum[40:0], 0), 2^{40}-1)

  qu[127:0] = load128({as[31:4],4{0}})
  as[31:0] = as[31:0] + {20{imm16[7]},imm16[7:0],4{0}}
\end{lstlisting}

\clearpage
\subsection{EE.VMULAS.U8.ACCX.LD.IP.QUP}\label{instruction:EEVMULASU8ACCXLDIPQUP}
\setlength\tabcolsep{5pt}
\textbf{Instruction Word} \\
\begin{tabular}{|c|c|c|c|c|c|c|c|c|c|c|c|c|}
	\hline
	0110 & imm16[5:4] & qu[2:1] & qy[0] & qs0[2:0] & qu[0] & qs1[2:0] & qx[1:0] & qy[2:1] & imm16[3:0] & as[3:0] & 111 & qx[2] \\
	\hline
\end{tabular}

\textbf{Assembler Syntax} \\
EE.VMULAS.U8.ACCX.LD.IP.QUP qu, as, imm16, qx, qy, qs0, qs1 \\
\textbf{Description} \\
This instruction divides registers \lstinline{qx} and \lstinline{qy} into 16 data segments by 8 bits. Then, it performs an unsigned multiply-accumulate operation on the 16 sets of segments respectively. The accumulated result is added to the value in special register \lstinline{ACCX}. Then, the sum obtained is saturated and then stored in \lstinline{ACCX}.\\
During the operation, the lower 4 bits of the access address in register \lstinline{as} are forced to be 0, and then the 16-byte data is loaded from the memory to register \lstinline{qu}. After the access is completed, the value in register \lstinline{as} is incremented by 6-bit sign-extended constant in the instruction code segment left-shifted by 4.\\
At the same time, this instruction also obtains 16-byte unaligned data by concatenating and shifting consecutive aligned data stored in the two registers \lstinline{qs0} and \lstinline{qs1} and stores it to \lstinline{qs0}. The shift byte is stored in special register \lstinline{SAR_BYTE}.\\
\textbf{Operation}
\begin{lstlisting}
  add0[15:0] = qx[  7:  0] * qy[  7:  0]
  add1[15:0] = qx[ 15:  8] * qy[ 15:  8]
  ...
  add15[15:0] = qx[127:120] * qy[127:120]
  sum[40:0] = ACCX[39:0] + ad[31:0]d0[15:0] + ad[31:0]d1[15:0] + ... + ad[31:0]d15[15:0]
  ACCX[40:0] = min(max(sum[40:0], 0), 2^{40}-1)

  qu[127:0] = load128({as[31:4],4{0}})
  as[31:0] = as[31:0] + {20{imm16[7]},imm16[7:0],4{0}}
  qs0[127:0]  = {qs1[127:0], qs0[127:0]} >> {SAR_BYTE[3:0] << 3}
\end{lstlisting}

\clearpage
\subsection{EE.VMULAS.U8.ACCX.LD.XP}\label{instruction:EEVMULASU8ACCXLDXP}
\setlength\tabcolsep{6pt}
\textbf{Instruction Word} \\
\begin{tabular}{|c|c|c|c|c|c|c|c|c|c|c|c|}
	\hline
	111100 & qu[2:1] & qy[0] & 001 & qu[0] & 110 & qx[1:0] & qy[2:1] & ad[3:0] & as[3:0] & 111 & qx[2] \\
	\hline
\end{tabular}

\textbf{Assembler Syntax} \\
EE.VMULAS.U8.ACCX.LD.XP qu, as, ad, qx, qy \\
\textbf{Description} \\
This instruction divides registers \lstinline{qx} and \lstinline{qy} into 16 data segments by 8 bits. Then, it performs an unsigned multiply-accumulate operation on the 16 sets of segments respectively. The accumulated result is added to the value in special register \lstinline{ACCX}. Then, the sum obtained is saturated and then stored in \lstinline{ACCX}.\\
During the operation, the lower 4 bits of the access address in register \lstinline{as} are forced to be 0, and then the 16-byte data is loaded from the memory to register \lstinline{qu}. After the access is completed, the value in register \lstinline{as} is incremented by the value in register \lstinline{ad}.\\
\textbf{Operation}
\begin{lstlisting}
  add0[15:0] = qx[  7:  0] * qy[  7:  0]
  add1[15:0] = qx[ 15:  8] * qy[ 15:  8]
  ...
  add15[15:0] = qx[127:120] * qy[127:120]
  sum[40:0] = ACCX[39:0] + ad[31:0]d0[15:0] + ad[31:0]d1[15:0] + ... + ad[31:0]d15[15:0]
  ACCX[40:0] = min(max(sum[40:0], 0), 2^{40}-1)

  qu[127:0] = load128({as[31:4],4{0}})
  as[31:0] = as[31:0] + ad[31:0]
\end{lstlisting}

\clearpage
\subsection{EE.VMULAS.U8.ACCX.LD.XP.QUP}\label{instruction:EEVMULASU8ACCXLDXPQUP}
\setlength\tabcolsep{6pt}
\textbf{Instruction Word} \\
\begin{tabular}{|c|c|c|c|c|c|c|c|c|c|c|c|}
	\hline
	110010 & qu[2:1] & qy[0] & qs0[2:0] & qu[0] & qs1[2:0] & qx[1:0] & qy[2:1] & ad[3:0] & as[3:0] & 111 & qx[2] \\
	\hline
\end{tabular}

\textbf{Assembler Syntax} \\
EE.VMULAS.U8.ACCX.LD.XP.QUP qu, as, ad, qx, qy, qs0, qs1 \\
\textbf{Description} \\
This instruction divides registers \lstinline{qx} and \lstinline{qy} into 16 data segments by 8 bits. Then, it performs an unsigned multiply-accumulate operation on the 16 sets of segments respectively. The accumulated result is added to the value in special register \lstinline{ACCX}. Then, the sum obtained is saturated and then stored in \lstinline{ACCX}.\\
During the operation, the lower 4 bits of the access address in register \lstinline{as} are forced to be 0, and then the 16-byte data is loaded from the memory to register \lstinline{qu}. After the access is completed, the value in register \lstinline{as} is incremented by the value in register \lstinline{ad}.\\
At the same time, this instruction also obtains 16-byte unaligned data by concatenating and shifting consecutive aligned data stored in the two registers \lstinline{qs0} and \lstinline{qs1} and stores it to \lstinline{qs0}. The shift byte is stored in special register \lstinline{SAR_BYTE}.\\
\textbf{Operation}
\begin{lstlisting}
  add0[15:0] = qx[  7:  0] * qy[  7:  0]
  add1[15:0] = qx[ 15:  8] * qy[ 15:  8]
  ...
  add15[15:0] = qx[127:120] * qy[127:120]
  sum[40:0] = ACCX[39:0] + ad[31:0]d0[15:0] + ad[31:0]d1[15:0] + ... + ad[31:0]d15[15:0]
  ACCX[40:0] = min(max(sum[40:0], 0), 2^{40}-1)

  qu[127:0] = load128({as[31:4],4{0}})
  as[31:0] = as[31:0] + ad[31:0]
  qs0[127:0] = {qs1[127: 0], qs0[127: 0]} >> {SAR_BYTE[3:0] << 3}
\end{lstlisting}

\clearpage
\subsection{EE.VMULAS.U8.QACC}\label{instruction:EEVMULASU8QACC}
\setlength\tabcolsep{6pt}
\textbf{Instruction Word} \\
\begin{tabular}{|c|c|c|c|c|c|}
	\hline
	000010100 & qy[2] & 1 & qy[1:0] & qx[2:0] & 11000100 \\
	\hline
\end{tabular}

\textbf{Assembler Syntax} \\
EE.VMULAS.U8.QACC qx, qy \\
\textbf{Description} \\
This instruction divides registers \lstinline{qx} and \lstinline{qy} into 16 data segments by 8 bits. Then, the unsigned multiplication result of the 16 sets of segments is added to the corresponding 20-bit data segment in special registers \lstinline{QACC_H} and \lstinline{QACC_L} respectively. The calculated result is saturated to a 20-bit unsigned number and then stored to the corresponding 20-bit data segment in \lstinline{QACC_H} and \lstinline{QACC_L}.\\
\textbf{Operation}
\begin{lstlisting}
  QACC_L[ 19:  0] = min(QACC_L[ 19:  0] + qx[  7:  0] * qy[  7:  0], 2^{20}-1)
  QACC_L[ 39: 20] = min(QACC_L[ 39: 20] + qx[ 15:  8] * qy[ 15:  8], 2^{w0}-1)
  ...
  QACC_L[159:140] = min(QACC_L[159:140] + qx[ 63: 56] * qy[ 63: 56], 2^{20}-1)
  QACC_H[ 19:  0] = min(QACC_H[ 19:  0] + qx[ 71: 64] * qy[ 71: 64], 2^{20}-1)
  QACC_H[ 39: 20] = min(QACC_H[ 39: 20] + qx[ 79: 72] * qy[ 79: 72], 2^{20}-1)
  ...
  QACC_H[159:140] = min(QACC_H[159:140] + qx[127:120] * qy[127:120], 2^{20}-1)
\end{lstlisting}

\clearpage
\subsection{EE.VMULAS.U8.QACC.LD.IP}\label{instruction:EEVMULASU8QACCLDIP}
\setlength\tabcolsep{6pt}
\textbf{Instruction Word} \\
\begin{tabular}{|c|c|c|c|c|c|c|c|c|c|c|c|c|}
	\hline
	1111 & imm16[5:4] & qu[2:1] & qy[0] & 000 & qu[0] & 111 & qx[1:0] & qy[2:1] & imm16[3:0] & as[3:0] & 111 & qx[2] \\
	\hline
\end{tabular}

\textbf{Assembler Syntax} \\
EE.VMULAS.U8.QACC.LD.IP qu, as, imm16, qx, qy \\
\textbf{Description} \\
This instruction divides registers \lstinline{qx} and \lstinline{qy} into 16 data segments by 8 bits. Then, the unsigned multiplication result of the 16 sets of segments is added to the corresponding 20-bit data segment in special registers \lstinline{QACC_H} and \lstinline{QACC_L} respectively. The calculated result is saturated to a 20-bit unsigned number and then stored to the corresponding 20-bit data segment in \lstinline{QACC_H} and \lstinline{QACC_L}.\\
During the operation, the lower 4 bits of the access address in register \lstinline{as} are forced to be 0, and then the 16-byte data is loaded from the memory to register \lstinline{qu}. After the access is completed, the value in register \lstinline{as} is incremented by 6-bit sign-extended constant in the instruction code segment left-shifted by 4.\\
\textbf{Operation}
\begin{lstlisting}
  QACC_L[ 19:  0] = min(QACC_L[ 19:  0] + qx[  7:  0] * qy[  7:  0], 2^{20}-1)
  QACC_L[ 39: 20] = min(QACC_L[ 39: 20] + qx[ 15:  8] * qy[ 15:  8], 2^{w0}-1)
  ...
  QACC_L[159:140] = min(QACC_L[159:140] + qx[ 63: 56] * qy[ 63: 56], 2^{20}-1)
  QACC_H[ 19:  0] = min(QACC_H[ 19:  0] + qx[ 71: 64] * qy[ 71: 64], 2^{20}-1)
  QACC_H[ 39: 20] = min(QACC_H[ 39: 20] + qx[ 79: 72] * qy[ 79: 72], 2^{20}-1)
  ...
  QACC_H[159:140] = min(QACC_H[159:140] + qx[127:120] * qy[127:120], 2^{20}-1)

  qu[127:0] = load128({as[31:4],4{0}})
  as[31:0] = as[31:0] + {22{imm16[5]},imm16[5:0],4{0}}
\end{lstlisting}

\clearpage
\subsection{EE.VMULAS.U8.QACC.LD.IP.QUP}\label{instruction:EEVMULASU8QACCLDIPQUP}
\setlength\tabcolsep{5pt}
\textbf{Instruction Word} \\
\begin{tabular}{|c|c|c|c|c|c|c|c|c|c|c|c|c|}
	\hline
	0111 & imm16[5:4] & qu[2:1] & qy[0] & qs0[2:0] & qu[0] & qs1[2:0] & qx[1:0] & qy[2:1] & imm16[3:0] & as[3:0] & 111 & qx[2] \\
	\hline
\end{tabular}

\textbf{Assembler Syntax} \\
EE.VMULAS.U8.QACC.LD.IP.QUP qu, as, imm16, qx, qy, qs0, qs1 \\
\textbf{Description} \\
This instruction divides registers \lstinline{qx} and \lstinline{qy} into 16 data segments by 8 bits. Then, the unsigned multiplication result of the 16 sets of segments is added to the corresponding 20-bit data segment in special registers \lstinline{QACC_H} and \lstinline{QACC_L} respectively. The calculated result is saturated to a 20-bit unsigned number and then stored to the corresponding 20-bit data segment in \lstinline{QACC_H} and \lstinline{QACC_L}.\\
During the operation, the lower 4 bits of the access address in register \lstinline{as} are forced to be 0, and then the 16-byte data is loaded from the memory to register \lstinline{qu}. After the access is completed, the value in register \lstinline{as} is incremented by 6-bit sign-extended constant in the instruction code segment left-shifted by 4.\\
At the same time, this instruction also obtains 16-byte unaligned data by concatenating and shifting consecutive aligned data stored in the two registers \lstinline{qs0} and \lstinline{qs1} and stores it to \lstinline{qs0}. The shift byte is stored in special register \lstinline{SAR_BYTE}.\\
\textbf{Operation}
\begin{lstlisting}
  QACC_L[ 19:  0] = min(QACC_L[ 19:  0] + qx[  7:  0] * qy[  7:  0], 2^{20}-1)
  QACC_L[ 39: 20] = min(QACC_L[ 39: 20] + qx[ 15:  8] * qy[ 15:  8], 2^{w0}-1)
  ...
  QACC_L[159:140] = min(QACC_L[159:140] + qx[ 63: 56] * qy[ 63: 56], 2^{20}-1)
  QACC_H[ 19:  0] = min(QACC_H[ 19:  0] + qx[ 71: 64] * qy[ 71: 64], 2^{20}-1)
  QACC_H[ 39: 20] = min(QACC_H[ 39: 20] + qx[ 79: 72] * qy[ 79: 72], 2^{20}-1)
  ...
  QACC_H[159:140] = min(QACC_H[159:140] + qx[127:120] * qy[127:120], 2^{20}-1)

  qu[127:0] = load128({as[31:4],4{0}})
  as[31:0] = as[31:0] + {22{imm16[5]},imm16[5:0],4{0}}
  qs0[127:0]  = {qs1[127:0], qs0[127:0]} >> {SAR_BYTE[3:0] << 3}
\end{lstlisting}

\clearpage
\subsection{EE.VMULAS.U8.QACC.LD.XP}\label{instruction:EEVMULASU8QACCLDXP}
\setlength\tabcolsep{6pt}
\textbf{Instruction Word} \\
\begin{tabular}{|c|c|c|c|c|c|c|c|c|c|c|c|}
	\hline
	111100 & qu[2:1] & qy[0] & 001 & qu[0] & 111 & qx[1:0] & qy[2:1] & ad[3:0] & as[3:0] & 111 & qx[2] \\
	\hline
\end{tabular}

\textbf{Assembler Syntax} \\
EE.VMULAS.U8.QACC.LD.XP qu, as, ad, qx, qy \\
\textbf{Description} \\
This instruction divides registers \lstinline{qx} and \lstinline{qy} into 16 data segments by 8 bits. Then, the unsigned multiplication result of the 16 sets of segments is added to the corresponding 20-bit data segment in special registers \lstinline{QACC_H} and \lstinline{QACC_L} respectively. The calculated result is saturated to a 20-bit unsigned number and then stored to the corresponding 20-bit data segment in \lstinline{QACC_H} and \lstinline{QACC_L}.\\
During the operation, the lower 4 bits of the access address in register \lstinline{as} are forced to be 0, and then the 16-byte data is loaded from the memory to register \lstinline{qu}. After the access is completed, the value in register \lstinline{as} is incremented by the value in register \lstinline{ad}.\\
\textbf{Operation}
\begin{lstlisting}
  QACC_L[ 19:  0] = min(QACC_L[ 19:  0] + qx[  7:  0] * qy[  7:  0], 2^{20}-1)
  QACC_L[ 39: 20] = min(QACC_L[ 39: 20] + qx[ 15:  8] * qy[ 15:  8], 2^{w0}-1)
  ...
  QACC_L[159:140] = min(QACC_L[159:140] + qx[ 63: 56] * qy[ 63: 56], 2^{20}-1)
  QACC_H[ 19:  0] = min(QACC_H[ 19:  0] + qx[ 71: 64] * qy[ 71: 64], 2^{20}-1)
  QACC_H[ 39: 20] = min(QACC_H[ 39: 20] + qx[ 79: 72] * qy[ 79: 72], 2^{20}-1)
  ...
  QACC_H[159:140] = min(QACC_H[159:140] + qx[127:120] * qy[127:120], 2^{20}-1)

  qu[127:0] = load128({as[31:4],4{0}})
  as[31:0] = as[31:0] + ad[31:0]
\end{lstlisting}

\clearpage
\subsection{EE.VMULAS.U8.QACC.LD.XP.QUP}\label{instruction:EEVMULASU8QACCLDXPQUP}
\setlength\tabcolsep{6pt}
\textbf{Instruction Word} \\
\begin{tabular}{|c|c|c|c|c|c|c|c|c|c|c|c|}
	\hline
	110011 & qu[2:1] & qy[0] & qs0[2:0] & qu[0] & qs1[2:0] & qx[1:0] & qy[2:1] & ad[3:0] & as[3:0] & 111 & qx[2] \\
	\hline
\end{tabular}

\textbf{Assembler Syntax} \\
EE.VMULAS.U8.QACC.LD.XP.QUP qu, as, ad, qx, qy, qs0, qs1 \\
\textbf{Description} \\
This instruction divides registers \lstinline{qx} and \lstinline{qy} into 16 data segments by 8 bits. Then, the unsigned multiplication result of the 16 sets of segments is added to the corresponding 20-bit data segment in special registers \lstinline{QACC_H} and \lstinline{QACC_L} respectively. The calculated result is saturated to a 20-bit unsigned number and then stored to the corresponding 20-bit data segment in \lstinline{QACC_H} and \lstinline{QACC_L}.\\
During the operation, the lower 4 bits of the access address in register \lstinline{as} are forced to be 0, and then the 16-byte data is loaded from the memory to register \lstinline{qu}. After the access is completed, the value in register \lstinline{as} is incremented by the value in register \lstinline{ad}.\\
At the same time, this instruction also obtains 16-byte unaligned data by concatenating and shifting consecutive aligned data stored in the two registers \lstinline{qs0} and \lstinline{qs1} and stores it to \lstinline{qs0}. The shift byte is stored in special register \lstinline{SAR_BYTE}.\\
\textbf{Operation}
\begin{lstlisting}
  QACC_L[ 19:  0] = min(QACC_L[ 19:  0] + qx[  7:  0] * qy[  7:  0], 2^{20}-1)
  QACC_L[ 39: 20] = min(QACC_L[ 39: 20] + qx[ 15:  8] * qy[ 15:  8], 2^{w0}-1)
  ...
  QACC_L[159:140] = min(QACC_L[159:140] + qx[ 63: 56] * qy[ 63: 56], 2^{20}-1)
  QACC_H[ 19:  0] = min(QACC_H[ 19:  0] + qx[ 71: 64] * qy[ 71: 64], 2^{20}-1)
  QACC_H[ 39: 20] = min(QACC_H[ 39: 20] + qx[ 79: 72] * qy[ 79: 72], 2^{20}-1)
  ...
  QACC_H[159:140] = min(QACC_H[159:140] + qx[127:120] * qy[127:120], 2^{20}-1)

  qu[127:0] = load128({as[31:4],4{0}})
  as[31:0] = as[31:0] + ad[31:0]
  qs0[127:0]  = {qs1[127:0], qs0[127:0]} >> {SAR_BYTE[3:0] << 3}
\end{lstlisting}

\clearpage
\subsection{EE.VMULAS.U8.QACC.LDBC.INCP}\label{instruction:EEVMULASU8QACCLDBCINCP}
\setlength\tabcolsep{6pt}
\textbf{Instruction Word} \\
\begin{tabular}{|c|c|c|c|c|c|c|c|c|c|}
	\hline
	111 & qu[2] & 0111 & qu[1] & qy[2] & qu[0] & qy[1:0] & qx[2:0] & as[3:0] & 0100 \\
	\hline
\end{tabular}

\textbf{Assembler Syntax} \\
EE.VMULAS.U8.QACC.LDBC.INCP qu, as, qx, qy \\
\textbf{Description} \\
This instruction divides registers \lstinline{qx} and \lstinline{qy} into 16 data segments by 8 bits. Then, the unsigned multiplication result of the 16 sets of segments is added to the corresponding 20-bit data segment in special registers \lstinline{QACC_H} and \lstinline{QACC_L} respectively. The calculated result is saturated to a 20-bit unsigned number and then stored to the corresponding 20-bit data segment in \lstinline{QACC_H} and \lstinline{QACC_L}.\\
At the same time, this instruction loads 8-bit data from memory at the address given by the access register \lstinline{as} and broadcasts it to the 16 8-bit data segments in register \lstinline{qu}. After the access, the value in register \lstinline{as} is incremented by 1.\\
\textbf{Operation}
\begin{lstlisting}
  QACC_L[ 19:  0] = min(QACC_L[ 19:  0] + qx[  7:  0] * qy[  7:  0], 2^{20}-1)
  QACC_L[ 39: 20] = min(QACC_L[ 39: 20] + qx[ 15:  8] * qy[ 15:  8], 2^{w0}-1)
  ...
  QACC_L[159:140] = min(QACC_L[159:140] + qx[ 63: 56] * qy[ 63: 56], 2^{20}-1)
  QACC_H[ 19:  0] = min(QACC_H[ 19:  0] + qx[ 71: 64] * qy[ 71: 64], 2^{20}-1)
  QACC_H[ 39: 20] = min(QACC_H[ 39: 20] + qx[ 79: 72] * qy[ 79: 72], 2^{20}-1)
  ...
  QACC_H[159:140] = min(QACC_H[159:140] + qx[127:120] * qy[127:120], 2^{20}-1)

  qu[127:0] = {16{load8(as[31:0])}}
  as[31:0] = as[31:0] + 1
\end{lstlisting}

\clearpage
\subsection{EE.VMULAS.U8.QACC.LDBC.INCP.QUP}\label{instruction:EEVMULASU8QACCLDBCINCPQUP}
\setlength\tabcolsep{6pt}
\textbf{Instruction Word} \\
\begin{tabular}{|c|c|c|c|c|c|c|c|c|c|c|c|}
	\hline
	111000 & qu[2:1] & qy[0] & qs0[2:0] & qu[0] & qs1[2:0] & qx[1:0] & qy[2:1] & 1011 & as[3:0] & 111 & qx[2] \\
	\hline
\end{tabular}

\textbf{Assembler Syntax} \\
EE.VMULAS.U8.QACC.LDBC.INCP.QUP qu, as, qx, qy, qs0, qs1 \\
\textbf{Description} \\
This instruction divides registers \lstinline{qx} and \lstinline{qy} into 16 data segments by 8 bits. Then, the unsigned multiplication result of the 16 sets of segments is added to the corresponding 20-bit data segment in special registers \lstinline{QACC_H} and \lstinline{QACC_L} respectively. The calculated result is saturated to a 20-bit unsigned number and then stored to the corresponding 20-bit data segment in \lstinline{QACC_H} and \lstinline{QACC_L}.\\
At the same time, this instruction loads 8-bit data from memory at the address given by the access register \lstinline{as} and broadcasts it to the 16 8-bit data segments in register \lstinline{qu}. After the access, the value in register \lstinline{as} is incremented by 1.\\
At the same time, this instruction also obtains 16-byte unaligned data by concatenating and shifting consecutive aligned data stored in the two registers \lstinline{qs0} and \lstinline{qs1} and stores it to \lstinline{qs0}. The shift byte is stored in special register \lstinline{SAR_BYTE}.\\
\textbf{Operation}
\begin{lstlisting}
  QACC_L[ 19:  0] = min(QACC_L[ 19:  0] + qx[  7:  0] * qy[  7:  0], 2^{20}-1)
  QACC_L[ 39: 20] = min(QACC_L[ 39: 20] + qx[ 15:  8] * qy[ 15:  8], 2^{w0}-1)
  ...
  QACC_L[159:140] = min(QACC_L[159:140] + qx[ 63: 56] * qy[ 63: 56], 2^{20}-1)
  QACC_H[ 19:  0] = min(QACC_H[ 19:  0] + qx[ 71: 64] * qy[ 71: 64], 2^{20}-1)
  QACC_H[ 39: 20] = min(QACC_H[ 39: 20] + qx[ 79: 72] * qy[ 79: 72], 2^{20}-1)
  ...
  QACC_H[159:140] = min(QACC_H[159:140] + qx[127:120] * qy[127:120], 2^{20}-1)

  qu[127:0] = {16{load8(as[31:0])}}
  as[31:0] = as[31:0] + 1
  qs0[127:0]  = {qs1[127:0], qs0[127:0]} >> {SAR_BYTE[3:0] << 3}
\end{lstlisting}

\clearpage
\subsection{EE.VPRELU.S16}\label{instruction:EEVPRELUS16}
\setlength\tabcolsep{6pt}
\textbf{Instruction Word} \\
\begin{tabular}{|c|c|c|c|c|c|c|c|c|c|}
	\hline
	10 & qz[2:1] & 1100 & qz[0] & qy[2] & 0 & qy[1:0] & qx[2:0] & ay[3:0] & 0100 \\
	\hline
\end{tabular}

\textbf{Assembler Syntax} \\
EE.VPRELU.S16 qz, qx, qy, ay \\
\textbf{Description} \\
This instruction divides register \lstinline{qx} into 8 data segments by 16 bits. If the value of the segment is not greater than 0, it will be multiplied by the value of the corresponding 16-bit segment in register \lstinline{qy} right-shifted by the lower 6-bit value in register \lstinline{ay}, and then the result obtained will be assigned to the corresponding 16-bit data segment in register \lstinline{qz}. Otherwise, the value of the segment in \lstinline{qx} will be assigned to \lstinline{qz}.\\
\textbf{Operation}
\begin{lstlisting}[escapeinside=`?]
  qz[ 15:  0] = (qx[ 15:  0]<=0) ? (qx[ 15:  0] * qy[ 15:  0]) >> ay[5:0] : qx[ 15:  0]
  qz[ 31: 16] = (qx[ 31: 16]<=0) ? (qx[ 31: 16] * qy[ 31: 16]) >> ay[5:0] : qx[ 31: 16]
  ...
  qz[127:112] = (qx[127:112]<=0) ? (qx[127:112] * qy[127:112]) >> ay[5:0] : qx[127:112]
\end{lstlisting}

\clearpage
\subsection{EE.VPRELU.S8}\label{instruction:EEVPRELUS8}
\setlength\tabcolsep{6pt}
\textbf{Instruction Word} \\
\begin{tabular}{|c|c|c|c|c|c|c|c|c|c|}
	\hline
	10 & qz[2:1] & 1100 & qz[0] & qy[2] & 1 & qy[1:0] & qx[2:0] & ay[3:0] & 0100 \\
	\hline
\end{tabular}

\textbf{Assembler Syntax} \\
EE.VPRELU.S8 qz, qx, qy, ay \\
\textbf{Description} \\
This instruction divides register \lstinline{qx} into 16 data segments by 8 bits. If the value of the segment is not greater than 0, it will be multiplied by the value of the corresponding 8-bit segment in register \lstinline{qy} and right-shifted by the lower 5-bit value in register \lstinline{ay}, and then the calculated result will be assigned to the corresponding 8-bit data segment in register \lstinline{qz}. Otherwise, the value of the segment in \lstinline{qx} will be assigned to \lstinline{qz}.\\
\textbf{Operation}
\begin{lstlisting}[escapeinside=`?]
  qz[  7:  0] = (qx[  7:  0]<=0) ? (qx[  7:  0] * qy[  7:  0]) >> ay[4:0] : qx[  7:  0]
  qz[ 15:  8] = (qx[ 15:  8]<=0) ? (qx[ 15:  8] * qy[ 15:  8]) >> ay[4:0] : qx[ 15:  8]
  ...
  qz[127:120] = (qx[127:120]<=0) ? (qx[127:120] * qy[127:120]) >> ay[4:0] : qx[127:120]
\end{lstlisting}

\clearpage
\subsection{EE.VRELU.S16}\label{instruction:EEVRELUS16}
\setlength\tabcolsep{6pt}
\textbf{Instruction Word} \\
\begin{tabular}{|c|c|c|c|c|c|c|c|}
	\hline
	11 & qs[2:1] & 1101 & qs[0] & 001 & ax[3:0] & ay[3:0] & 0100 \\
	\hline
\end{tabular}

\textbf{Assembler Syntax} \\
EE.VRELU.S16 qs, ax, ay \\
\textbf{Description} \\
This instruction divides register \lstinline{qs} into 8 data segments by 16 bits. If the value of the segment is not greater than 0, it will be multiplied by the value of the lower 16 bits in register \lstinline{ax} and right-shifted by the value of the lower 6 bits in register \lstinline{ay}, and then the result obtained will overwrite the value of the segment. Otherwise, the value of the segment will remain unchanged.\\
\textbf{Operation}
\begin{lstlisting}[escapeinside=`?]
  qs[ 15:  0] = (qs[ 15:  0]<=0) ? (qs[ 15:  0] * ax[15:0]) >> ay[5:0] : qs[ 15:  0]
  qs[ 31: 16] = (qs[ 31: 16]<=0) ? (qs[ 31: 16] * ax[15:0]) >> ay[5:0] : qs[ 31: 16]
  ...
  qs[127:112] = (qs[127:112]<=0) ? (qs[127:112] * ax[15:0]) >> ay[5:0] : qs[127:112]
\end{lstlisting}

\clearpage
\subsection{EE.VRELU.S8}\label{instruction:EEVRELUS8}
\setlength\tabcolsep{6pt}
\textbf{Instruction Word} \\
\begin{tabular}{|c|c|c|c|c|c|c|c|}
	\hline
	11 & qs[2:1] & 1101 & qs[0] & 101 & ax[3:0] & ay[3:0] & 0100 \\
	\hline
\end{tabular}

\textbf{Assembler Syntax} \\
EE.VRELU.S8 qs, ax, ay \\
\textbf{Description} \\
This instruction divides register \lstinline{qs} into 16 data segments by 8 bits. If the value of the segment is not greater than 0, it will be multiplied by the value of the lower 8 bits in register \lstinline{ax} and right-shifted by the value of the lower 5 bits in register \lstinline{ay}, and then the result obtained will overwrite the value of the segment. Otherwise, the value of the segment will remain unchanged.\\
\textbf{Operation}
\begin{lstlisting}[escapeinside=`?]
  qs[  7:  0] = (qs[  7:  0]<=0) ? (qs[  7:  0] * ax[7:0]) >> ay[4:0] : qs[  7:  0]
  qs[ 15:  8] = (qs[ 15:  8]<=0) ? (qs[ 15:  8] * ax[7:0]) >> ay[4:0] : qs[ 15:  8]
  ...
  qs[127:120] = (qs[127:120]<=0) ? (qs[127:120] * ax[7:0]) >> ay[4:0] : qs[127:120]
\end{lstlisting}

\clearpage
\subsection{EE.VSL.32}\label{instruction:EEVSL32}
\setlength\tabcolsep{6pt}
\textbf{Instruction Word} \\
\begin{tabular}{|c|c|c|c|c|c|c|}
	\hline
	11 & qs[2:1] & 1101 & qs[0] & 01111110 & qa[2:0] & 0100 \\
	\hline
\end{tabular}

\textbf{Assembler Syntax} \\
EE.VSL.32 qa, qs \\
\textbf{Description} \\
This instruction performs a left shift on the four 32-bit data segments in register \lstinline{qs} respectively. The left shift amount is the value in the 6-bit special register \lstinline{SAR}. During the shift process, the lower bits are padded with 0. The lower 32 bits of the shift result are stored in the corresponding data segment in register \lstinline{qa}.\\
\textbf{Operation}
\begin{lstlisting}
  qa[ 31:  0] = (qs[ 31:  0] << SAR[5:0])
  qa[ 63: 32] = (qs[ 63: 32] << SAR[5:0])
  qa[ 95: 64] = (qs[ 95: 64] << SAR[5:0])
  qa[127: 96] = (qs[127: 96] << SAR[5:0])
\end{lstlisting}




\clearpage
\subsection{EE.VSMULAS.S16.QACC}\label{instruction:EEVSMULASS16QACC}
\setlength\tabcolsep{6pt}
\textbf{Instruction Word} \\
\begin{tabular}{|c|c|c|c|c|c|c|c|c|}
	\hline
	10 & sel8[2:1] & 1110 & sel8[0] & qy[2] & 1 & qy[1:0] & qx[2:0] & 11000100 \\
	\hline
\end{tabular}

\textbf{Assembler Syntax} \\
EE.VSMULAS.S16.QACC qx, qy, sel8 \\
\textbf{Description} \\
This instruction selects one out of the eight 16-bit data segments in register \lstinline{qy} according to immediate number \lstinline{sel8} and performs a signed multiplication on it and the eight 16-bit data segments in register \lstinline{qx} respectively. The 8 results obtained are added to the corresponding 40-bit data segment in special registers \lstinline{QACC_H} and \lstinline{QACC_L} respectively. Then, the result is saturated to a 40-bit signed number and then stored to the corresponding 40-bit data segment in \lstinline{QACC_H} and \lstinline{QACC_L}.\\
\textbf{Operation}
\begin{lstlisting}
  temp[15:0] = qy[sel8*16+15:sel8*16]
  QACC_L[ 39:  0] = min(max(QACC_L[ 39:  0] + qx[ 15:  0] * temp[15:0], -2^{39}), 2^{39}-1)
  QACC_L[ 79: 40] = min(max(QACC_L[ 79: 40] + qx[ 31: 16] * temp[15:0], -2^{39}), 2^{39}-1)
  QACC_L[119: 80] = min(max(QACC_L[119: 80] + qx[ 47: 32] * temp[15:0], -2^{39}), 2^{39}-1)
  QACC_L[159:120] = min(max(QACC_L[159:120] + qx[ 63: 48] * temp[15:0], -2^{39}), 2^{39}-1)
  QACC_H[ 39:  0] = min(max(QACC_H[ 39:  0] + qx[ 79: 64] * temp[15:0], -2^{39}), 2^{39}-1)
  QACC_H[ 79: 40] = min(max(QACC_H[ 79: 40] + qx[ 95: 80] * temp[15:0], -2^{39}), 2^{39}-1)
  QACC_H[119: 80] = min(max(QACC_H[119: 80] + qx[111: 96] * temp[15:0], -2^{39}), 2^{39}-1)
  QACC_H[159:120] = min(max(QACC_H[159:120] + qx[127:112] * temp[15:0], -2^{39}), 2^{39}-1)
\end{lstlisting}

\clearpage
\subsection{EE.VSMULAS.S16.QACC.LD.INCP}\label{instruction:EEVSMULASS16QACCLDINCP}
\setlength\tabcolsep{6pt}
\textbf{Instruction Word} \\
\begin{tabular}{|c|c|c|c|c|c|c|c|c|c|c|c|}
	\hline
	111000 & qu[2:1] & qy[0] & 111 & qu[0] & sel8[2:0] & qx[1:0] & qy[2:1] & 1100 & as[3:0] & 111 & qx[2] \\
	\hline
\end{tabular}

\textbf{Assembler Syntax} \\
EE.VSMULAS.S16.QACC.LD.INCP qu, as, qx, qy, sel8 \\
\textbf{Description} \\
This instruction selects one out of the eight 16-bit data segments in register \lstinline{qy} according to immediate number \lstinline{sel8} and performs a signed multiplication on it and the eight 16-bit data segments in register \lstinline{qx} respectively. The 8 results obtained are added to the corresponding 40-bit data segment in special registers \lstinline{QACC_H} and \lstinline{QACC_L} respectively. Then, the result is saturated to a 40-bit signed number and then stored to the corresponding 40-bit data segment in \lstinline{QACC_H} and \lstinline{QACC_L}.\\
During the operation, the lower 4 bits of the access address in register \lstinline{as} are forced to be 0, and then the 16-byte data is loaded from the memory to register \lstinline{qu}. After the access, the value in register \lstinline{as} is incremented by 16.\\
\textbf{Operation}
\begin{lstlisting}
  temp[15:0] = qy[sel8*16+15:sel8*16]
  QACC_L[ 39:  0] = min(max(QACC_L[ 39:  0] + qx[ 15:  0] * temp[15:0], -2^{39}), 2^{39}-1)
  QACC_L[ 79: 40] = min(max(QACC_L[ 79: 40] + qx[ 31: 16] * temp[15:0], -2^{39}), 2^{39}-1)
  QACC_L[119: 80] = min(max(QACC_L[119: 80] + qx[ 47: 32] * temp[15:0], -2^{39}), 2^{39}-1)
  QACC_L[159:120] = min(max(QACC_L[159:120] + qx[ 63: 48] * temp[15:0], -2^{39}), 2^{39}-1)
  QACC_H[ 39:  0] = min(max(QACC_H[ 39:  0] + qx[ 79: 64] * temp[15:0], -2^{39}), 2^{39}-1)
  QACC_H[ 79: 40] = min(max(QACC_H[ 79: 40] + qx[ 95: 80] * temp[15:0], -2^{39}), 2^{39}-1)
  QACC_H[119: 80] = min(max(QACC_H[119: 80] + qx[111: 96] * temp[15:0], -2^{39}), 2^{39}-1)
  QACC_H[159:120] = min(max(QACC_H[159:120] + qx[127:112] * temp[15:0], -2^{39}), 2^{39}-1)

  qu[127:0] = load128({as[31:4],4{0}})
  as[31:0] = as[31:0] + 16
\end{lstlisting}

\clearpage
\subsection{EE.VSMULAS.S8.QACC}\label{instruction:EEVSMULASS8QACC}
\setlength\tabcolsep{6pt}
\textbf{Instruction Word} \\
\begin{tabular}{|c|c|c|c|c|c|c|c|c|c|c|}
	\hline
	10 & sel16[3:2] & 1110 & sel16[1] & qy[2] & 0 & qy[1:0] & qx[2:0] & 010 & sel16[0] & 0100 \\
	\hline
\end{tabular}

\textbf{Assembler Syntax} \\
EE.VSMULAS.S8.QACC qx, qy, sel16 \\
\textbf{Description} \\
This instruction selects one out of the 16 8-bit data segments in register \lstinline{qy} according to immediate number \lstinline{sel16} and performs a signed multiplication on it and the 16 8-bit data segments in register \lstinline{qx} respectively. The 16 results obtained are added to the corresponding 20-bit data segment in special registers \lstinline{QACC_H} and \lstinline{QACC_L} respectively. Then, the result is saturated to a 20-bit signed number and then stored to the corresponding 20-bit data segment in \lstinline{QACC_H} and \lstinline{QACC_L}.\\
\textbf{Operation}
\begin{lstlisting}
  temp[7:0] = qy[sel16*8+7:sel16*8]
  QACC_L[ 19:  0] = min(max(QACC_L[ 19:  0] + qx[  7:  0] * temp[7:0], -2^{19}), 2^{19}-1)
  QACC_L[ 39: 20] = min(max(QACC_L[ 39: 20] + qx[ 15:  8] * temp[7:0], -2^{19}), 2^{19}-1)
  QACC_L[ 59: 40] = min(max(QACC_L[ 59: 40] + qx[ 23: 16] * temp[7:0], -2^{19}), 2^{19}-1)
  QACC_L[ 79: 60] = min(max(QACC_L[ 79: 60] + qx[ 31: 24] * temp[7:0], -2^{19}), 2^{19}-1)
  QACC_L[ 99: 80] = min(max(QACC_L[ 99: 80] + qx[ 39: 32] * temp[7:0], -2^{19}), 2^{19}-1)
  QACC_L[119:100] = min(max(QACC_L[119:100] + qx[ 47: 40] * temp[7:0], -2^{19}), 2^{19}-1)
  QACC_L[139:120] = min(max(QACC_L[139:120] + qx[ 55: 48] * temp[7:0], -2^{19}), 2^{19}-1)
  QACC_L[159:140] = min(max(QACC_L[159:140] + qx[ 63: 56] * temp[7:0], -2^{19}), 2^{19}-1)
  QACC_H[ 19:  0] = min(max(QACC_H[ 19:  0] + qx[ 71: 64] * temp[7:0], -2^{19}), 2^{19}-1)
  QACC_H[ 39: 20] = min(max(QACC_H[ 39: 20] + qx[ 79: 72] * temp[7:0], -2^{19}), 2^{19}-1)
  QACC_H[ 59: 40] = min(max(QACC_H[ 59: 40] + qx[ 87: 80] * temp[7:0], -2^{19}), 2^{19}-1)
  QACC_H[ 79: 60] = min(max(QACC_H[ 79: 60] + qx[ 95: 88] * temp[7:0], -2^{19}), 2^{19}-1)
  QACC_H[ 99: 80] = min(max(QACC_H[ 99: 80] + qx[103: 96] * temp[7:0], -2^{19}), 2^{19}-1)
  QACC_H[119:100] = min(max(QACC_H[119:100] + qx[111:104] * temp[7:0], -2^{19}), 2^{19}-1)
  QACC_H[139:120] = min(max(QACC_H[139:120] + qx[119:112] * temp[7:0], -2^{19}), 2^{19}-1)
  QACC_H[159:140] = min(max(QACC_H[159:140] + qx[127:120] * temp[7:0], -2^{19}), 2^{19}-1)
\end{lstlisting}

\clearpage
\subsection{EE.VSMULAS.S8.QACC.LD.INCP}\label{instruction:EEVSMULASS8QACCLDINCP}
\setlength\tabcolsep{6pt}
\textbf{Instruction Word} \\
\begin{tabular}{|c|c|c|c|c|c|c|c|c|c|c|c|c|}
	\hline
	111000 & qu[2:1] & qy[0] & 01 & sel16[0] & qu[0] & sel16[3:1] & qx[1:0] & qy[2:1] & 1100 & as[3:0] & 111 & qx[2] \\
	\hline
\end{tabular}

\textbf{Assembler Syntax} \\
EE.VSMULAS.S8.QACC.LD.INCP qu, as, qx, qy, sel16 \\
\textbf{Description} \\
This instruction selects one out of the 16 8-bit data segments in register \lstinline{qy} according to immediate number \lstinline{sel16} and performs a signed multiplication on it and the 16 8-bit data segments in register \lstinline{qx} respectively. The 16 results obtained are added to the corresponding 20-bit data segment in special registers \lstinline{QACC_H} and \lstinline{QACC_L} respectively. Then, the result is saturated to a 20-bit signed number and then stored to the corresponding 20-bit data segment in \lstinline{QACC_H} and \lstinline{QACC_L}.\\
During the operation, the lower 4 bits of the access address in register \lstinline{as} are forced to be 0, and then the 16-byte data is loaded from the memory to register \lstinline{qu}. After the access, the value in register \lstinline{as} is incremented by 16.\\
\textbf{Operation}
\begin{lstlisting}
  temp[7:0] = qy[sel16*8+7:sel16*8]
  QACC_L[ 19:  0] = min(max(QACC_L[ 19:  0] + qx[  7:  0] * temp[7:0], -2^{19}), 2^{19}-1)
  QACC_L[ 39: 20] = min(max(QACC_L[ 39: 20] + qx[ 15:  8] * temp[7:0], -2^{19}), 2^{19}-1)
  QACC_L[ 59: 40] = min(max(QACC_L[ 59: 40] + qx[ 23: 16] * temp[7:0], -2^{19}), 2^{19}-1)
  QACC_L[ 79: 60] = min(max(QACC_L[ 79: 60] + qx[ 31: 24] * temp[7:0], -2^{19}), 2^{19}-1)
  QACC_L[ 99: 80] = min(max(QACC_L[ 99: 80] + qx[ 39: 32] * temp[7:0], -2^{19}), 2^{19}-1)
  QACC_L[119:100] = min(max(QACC_L[119:100] + qx[ 47: 40] * temp[7:0], -2^{19}), 2^{19}-1)
  QACC_L[139:120] = min(max(QACC_L[139:120] + qx[ 55: 48] * temp[7:0], -2^{19}), 2^{19}-1)
  QACC_L[159:140] = min(max(QACC_L[159:140] + qx[ 63: 56] * temp[7:0], -2^{19}), 2^{19}-1)
  QACC_H[ 19:  0] = min(max(QACC_H[ 19:  0] + qx[ 71: 64] * temp[7:0], -2^{19}), 2^{19}-1)
  QACC_H[ 39: 20] = min(max(QACC_H[ 39: 20] + qx[ 79: 72] * temp[7:0], -2^{19}), 2^{19}-1)
  QACC_H[ 59: 40] = min(max(QACC_H[ 59: 40] + qx[ 87: 80] * temp[7:0], -2^{19}), 2^{19}-1)
  QACC_H[ 79: 60] = min(max(QACC_H[ 79: 60] + qx[ 95: 88] * temp[7:0], -2^{19}), 2^{19}-1)
  QACC_H[ 99: 80] = min(max(QACC_H[ 99: 80] + qx[103: 96] * temp[7:0], -2^{19}), 2^{19}-1)
  QACC_H[119:100] = min(max(QACC_H[119:100] + qx[111:104] * temp[7:0], -2^{19}), 2^{19}-1)
  QACC_H[139:120] = min(max(QACC_H[139:120] + qx[119:112] * temp[7:0], -2^{19}), 2^{19}-1)
  QACC_H[159:140] = min(max(QACC_H[159:140] + qx[127:120] * temp[7:0], -2^{19}), 2^{19}-1)

  qu[127:0] = load128({as[31:4],4{0}})
  as[31:0] = as[31:0] + 16
\end{lstlisting}

\clearpage
\subsection{EE.VSR.32}\label{instruction:EEVSR32}
\setlength\tabcolsep{6pt}
\textbf{Instruction Word} \\
\begin{tabular}{|c|c|c|c|c|c|c|}
	\hline
	11 & qs[2:1] & 1101 & qs[0] & 01111111 & qa[2:0] & 0100 \\
	\hline
\end{tabular}

\textbf{Assembler Syntax} \\
EE.VSR.32 qa, qs \\
\textbf{Description} \\
This instruction performs an arithmetic right shift on the four 32-bit data segments in register \lstinline{qs} respectively. The shift amount is the value in the 6-bit special register \lstinline{SAR}. During the shift process, the higher bits are padded with signed bit. The lower 32 bits of the shift result are stored in the corresponding data segment in register \lstinline{qa}.\\
\textbf{Operation}
\begin{lstlisting}
  qa[ 31:  0] = (qs[ 31:  0] >> SAR[5:0])
  qa[ 63: 32] = (qs[ 63: 32] >> SAR[5:0])
  qa[ 95: 64] = (qs[ 95: 64] >> SAR[5:0])
  qa[127: 96] = (qs[127: 96] >> SAR[5:0])
\end{lstlisting}

\clearpage
\subsection{EE.VST.128.IP}\label{instruction:EEVST128IP}
\setlength\tabcolsep{6pt}
\textbf{Instruction Word} \\
\begin{tabular}{|c|c|c|c|c|c|c|c|}
	\hline
	1 & imm16[7] & qv[2:1] & 1010 & qv[0] & imm16[6:0] & as[3:0] & 0100 \\
	\hline
\end{tabular}

\textbf{Assembler Syntax} \\
EE.VST.128.IP qv, as, -2048..2032 \\
\textbf{Description} \\
This instruction forces the lower 4 bits of the access address in register \lstinline{as} to 0 and stores the 128 bits in register \lstinline{qv} to memory. After the access is completed, the value in register \lstinline{as} is incremented by 8-bit sign-extended constant in the instruction code segment left-shifted by 4.\\
\textbf{Operation}
\begin{lstlisting}
  qv[127:0] => store128({as[31:4],4{0}})
  as[31:0] = as[31:0] + {20{imm16[7]},imm16[7:0],4{0}}
\end{lstlisting}

\clearpage
\subsection{EE.VST.128.XP}\label{instruction:EEVST128XP}
\setlength\tabcolsep{6pt}
\textbf{Instruction Word} \\
\begin{tabular}{|c|c|c|c|c|c|c|c|}
	\hline
	10 & qv[2:1] & 1101 & qv[0] & 111 & ad[3:0] & as[3:0] & 0100 \\
	\hline
\end{tabular}

\textbf{Assembler Syntax} \\
EE.VST.128.XP qv, as, ad \\
\textbf{Description} \\
This instruction forces the lower 4 bits of the access address in register \lstinline{as} to 0 and stores the 128 bits in register \lstinline{qv} to memory. After the access is completed, the value in register \lstinline{as} is incremented by the value in register \lstinline{ad}.\\
\textbf{Operation}
\begin{lstlisting}
  qv[127:0] => store128({as[31:4],4{0}})
  as[31:0] = as[31:0] + ad[31:0]
\end{lstlisting}

\clearpage
\subsection{EE.VST.H.64.IP}\label{instruction:EEVSTH64IP}
\setlength\tabcolsep{6pt}
\textbf{Instruction Word} \\
\begin{tabular}{|c|c|c|c|c|c|c|c|}
	\hline
	1 & imm8[7] & qv[2:1] & 1011 & qv[0] & imm8[6:0] & as[3:0] & 0100 \\
	\hline
\end{tabular}

\textbf{Assembler Syntax} \\
EE.VST.H.64.IP qv, as, -1024..1016 \\
\textbf{Description} \\
This instruction forces the lower 3 bits of the access address in register \lstinline{as} to 0 and stores the upper 64 bits in register \lstinline{qv} to memory. After the access is completed, the value in register \lstinline{as} is incremented by 8-bit sign-extended constant in the instruction code segment left-shifted by 3.\\
\textbf{Operation}
\begin{lstlisting}
  qv[127: 64] => store64({as[31:3],3{0}})
  as[31:0] = as[31:0] + {21{imm8[7]},imm8[7:0],3{0}}
\end{lstlisting}

\clearpage
\subsection{EE.VST.H.64.XP}\label{instruction:EEVSTH64XP}
\setlength\tabcolsep{6pt}
\textbf{Instruction Word} \\
\begin{tabular}{|c|c|c|c|c|c|c|c|}
	\hline
	11 & qv[2:1] & 1101 & qv[0] & 000 & ad[3:0] & as[3:0] & 0100 \\
	\hline
\end{tabular}

\textbf{Assembler Syntax} \\
EE.VST.H.64.XP qv, as, ad \\
\textbf{Description} \\
This instruction forces the lower 3 bits of the access address in register \lstinline{as} to 0 and stores the upper 64 bits in register \lstinline{qv} to memory. After the access is completed, the value in register \lstinline{as} is incremented by the value in register \lstinline{ad}.\\
\textbf{Operation}
\begin{lstlisting}
  qv[127: 64] => store64({as[31:3],3{0}})
  as  = as + ad[31:0]
\end{lstlisting}

\clearpage
\subsection{EE.VST.L.64.IP}\label{instruction:EEVSTL64IP}
\setlength\tabcolsep{6pt}
\textbf{Instruction Word} \\
\begin{tabular}{|c|c|c|c|c|c|c|c|}
	\hline
	1 & imm8[7] & qv[2:1] & 0100 & qv[0] & imm8[6:0] & as[3:0] & 0100 \\
	\hline
\end{tabular}

\textbf{Assembler Syntax} \\
EE.VST.L.64.IP qv, as, -1024..1016 \\
\textbf{Description} \\
This instruction forces the lower 3 bits of the access address in register \lstinline{as} to 0 and stores the lower 64 bits in register \lstinline{qv} to memory. After the access is completed, the value in register \lstinline{as} is incremented by 8-bit sign-extended constant in the instruction code segment left-shifted by 3.\\
\textbf{Operation}
\begin{lstlisting}
  qv[ 63:  0] => store64({as[31:3],3{0}})
  as[31:0] = as[31:0] + {21{imm8[7]},imm8[7:0],3{0}}
\end{lstlisting}

\clearpage
\subsection{EE.VST.L.64.XP}\label{instruction:EEVSTL64XP}
\setlength\tabcolsep{6pt}
\textbf{Instruction Word} \\
\begin{tabular}{|c|c|c|c|c|c|c|c|}
	\hline
	11 & qv[2:1] & 1101 & qv[0] & 100 & ad[3:0] & as[3:0] & 0100 \\
	\hline
\end{tabular}

\textbf{Assembler Syntax} \\
EE.VST.L.64.XP qv, as, ad \\
\textbf{Description} \\
This instruction forces the lower 3 bits of the access address in register \lstinline{as} to 0 and stores the lower 64 bits in register \lstinline{qv} to memory. After the access is completed, the value in register \lstinline{as} is incremented by the value in register \lstinline{ad}.\\
\textbf{Operation}
\begin{lstlisting}
  qv[63:0] => store64({as[31:3],3{0}})
  as  = as + ad[31:0]
\end{lstlisting}

\clearpage
\subsection{EE.VSUBS.S16}\label{instruction:EEVSUBSS16}
\setlength\tabcolsep{6pt}
\textbf{Instruction Word} \\
\begin{tabular}{|c|c|c|c|c|c|c|c|c|}
	\hline
	10 & qa[2:1] & 1110 & qa[0] & qy[2] & 1 & qy[1:0] & qx[2:0] & 11010100 \\
	\hline
\end{tabular}

\textbf{Assembler Syntax} \\
EE.VSUBS.S16 qa, qx, qy \\
\textbf{Description} \\
This instruction performs a vector subtraction on 16-bit data. Registers \lstinline{qx} and \lstinline{qy} are the subtrahend and the minuend respectively. Then, the 8 results obtained from the calculation are saturated and then written into register \lstinline{qa}.\\
\textbf{Operation}
\begin{lstlisting}
  qa[ 15:  0] = min(max(qx[ 15:  0] - qy[ 15:  0], -2^{15}), 2^{15}-1)
  qa[ 31: 16] = min(max(qx[ 31: 16] - qy[ 31: 16], -2^{15}), 2^{15}-1)
  ...
\end{lstlisting}

\clearpage
\subsection{EE.VSUBS.S16.LD.INCP}\label{instruction:EEVSUBSS16LDINCP}
\setlength\tabcolsep{6pt}
\textbf{Instruction Word} \\
\begin{tabular}{|c|c|c|c|c|c|c|c|c|c|c|c|}
	\hline
	111000 & qu[2:1] & qy[0] & 100 & qu[0] & qa[2:0] & qx[1:0] & qy[2:1] & 1101 & as[3:0] & 111 & qx[2] \\
	\hline
\end{tabular}

\textbf{Assembler Syntax} \\
EE.VSUBS.S16.LD.INCP qu, as, qa, qx, qy \\
\textbf{Description} \\
This instruction performs a vector subtraction on 16-bit data. Registers \lstinline{qx} and \lstinline{qy} are the subtrahend and the minuend respectively. Then, the 8 results obtained from the calculation are saturated and then written into register \lstinline{qa}. \\ During the operation, the lower 4 bits of the access address in register \lstinline{as} are forced to be 0, and then the 16-byte data is loaded from the memory to register \lstinline{qu}. After the access, the value in register \lstinline{as} is incremented by 16.\\
\textbf{Operation}
\begin{lstlisting}
  qa[ 15:  0] = min(max(qx[ 15:  0] - qy[ 15:  0], -2^{15}), 2^{15}-1)
  qa[ 31: 16] = min(max(qx[ 31: 16] - qy[ 31: 16], -2^{15}), 2^{15}-1)
  ...

  qu[127:0] = load128({as[31:4],4{0}})
  as[31:0] = as[31:0] + 16
\end{lstlisting}

\clearpage
\subsection{EE.VSUBS.S16.ST.INCP}\label{instruction:EEVSUBSS16STINCP}
\setlength\tabcolsep{6pt}
\textbf{Instruction Word} \\
\begin{tabular}{|c|c|c|c|c|c|c|c|c|c|c|}
	\hline
	11101000 & qy[0] & qv[2:0] & 1 & qa[2:0] & qx[1:0] & qy[2:1] & 0001 & as[3:0] & 111 & qx[2] \\
	\hline
\end{tabular}

\textbf{Assembler Syntax} \\
EE.VSUBS.S16.ST.INCP qv, as, qa, qx, qy \\
\textbf{Description} \\
This instruction performs a vector subtraction on 16-bit data. Registers \lstinline{qx} and \lstinline{qy} are the subtrahend and the minuend respectively. Then, the 8 results obtained from the calculation are saturated and then written into register \lstinline{qa}. \\ During the operation, the instruction forces the lower 4 bits of the access address in register \lstinline{as} to 0 and stores the value in register \lstinline{qv} to memory. After the access, the value in register \lstinline{as} is incremented by 16.\\
\textbf{Operation}
\begin{lstlisting}
  qa[ 15:  0] = min(max(qx[ 15:  0] - qy[ 15:  0], -2^{15}), 2^{15}-1)
  qa[ 31: 16] = min(max(qx[ 31: 16] - qy[ 31: 16], -2^{15}), 2^{15}-1)
  ...

  qv[127:0] => store128({as[31:4],4{0}})
  as[31:0] = as[31:0] + 16
\end{lstlisting}

\clearpage
\subsection{EE.VSUBS.S32}\label{instruction:EEVSUBSS32}
\setlength\tabcolsep{6pt}
\textbf{Instruction Word} \\
\begin{tabular}{|c|c|c|c|c|c|c|c|c|}
	\hline
	10 & qa[2:1] & 1110 & qa[0] & qy[2] & 1 & qy[1:0] & qx[2:0] & 11100100 \\
	\hline
\end{tabular}

\textbf{Assembler Syntax} \\
EE.VSUBS.S32 qa, qx, qy \\
\textbf{Description} \\
This instruction performs a vector subtraction on 32-bit data. Registers \lstinline{qx} and \lstinline{qy} are the subtrahend and the minuend respectively. Then, the 4 results obtained from the calculation are saturated and then written into register \lstinline{qa}.\\
\textbf{Operation}
\begin{lstlisting}
  qa[ 31:  0] = min(max(qx[ 31:  0] - qy[ 31:  0], -2^{31}), 2^{31}-1)
  qa[ 63: 32] = min(max(qx[ 63: 32] - qy[ 63: 32], -2^{31}), 2^{31}-1)
  ...
\end{lstlisting}

\clearpage
\subsection{EE.VSUBS.S32.LD.INCP}\label{instruction:EEVSUBSS32LDINCP}
\setlength\tabcolsep{6pt}
\textbf{Instruction Word} \\
\begin{tabular}{|c|c|c|c|c|c|c|c|c|c|c|c|}
	\hline
	111000 & qu[2:1] & qy[0] & 101 & qu[0] & qa[2:0] & qx[1:0] & qy[2:1] & 1101 & as[3:0] & 111 & qx[2] \\
	\hline
\end{tabular}

\textbf{Assembler Syntax} \\
EE.VSUBS.S32.LD.INCP qu, as, qa, qx, qy \\
\textbf{Description} \\
This instruction performs a vector subtraction on 32-bit data. Registers \lstinline{qx} and \lstinline{qy} are the subtrahend and the minuend respectively. Then, the 4 results obtained from the calculation are saturated and then written into register \lstinline{qa}. \\ During the operation, the lower 4 bits of the access address in register \lstinline{as} are forced to be 0, and then the 16-byte data is loaded from the memory to register \lstinline{qu}. After the access, the value in register \lstinline{as} is incremented by 16.\\
\textbf{Operation}
\begin{lstlisting}
  qa[31: 0] = min(max(qx[31: 0] - qy[31: 0], -2^{31}), 2^{31}-1)
  qa[63:32] = min(max(qx[63:32] - qy[63:32], -2^{31}), 2^{31}-1)
  ...

  qu[127:0] = load128({as[31:4],4{0}})
  as[31:0] = as[31:0] + 16
\end{lstlisting}

\clearpage
\subsection{EE.VSUBS.S32.ST.INCP}\label{instruction:EEVSUBSS32STINCP}
\setlength\tabcolsep{6pt}
\textbf{Instruction Word} \\
\begin{tabular}{|c|c|c|c|c|c|c|c|c|c|c|}
	\hline
	11101000 & qy[0] & qv[2:0] & 1 & qa[2:0] & qx[1:0] & qy[2:1] & 0010 & as[3:0] & 111 & qx[2] \\
	\hline
\end{tabular}

\textbf{Assembler Syntax} \\
EE.VSUBS.S32.ST.INCP qv, as, qa, qx, qy \\
\textbf{Description} \\
This instruction performs a vector subtraction on 32-bit data. Registers \lstinline{qx} and \lstinline{qy} are the subtrahend and the minuend respectively. Then, the 4 results obtained from the calculation are saturated and then written into register \lstinline{qa}. \\ During the operation, the instruction forces the lower 4 bits of the access address in register \lstinline{as} to 0 and stores the value in register \lstinline{qv} to memory. After the access, the value in register \lstinline{as} is incremented by 16.\\
\textbf{Operation}
\begin{lstlisting}
  qa[ 31:  0] = min(max(qx[ 31:  0] - qy[ 31:  0], -2^{31}), 2^{31}-1)
  qa[ 63: 32] = min(max(qx[ 63: 32] - qy[ 63: 32], -2^{31}), 2^{31}-1)
  ...

  qv[127:0] => store128({as[31:4],4{0}})
  as[31:0] = as[31:0] + 16
\end{lstlisting}

\clearpage
\subsection{EE.VSUBS.S8}\label{instruction:EEVSUBSS8}
\setlength\tabcolsep{6pt}
\textbf{Instruction Word} \\
\begin{tabular}{|c|c|c|c|c|c|c|c|c|}
	\hline
	10 & qa[2:1] & 1110 & qa[0] & qy[2] & 1 & qy[1:0] & qx[2:0] & 11110100 \\
	\hline
\end{tabular}

\textbf{Assembler Syntax} \\
EE.VSUBS.S8 qa, qx, qy \\
\textbf{Description} \\
This instruction performs a vector subtraction on 8-bit data. Registers \lstinline{qx} and \lstinline{qy} are the subtrahend and the minuend respectively. Then, the 16 results obtained from the calculation are saturated and then written into register \lstinline{qa}.\\
\textbf{Operation}
\begin{lstlisting}
  qa[  7:  0] = min(max(qx[  7:  0] - qy[  7:  0], -2^{7}), 2^{7}-1)
  qa[ 15:  8] = min(max(qx[ 15:  8] - qy[ 15:  8], -2^{7}), 2^{7}-1)
  ...
\end{lstlisting}

\clearpage
\subsection{EE.VSUBS.S8.LD.INCP}\label{instruction:EEVSUBSS8LDINCP}
\setlength\tabcolsep{6pt}
\textbf{Instruction Word} \\
\begin{tabular}{|c|c|c|c|c|c|c|c|c|c|c|c|}
	\hline
	111000 & qu[2:1] & qy[0] & 110 & qu[0] & qa[2:0] & qx[1:0] & qy[2:1] & 1101 & as[3:0] & 111 & qx[2] \\
	\hline
\end{tabular}

\textbf{Assembler Syntax} \\
EE.VSUBS.S8.LD.INCP qu, as, qa, qx, qy \\
\textbf{Description} \\
This instruction performs a vector subtraction on 8-bit data. Registers \lstinline{qx} and \lstinline{qy} are the subtrahend and the minuend respectively. Then, the 16 results obtained from the calculation are saturated and then written into register \lstinline{qa}. \\ During the operation, the lower 4 bits of the access address in register \lstinline{as} are forced to be 0, and then the 16-byte data is loaded from the memory to register \lstinline{qu}. After the access, the value in register \lstinline{as} is incremented by 16.\\
\textbf{Operation}
\begin{lstlisting}
  qa[  7:  0] = min(max(qx[  7:  0] - qy[  7:  0], -2^{7}), 2^{7}-1)
  qa[ 15:  8] = min(max(qx[ 15:  8] - qy[ 15:  8], -2^{7}), 2^{7}-1)
  ...

  qu[127:0] = load128({as[31:4],4{0}})
  as[31:0] = as[31:0] + 16
\end{lstlisting}

\clearpage
\subsection{EE.VSUBS.S8.ST.INCP}\label{instruction:EEVSUBSS8STINCP}
\setlength\tabcolsep{6pt}
\textbf{Instruction Word} \\
\begin{tabular}{|c|c|c|c|c|c|c|c|c|c|c|}
	\hline
	11101000 & qy[0] & qv[2:0] & 1 & qa[2:0] & qx[1:0] & qy[2:1] & 0011 & as[3:0] & 111 & qx[2] \\
	\hline
\end{tabular}

\textbf{Assembler Syntax} \\
EE.VSUBS.S8.ST.INCP qv, as, qa, qx, qy \\
\textbf{Description} \\
This instruction performs a vector subtraction on 8-bit data. Registers \lstinline{qx} and \lstinline{qy} are the subtrahend and the minuend respectively. Then, the 16 results obtained from the calculation are saturated and then written into register \lstinline{qa}. \\ During the operation, the instruction forces the lower 4 bits of the access address in register \lstinline{as} to 0 and stores the value in register \lstinline{qv} to memory. After the access, the value in register \lstinline{as} is incremented by 16.\\
\textbf{Operation}
\begin{lstlisting}
  qa[  7:  0] = min(max(qx[  7:  0] - qy[  7:  0], -2^{7}), 2^{7}-1)
  qa[ 15:  8] = min(max(qx[ 15:  8] - qy[ 15:  8], -2^{7}), 2^{7}-1)
  ...

  qv[127:0] => store128({as[31:4],4{0}})
  as[31:0] = as[31:0] + 16
\end{lstlisting}

\clearpage
\subsection{EE.VUNZIP.16}\label{instruction:EEVUNZIP16}
\setlength\tabcolsep{6pt}
\textbf{Instruction Word} \\
\begin{tabular}{|c|c|c|c|c|c|}
	\hline
	11 & qs1[2:1] & 1100 & qs1[0] & qs0[2:0] & 001110000100 \\
	\hline
\end{tabular}

\textbf{Assembler Syntax} \\
EE.VUNZIP.16 qs0, qs1 \\
\textbf{Description} \\
This instruction implements the unzip algorithm on 16-bit vector data.\\
\textbf{Operation}
\begin{lstlisting}
  qs0[ 15:  0] = qs0[ 15:  0]
  qs0[ 31: 16] = qs0[ 47: 32]
  qs0[ 47: 32] = qs0[ 79: 64]
  qs0[ 63: 48] = qs0[111: 96]
  qs0[ 79: 64] = qs1[ 15:  0]
  qs0[ 95: 80] = qs1[ 47: 32]
  qs0[111: 96] = qs1[ 79: 64]
  qs0[127:112] = qs1[111: 96]
  qs1[ 15:  0] = qs0[ 31: 16]
  qs1[ 31: 16] = qs0[ 63: 48]
  qs1[ 47: 32] = qs0[ 95: 80]
  qs1[ 63: 48] = qs0[127:112]
  qs1[ 79: 64] = qs1[ 31: 16]
  qs1[ 95: 80] = qs1[ 63: 48]
  qs1[111: 96] = qs1[ 95: 80]
  qs1[127:112] = qs1[127:112]
\end{lstlisting}

\clearpage
\subsection{EE.VUNZIP.32}\label{instruction:EEVUNZIP32}
\setlength\tabcolsep{6pt}
\textbf{Instruction Word} \\
\begin{tabular}{|c|c|c|c|c|c|}
	\hline
	11 & qs1[2:1] & 1100 & qs1[0] & qs0[2:0] & 001110010100 \\
	\hline
\end{tabular}

\textbf{Assembler Syntax} \\
EE.VUNZIP.32 qs0, qs1 \\
\textbf{Description} \\
This instruction implements the unzip algorithm on 32-bit vector data.\\
\textbf{Operation}
\begin{lstlisting}
  qs0[ 31:  0] = qs0[ 31:  0]
  qs0[ 63: 32] = qs0[ 95: 64]
  qs0[ 95: 64] = qs1[ 31:  0]
  qs0[127: 96] = qs1[ 95: 64]
  qs1[ 31:  0] = qs0[ 63: 32]
  qs1[ 63: 32] = qs0[127: 96]
  qs1[ 95: 64] = qs1[ 63: 32]
  qs1[127: 96] = qs1[127: 96]
\end{lstlisting}

\clearpage
\subsection{EE.VUNZIP.8}\label{instruction:EEVUNZIP8}
\setlength\tabcolsep{6pt}
\textbf{Instruction Word} \\
\begin{tabular}{|c|c|c|c|c|c|}
	\hline
	11 & qs1[2:1] & 1100 & qs1[0] & qs0[2:0] & 001110100100 \\
	\hline
\end{tabular}

\textbf{Assembler Syntax} \\
EE.VUNZIP.8 qs0, qs1 \\
\textbf{Description} \\
This instruction implements the unzip algorithm on 8-bit vector data.\\
\textbf{Operation}
\begin{lstlisting}
  qs0[  7:  0] = qs0[  7:  0]
  qs0[ 15:  8] = qs0[ 23: 16]
  qs0[ 23: 16] = qs0[ 39: 32]
  qs0[ 31: 24] = qs0[ 55: 48]
  qs0[ 39: 32] = qs0[ 71: 64]
  qs0[ 47: 40] = qs0[ 87: 80]
  qs0[ 55: 48] = qs0[103: 96]
  qs0[ 63: 56] = qs0[119:112]
  qs0[ 71: 64] = qs1[  7:  0]
  qs0[ 79: 72] = qs1[ 23: 16]
  qs0[ 87: 80] = qs1[ 39: 32]
  qs0[ 95: 88] = qs1[ 55: 48]
  qs0[103: 96] = qs1[ 71: 64]
  qs0[111:104] = qs1[ 87: 80]
  qs0[119:112] = qs1[103: 96]
  qs0[127:120] = qs1[119:112]
  qs1[  7:  0] = qs0[ 15:  8]
  qs1[ 15:  8] = qs0[ 31: 24]
  qs1[ 23: 16] = qs0[ 47: 40]
  qs1[ 31: 24] = qs0[ 63: 56]
  qs1[ 39: 32] = qs0[ 79: 72]
  qs1[ 47: 40] = qs0[ 95: 88]
  qs1[ 55: 48] = qs0[111:104]
  qs1[ 63: 56] = qs0[127:120]
  qs1[ 71: 64] = qs1[ 15:  8]
  qs1[ 79: 72] = qs1[ 31: 24]
  qs1[ 87: 80] = qs1[ 47: 40]
  qs1[ 95: 88] = qs1[ 63: 56]
  qs1[103: 96] = qs1[ 79: 72]
  qs1[111:104] = qs1[ 95: 88]
  qs1[119:112] = qs1[111:104]
  qs1[127:120] = qs1[127:120]
\end{lstlisting}

\clearpage
\subsection{EE.VZIP.16}\label{instruction:EEVZIP16}
\setlength\tabcolsep{6pt}
\textbf{Instruction Word} \\
\begin{tabular}{|c|c|c|c|c|c|}
	\hline
	11 & qs1[2:1] & 1100 & qs1[0] & qs0[2:0] & 001110110100 \\
	\hline
\end{tabular}

\textbf{Assembler Syntax} \\
EE.VZIP.16 qs0, qs1 \\
\textbf{Description} \\
This instruction implements the zip algorithm on 16-bit vector data.\\
\textbf{Operation}
\begin{lstlisting}
  qs0[ 15:  0] = qs0[ 15:  0]
  qs0[ 31: 16] = qs1[ 15:  0]
  qs0[ 47: 32] = qs0[ 31: 16]
  qs0[ 63: 48] = qs1[ 31: 16]
  qs0[ 79: 64] = qs0[ 47: 32]
  qs0[ 95: 80] = qs1[ 47: 32]
  qs0[111: 96] = qs0[ 63: 48]
  qs0[127:112] = qs1[ 63: 48]
  qs1[ 15:  0] = qs0[ 79: 64]
  qs1[ 31: 16] = qs1[ 79: 64]
  qs1[ 47: 32] = qs0[ 95: 80]
  qs1[ 63: 48] = qs1[ 95: 80]
  qs1[ 79: 64] = qs0[111: 96]
  qs1[ 95: 80] = qs1[111: 96]
  qs1[111: 96] = qs0[127:112]
  qs1[127:112] = qs1[127:112]
\end{lstlisting}

\clearpage
\subsection{EE.VZIP.32}\label{instruction:EEVZIP32}
\setlength\tabcolsep{6pt}
\textbf{Instruction Word} \\
\begin{tabular}{|c|c|c|c|c|c|}
	\hline
	11 & qs1[2:1] & 1100 & qs1[0] & qs0[2:0] & 001111000100 \\
	\hline
\end{tabular}

\textbf{Assembler Syntax} \\
EE.VZIP.32 qs0, qs1 \\
\textbf{Description} \\
This instruction implements the zip algorithm on 32-bit vector data.\\
\textbf{Operation}
\begin{lstlisting}
  qs0[ 31:  0] = qs0[ 31:  0]
  qs0[ 63: 32] = qs1[ 31:  0]
  qs0[ 95: 64] = qs0[ 63: 32]
  qs0[127: 96] = qs1[ 63: 32]
  qs1[ 31:  0] = qs0[ 95: 64]
  qs1[ 63: 32] = qs1[ 95: 64]
  qs1[ 95: 64] = qs0[127: 96]
  qs1[127: 96] = qs1[127: 96]
\end{lstlisting}

\clearpage
\subsection{EE.VZIP.8}\label{instruction:EEVZIP8}
\setlength\tabcolsep{6pt}
\textbf{Instruction Word} \\
\begin{tabular}{|c|c|c|c|c|c|}
	\hline
	11 & qs1[2:1] & 1100 & qs1[0] & qs0[2:0] & 001111010100 \\
	\hline
\end{tabular}

\textbf{Assembler Syntax} \\
EE.VZIP.8 qs0, qs1 \\
\textbf{Description} \\
This instruction implements the zip algorithm on 8-bit vector data.\\
\textbf{Operation}
\begin{lstlisting}
  qs0[  7:  0] = qs0[  7:  0]
  qs0[ 15:  8] = qs1[  7:  0]
  qs0[ 23: 16] = qs0[ 15:  8]
  qs0[ 31: 24] = qs1[ 15:  8]
  qs0[ 39: 32] = qs0[ 23: 16]
  qs0[ 47: 40] = qs1[ 23: 16]
  qs0[ 55: 48] = qs0[ 31: 24]
  qs0[ 63: 56] = qs1[ 31: 24]
  qs0[ 71: 64] = qs0[ 39: 32]
  qs0[ 79: 72] = qs1[ 39: 32]
  qs0[ 87: 80] = qs0[ 47: 40]
  qs0[ 95: 88] = qs1[ 47: 40]
  qs0[103: 96] = qs0[ 55: 48]
  qs0[111:104] = qs1[ 55: 48]
  qs0[119:112] = qs0[ 63: 56]
  qs0[127:120] = qs1[ 63: 56]
  qs1[  7:  0] = qs0[ 71: 64]
  qs1[ 15:  8] = qs1[ 71: 64]
  qs1[ 23: 16] = qs0[ 79: 72]
  qs1[ 31: 24] = qs1[ 79: 72]
  qs1[ 39: 32] = qs0[ 87: 80]
  qs1[ 47: 40] = qs1[ 87: 80]
  qs1[ 55: 48] = qs0[ 95: 88]
  qs1[ 63: 56] = qs1[ 95: 88]
  qs1[ 71: 64] = qs0[103: 96]
  qs1[ 79: 72] = qs1[103: 96]
  qs1[ 87: 80] = qs0[111:104]
  qs1[ 95: 88] = qs1[111:104]
  qs1[103: 96] = qs0[119:112]
  qs1[111:104] = qs1[119:112]
  qs1[119:112] = qs0[127:120]
  qs1[127:120] = qs1[127:120]
\end{lstlisting}

\clearpage
\subsection{EE.WR\_MASK\_GPIO\_OUT}\label{instruction:EEWRMASKGPIOOUT}
\setlength\tabcolsep{6pt}
\textbf{Instruction Word} \\
\begin{tabular}{|c|c|c|c|}
	\hline
	011100100100 & ax[3:0] & as[3:0] & 0100 \\
	\hline
\end{tabular}

\textbf{Assembler Syntax} \\
EE.WR\_MASK\_GPIO\_OUT as, ax \\
\textbf{Description} \\
It is a dedicated CPU GPIO instruction to set specified bits in \lstinline{GPIO_OUT}. The lower 8 bits in register \lstinline{ax} store the mask, and the lower 8 bits in register \lstinline{as} store the assignment content.\\
\textbf{Operation}
\begin{lstlisting}
  GPIO_OUT[7:0] =  (GPIO_OUT[7:0] & ~ax[7:0]) | (as[7:0] & ax[7:0])
\end{lstlisting}

\clearpage
\subsection{EE.XORQ}\label{instruction:EEXORQ}
\setlength\tabcolsep{6pt}
\textbf{Instruction Word} \\
\begin{tabular}{|c|c|c|c|c|c|c|c|c|c|c|}
	\hline
	11 & qa[2:1] & 1101 & qa[0] & 011 & qy[2:1] & 01 & qx[2:1] & qy[0] & qx[0] & 0100 \\
	\hline
\end{tabular}

\textbf{Assembler Syntax} \\
EE.XORQ qa, qx, qy \\
\textbf{Description} \\
This instruction performs a bitwise XOR operation on registers \lstinline{qx} and \lstinline{qy} and writes the result of the logical operation to register \lstinline{qa}.\\
\textbf{Operation}
\begin{lstlisting}
  qa = qx ^ qy
\end{lstlisting}

\clearpage
\subsection{EE.ZERO.ACCX}\label{instruction:EEZEROACCX}
\setlength\tabcolsep{6pt}
\textbf{Instruction Word} \\
\begin{tabular}{|c|}
	\hline
	001001010000100000000100 \\
	\hline
\end{tabular}

\textbf{Assembler Syntax} \\
EE.ZERO.ACCX  \\
\textbf{Description} \\
This instruction clears the value in special register \lstinline{ACCX} to 0.\\
\textbf{Operation}
\begin{lstlisting}
  ACCX = 0
\end{lstlisting}

\clearpage
\subsection{EE.ZERO.Q}\label{instruction:EEZEROQ}
\setlength\tabcolsep{6pt}
\textbf{Instruction Word} \\
\begin{tabular}{|c|c|c|c|c|}
	\hline
	11 & qa[2:1] & 1101 & qa[0] & 111111110100100 \\
	\hline
\end{tabular}

\textbf{Assembler Syntax} \\
EE.ZERO.Q qa \\
\textbf{Description} \\
This instruction clears the value in register \lstinline{qa} to 0.\\
\textbf{Operation}
\begin{lstlisting}
  qa = 0
\end{lstlisting}

\clearpage
\subsection{EE.ZERO.QACC}\label{instruction:EEZEROQACC}
\setlength\tabcolsep{6pt}
\textbf{Instruction Word} \\
\begin{tabular}{|c|}
	\hline
	001001010000100001000100 \\
	\hline
\end{tabular}

\textbf{Assembler Syntax} \\
EE.ZERO.QACC  \\
\textbf{Description} \\
This instruction clears the values in special registers \lstinline{QACC_L} and \lstinline{QACC_H} to 0.\\
\textbf{Operation}
\begin{lstlisting}
  QACC_L = 0
  QACC_H = 0
\end{lstlisting}


\clearpage
\subsection{LD.QR}\label{instruction:LDQR}
\setlength\tabcolsep{6pt}
\textbf{Instruction Word} \\
\begin{tabular}{|c|c|c|c|c|c|c|c|}
	\hline
	11 & qu[2:1] & 1101 & qu[0] & 010 & imm[3:0] & as[3:0] & 0100 \\
	\hline
\end{tabular}

\textbf{Assembler Syntax} \\
LD.QR qu, as, imm, -128..112  \\
\textbf{Description} \\
This instruction loads 128 bits from memory to the target QR register \lstinline{qu}.\\
\textbf{Operation}
\begin{lstlisting}
  qu = load128(as + imm)
\end{lstlisting}

\clearpage
\subsection{ST.QR}\label{instruction:STQR}
\setlength\tabcolsep{6pt}
\textbf{Instruction Word} \\
\begin{tabular}{|c|c|c|c|c|c|c|c|}
	\hline
	11 & qs[2:1] & 1101 & qs[0] & 110 & imm[3:0] & as[3:0] & 0100 \\
	\hline
\end{tabular}

\textbf{Assembler Syntax} \\
LD.QR qs, as, imm, -128..112  \\
\textbf{Description} \\
This instruction stores 128 bits from the source QR register \lstinline{qs} to memory.\\
\textbf{Operation}
\begin{lstlisting}
  qs => store128(as + imm)
\end{lstlisting}


\clearpage
\subsection{MV.QR}\label{instruction:MVQR}
\setlength\tabcolsep{6pt}
\textbf{Instruction Word} \\
\begin{tabular}{|c|c|c|c|c|c|c|c|c|}
	\hline
	10 & qu[2:1] & 1111 & qu[0] & 000 & qs[2:1] & 000 & qs[0] & 00100 \\
	\hline
\end{tabular}

\textbf{Assembler Syntax} \\
MV.QR qu, qs \\
\textbf{Description} \\
This instruction moves the value from the source QR register \lstinline{qs} to the target QR register \lstinline{qu}.\\
\textbf{Operation}
\begin{lstlisting}
  qu = qs
\end{lstlisting}
